% Équations et inéquations polynomiales — Mathématiques 1ère D — Cours signé OnBuch+
% Programme officiel MINESEC. Compiler avec : tectonic N01-equations-inequations-polynomiales.tex
\newcommand{\DOCMATIERE}{Mathématiques}
\newcommand{\DOCNIVEAU}{1ère D}
\newcommand{\DOCMODULE}{Relations et opérations fondamentales dans l'ensemble des nombres réels}
\newcommand{\DOCLECON}{N01}
\newcommand{\DOCTITRE}{Équations et inéquations polynomiales}
% OnBuch+ — préambule « cours signés » ENRICHI (compilé avec tectonic / XeLaTeX)
% Système visuel « Pop » d'OnBuch+ : fond crème (jamais blanc), bordures encre
% épaisses, ombres « dures » décalées, titres Archivo Black en capitales.
% Version MATHÉMATIQUES : graphes (pgfplots/tikz), boîtes pédagogiques
% (définition, propriété, méthode, attention, le savais-tu, exercice résolu,
% exercice, TP, bilan), page de garde et signature OnBuch+ ancrée en bas.
%
% Macros attendues dans main.tex AVANT \input{preamble} :
%   \DOCMATIERE  \DOCNIVEAU  \DOCMODULE  \DOCLECON (numéro)  \DOCTITRE
\documentclass[11pt]{article}

\usepackage{fontspec}
\usepackage[french]{babel}
\usepackage[a4paper,margin=2cm,top=3.4cm,bottom=2.8cm,headheight=1.4cm,footskip=1.3cm]{geometry}
\usepackage{amsmath,amssymb,amsfonts}
\usepackage{mathtools} % \xmapsto, \coloneqq... souvent utilisés par les modèles
\usepackage{cancel} % simplifications barrées
% \abs{z} (module, valeur absolue) et \norm{u} (norme), utilisés par les modèles
\providecommand{\abs}{}\let\abs\relax\DeclarePairedDelimiter{\abs}{\lvert}{\rvert}
\providecommand{\norm}{}\let\norm\relax\DeclarePairedDelimiter{\norm}{\lVert}{\rVert}
\usepackage[table]{xcolor} % [table] : fournit aussi \rowcolors (alternance de lignes)
\usepackage{pagecolor}
\usepackage{tikz}
\usetikzlibrary{arrows.meta,positioning,calc,shapes.geometric,shapes.misc,shapes.arrows,decorations.pathmorphing,decorations.pathreplacing,decorations.markings,patterns,shadows,mindmap,trees,fit,backgrounds,matrix,intersections,angles,quotes}
\usepackage{pgfplots}
\usepackage{pgf-pie} % diagrammes circulaires \pie (statistiques)
\pgfplotsset{compat=1.18}
\usepgfplotslibrary{fillbetween,groupplots} % aires entre courbes (calcul intégral)
\usepackage{enumitem}
\usepackage{fancyhdr}
% [most] charge toutes les bibliothèques tcolorbox (listings, minted, raster,
% documentation...) et gonfle inutilement la mémoire TeX sur un document long
% (~300 boîtes) : on ne charge que ce qui est réellement utilisé.
\usepackage[skins,breakable]{tcolorbox}
\usepackage{graphicx}
\usepackage{colortbl}
\usepackage{booktabs}
\usepackage{tabularx}
\usepackage{diagbox} % case d'en-tête diagonale des tableaux à double entrée
% Colonnes extensibles centrée / gauche / droite (C, L, R), souvent utilisées par les modèles
\newcolumntype{C}{>{\centering\arraybackslash}X}
\newcolumntype{L}{>{\raggedright\arraybackslash}X}
\newcolumntype{R}{>{\raggedleft\arraybackslash}X}
\usepackage{array}
\usepackage{multicol}
\usepackage{siunitx}
% parse-numbers=false : accepte aussi \SI{2,5 \times 10^{-3}}{...}
\sisetup{locale=FR,output-decimal-marker={,},group-minimum-digits=4,parse-numbers=false}
\DeclareSIUnit\ohm{\ensuremath{\symup{\Omega}}} % U+2126 absent de la police
\usepackage{qrcode}
% \degree : commande fréquemment utilisée par les modèles pour un angle ou une
% température en dehors de \SI{}{} (non fournie par les paquets chargés).
\providecommand{\degree}{^{\circ}}
% \qed (fin de démonstration), utilisé par les modèles sans amsthm
\providecommand{\qed}{\hfill$\square$}
% \barre : double barre des valeurs interdites dans un tableau de variations (tkz-tab)
\providecommand{\barre}{\ensuremath{\|}}
% environnement proof (amsthm non chargé) utilisé par les modèles
\ifdefined\proof\else\newenvironment{proof}[1][Démonstration]{\par\noindent\textit{#1.}\ }{\qed\par}\fi
\usepackage{lastpage}
\usepackage{hyperref}
\hypersetup{hidelinks}

% ============================================================
% Palette « Pop » OnBuch+ — mêmes hex que la classe P (plus_ui.dart)
% ============================================================
\definecolor{poporange}{HTML}{F59321}
\definecolor{popdark}{HTML}{E07A0C}
\definecolor{popcream}{HTML}{FFF6E8}
\definecolor{popink}{HTML}{1C1714}
\definecolor{poppurple}{HTML}{7A5AE0}
\definecolor{popgreen}{HTML}{2E9E6A}
\definecolor{poppink}{HTML}{E0457B}
\definecolor{popblue}{HTML}{2F7DE1}
\definecolor{popgold}{HTML}{C98A12}
\definecolor{popmuted}{HTML}{8A7F76}
\definecolor{popline}{HTML}{EADFD2}
\definecolor{popgray}{HTML}{8A7F76}
\colorlet{popgrey}{popgray}
% Alias tolérants (le modèle nomme parfois les couleurs différemment)
\colorlet{popwhite}{white}
\colorlet{popblack}{popink}
\colorlet{popred}{poppink}
\colorlet{popyellow}{popgold}
% Teintes claires pour fonds de tableaux / zones
\colorlet{popblueL}{popblue!10!white}
\colorlet{poporangeL}{poporange!14!white}
\colorlet{popgreenL}{popgreen!12!white}
\colorlet{poppurpleL}{poppurple!12!white}
\colorlet{poppinkL}{poppink!12!white}
\colorlet{popgoldL}{popgold!14!white}

\pagecolor{popcream}

% ============================================================
% Typographie
% ============================================================
\defaultfontfeatures{Ligatures=TeX}
\setmainfont{PlusJakartaSans-Medium.ttf}[Path=../../../fonts/,BoldFont=PlusJakartaSans-Bold.ttf]
% Maths en sans-serif (Fira Math) avec chiffres et lettres droites de Plus
% Jakarta, pour que les formules se fondent dans le texte.
\usepackage[math-style=ISO,bold-style=ISO]{unicode-math}
\setmathfont{FiraMath-Regular.otf}[Path=../../../fonts/]
\setmathfont{PlusJakartaSans-Medium.ttf}[Path=../../../fonts/,range={up/num,up/{Latin,latin},\mathup}]
\usepackage{icomma} % virgule décimale sans espace en mode math
% Caractères mathématiques Unicode absents des polices de texte (Plus Jakarta,
% Archivo Black) : rendus par leur équivalent mathématique, en texte comme en
% formule — sinon ils s'affichent en carré vide (ex. « Arithmétique dans ℤ »).
\usepackage{newunicodechar}
\newunicodechar{ℝ}{\ensuremath{\mathbb{R}}}
\newunicodechar{ℤ}{\ensuremath{\mathbb{Z}}}
\newunicodechar{ℂ}{\ensuremath{\mathbb{C}}}
\newunicodechar{ℕ}{\ensuremath{\mathbb{N}}}
\newunicodechar{ℚ}{\ensuremath{\mathbb{Q}}}
\newunicodechar{𝔻}{\ensuremath{\mathbb{D}}}
\newunicodechar{α}{\ensuremath{\alpha}}
\newunicodechar{β}{\ensuremath{\beta}}
\newunicodechar{γ}{\ensuremath{\gamma}}
\newunicodechar{δ}{\ensuremath{\delta}}
\newunicodechar{ε}{\ensuremath{\varepsilon}}
\newunicodechar{θ}{\ensuremath{\theta}}
\newunicodechar{λ}{\ensuremath{\lambda}}
\newunicodechar{μ}{\ensuremath{\mu}}
\newunicodechar{σ}{\ensuremath{\sigma}}
\newunicodechar{τ}{\ensuremath{\tau}}
\newunicodechar{φ}{\ensuremath{\varphi}}
\newunicodechar{ψ}{\ensuremath{\psi}}
\newunicodechar{ω}{\ensuremath{\omega}}
\newunicodechar{Σ}{\ensuremath{\Sigma}}
% majuscules produites par \MakeUppercase dans les titres (α-aminés -> Α)
\newunicodechar{Α}{\ensuremath{\alpha}}
\newunicodechar{Β}{\ensuremath{\beta}}
\newunicodechar{Γ}{\ensuremath{\Gamma}}
\newunicodechar{Δ}{\ensuremath{\Delta}}
\newunicodechar{Ε}{\ensuremath{\varepsilon}}
\newunicodechar{Θ}{\ensuremath{\Theta}}
\newunicodechar{Λ}{\ensuremath{\Lambda}}
\newunicodechar{Μ}{\ensuremath{\mu}}
\newunicodechar{Τ}{\ensuremath{\tau}}
\newunicodechar{Φ}{\ensuremath{\Phi}}
\newunicodechar{Ψ}{\ensuremath{\Psi}}
\newunicodechar{Ω}{\ensuremath{\Omega}}
\newunicodechar{↦}{\ensuremath{\mapsto}}
\newunicodechar{∈}{\ensuremath{\in}}
\newunicodechar{∉}{\ensuremath{\notin}}
\newunicodechar{∧}{\ensuremath{\wedge}}
\newunicodechar{⇔}{\ensuremath{\Leftrightarrow}}
\newunicodechar{⟺}{\ensuremath{\Longleftrightarrow}}
\newunicodechar{⇒}{\ensuremath{\Rightarrow}}
\newunicodechar{⟹}{\ensuremath{\Longrightarrow}}
\newunicodechar{∘}{\ensuremath{\circ}}
\newunicodechar{∪}{\ensuremath{\cup}}
\newunicodechar{∩}{\ensuremath{\cap}}
\newunicodechar{≡}{\ensuremath{\equiv}}
\newunicodechar{⊂}{\ensuremath{\subset}}
\newunicodechar{⊥}{\ensuremath{\perp}}
\newunicodechar{∃}{\ensuremath{\exists}}
\newunicodechar{∀}{\ensuremath{\forall}}
\newunicodechar{∛}{\ensuremath{\sqrt[3]{\,}}}
\newunicodechar{∜}{\ensuremath{\sqrt[4]{\,}}}
\newunicodechar{⃗}{\ensuremath{{}^{\to}}}
\newunicodechar{ⁿ}{\ifmmode^{n}\else\textsuperscript{\ensuremath{n}}\fi}
\newunicodechar{ˣ}{\ifmmode^{x}\else\textsuperscript{\ensuremath{x}}\fi}
\newunicodechar{ᵇ}{\ifmmode^{b}\else\textsuperscript{\ensuremath{b}}\fi}
\newunicodechar{ʳ}{\ifmmode^{r}\else\textsuperscript{\ensuremath{r}}\fi}
\newunicodechar{ᵘ}{\ifmmode^{u}\else\textsuperscript{\ensuremath{u}}\fi}
\newunicodechar{ʸ}{\ifmmode^{y}\else\textsuperscript{\ensuremath{y}}\fi}
\newunicodechar{ᵃ}{\ifmmode^{a}\else\textsuperscript{\ensuremath{a}}\fi}
\newunicodechar{ᵉ}{\ifmmode^{e}\else\textsuperscript{\ensuremath{e}}\fi}
\newunicodechar{ᵏ}{\ifmmode^{k}\else\textsuperscript{\ensuremath{k}}\fi}
\newunicodechar{ᵖ}{\ifmmode^{p}\else\textsuperscript{\ensuremath{p}}\fi}
\newunicodechar{ᴺ}{\ifmmode^{N}\else\textsuperscript{\ensuremath{N}}\fi}
\newunicodechar{ᶜ}{\ifmmode^{c}\else\textsuperscript{\ensuremath{c}}\fi}
\newunicodechar{ʰ}{\ifmmode^{h}\else\textsuperscript{\ensuremath{h}}\fi}
\newunicodechar{ᵠ}{\ifmmode^{q}\else\textsuperscript{\ensuremath{q}}\fi}
\newunicodechar{ₙ}{\ifmmode_{n}\else\textsubscript{\ensuremath{n}}\fi}
\newunicodechar{ₐ}{\ifmmode_{a}\else\textsubscript{\ensuremath{a}}\fi}
\newunicodechar{ᵢ}{\ifmmode_{i}\else\textsubscript{\ensuremath{i}}\fi}
\newunicodechar{ᵣ}{\ifmmode_{r}\else\textsubscript{\ensuremath{r}}\fi}
\newunicodechar{ₖ}{\ifmmode_{k}\else\textsubscript{\ensuremath{k}}\fi}
\newunicodechar{ᵦ}{\ifmmode_{\beta}\else\textsubscript{\ensuremath{\beta}}\fi}
\newunicodechar{ₓ}{\ifmmode_{x}\else\textsubscript{\ensuremath{x}}\fi}
\newfontfamily\popdisplay{ArchivoBlack-Regular.ttf}[Path=../../../fonts/]
\newfontfamily\poplogo{SpaceGrotesk-Bold.ttf}[Path=../../../fonts/]
\newfontfamily\popsemi{PlusJakartaSans-SemiBold.ttf}[Path=../../../fonts/]
\newfontfamily\popxbold{PlusJakartaSans-ExtraBold.ttf}[Path=../../../fonts/]
\newfontfamily\popxboldls{PlusJakartaSans-ExtraBold.ttf}[Path=../../../fonts/,LetterSpace=3.0]

\setlist[enumerate]{leftmargin=1.7em,itemsep=5pt,topsep=4pt}
\setlist[itemize]{leftmargin=1.7em,itemsep=4pt,topsep=4pt,label={\color{poporange}\textbullet}}
\setlength{\parindent}{0pt}
\setlength{\parskip}{5pt}
\renewcommand{\arraystretch}{1.25}

% pgfplots : style Pop par défaut
\pgfplotsset{
  every axis/.append style={axis line style={popink,line width=1pt},tick style={popink},
    grid style={popline,line width=0.5pt},label style={font=\small},tick label style={font=\footnotesize}},
  popcurve/.style={poporange,line width=1.8pt,no markers,smooth},
}
% Même style disponible dans un \draw TikZ simple (hors pgfplots)
\tikzset{popcurve/.style={poporange,line width=1.8pt}}

% ============================================================
% Pill / squiggle
% ============================================================
\newcommand{\poppill}[3]{%
  \tikz[baseline=-0.55ex]\node[draw=popink,line width=1.2pt,rounded corners=2.6pt,%
    fill=#2,inner xsep=6.5pt,inner ysep=3pt]{%
    \popxboldls\fontsize{7}{7}\selectfont\color{#3}\MakeUppercase{#1}};%
}
\newcommand{\popsquiggle}[1]{%
  \tikz[baseline=-0.4ex]{
    \draw[#1,line width=2.6pt,line cap=round]
      plot [smooth] coordinates {(0,0.09) (0.34,0.01) (0.68,0.21) (1.02,0.01) (1.36,0.10)};
  }%
}
% Mot-clé surligné (marqueur orange sous le texte)
\newcommand{\cle}[1]{{\popxbold\color{popdark}#1}}

% ============================================================
% En-tête / pied
% ============================================================
\pagestyle{fancy}
\fancyhf{}
\renewcommand{\headrulewidth}{0pt}
\renewcommand{\footrulewidth}{0pt}
\lhead{\raisebox{-0.15ex}{\poplogo\fontsize{13}{13}\selectfont\color{popink}OnBuch\color{poporange}+}}
\rhead{\poppill{\DOCMATIERE\ \textbullet\ \DOCNIVEAU\ \textbullet\ Leçon \DOCLECON}{white}{popink}}
\lfoot{\popxbold\fontsize{7.6}{7.6}\selectfont\color{popmuted}\MakeUppercase{Cours signé OnBuch+}\ \textbullet\ {\popsemi\DOCTITRE}}
\rfoot{\tikz[baseline=-0.6ex]\node[rounded corners=8.5pt,fill=popink,minimum height=17pt,minimum width=17pt,inner xsep=5pt,inner sep=0]{\popxbold\fontsize{8}{8}\selectfont\color{popcream}\thepage\,/\,\pageref*{LastPage}};}

% ============================================================
% Titres
% ============================================================
\newcommand{\chaptitre}[2]{%
  \begin{tcolorbox}[enhanced,colback=poporange,colframe=popink,boxrule=2.6pt,arc=11pt,
      drop shadow=popink,left=15pt,right=15pt,top=12pt,bottom=11pt,before skip=3pt,after skip=18pt]
    \poppill{#2}{popink}{popcream}\par\vspace{9pt}
    {\raggedright\hyphenpenalty=10000\exhyphenpenalty=10000\popdisplay\fontsize{22}{24}\selectfont\color{popcream}\MakeUppercase{#1}\par}
  \end{tcolorbox}
}
% Section numérotée : pastille orange avec le numéro + titre Archivo
\newcounter{popsec}
\newcommand{\coursec}[1]{%
  \stepcounter{popsec}\setcounter{popsub}{0}%
  \par\vspace{16pt}\noindent
  \begin{minipage}{\linewidth}
  \tikz[baseline=-0.7ex]\node[draw=popink,line width=1.6pt,fill=poporange,rounded corners=4pt,minimum size=22pt,inner sep=0]{\popdisplay\fontsize{12}{12}\selectfont\color{popcream}\thepopsec};%
  \hspace{8pt}{\popdisplay\fontsize{15}{17}\selectfont\color{popink}\MakeUppercase{#1}}\par
  \vspace{4pt}\noindent{\color{poporange}\rule{36pt}{3.6pt}}
  \end{minipage}\par\nopagebreak\vspace{8pt}
}
\newcounter{popsub}
\newcommand{\courssub}[1]{%
  \stepcounter{popsub}%
  \par\vspace{9pt}\noindent{\popxbold\fontsize{11.5}{13}\selectfont\color{popink}{\color{poporange}\thepopsec.\thepopsub}\hspace{6pt}#1}\par\nopagebreak\vspace{4pt}
}

% ============================================================
% Boîtes pédagogiques — même recette : fond blanc, bordure épaisse colorée,
% ombre dure encre, badge plein. Titre optionnel [#1].
% ============================================================
\tcbset{popbase/.style={breakable,enhanced,colback=white,boxrule=2.3pt,arc=9pt,
  shadow={3pt}{-3pt}{0pt}{popink},left=14pt,right=14pt,top=11pt,bottom=11pt,before skip=11pt,after skip=15pt,
  fontupper=\popsemi\fontsize{10.6}{14.5}\selectfont\color{popink},
  fontlower=\popsemi\fontsize{10.6}{14.5}\selectfont\color{popink}}}
% Coche dessinée (le glyphe ✓ n'existe pas dans les polices du document)
\newcommand{\popcheck}[1]{\tikz[baseline=-0.5ex]\draw[#1,line width=1.6pt,line cap=round,line join=round](-0.13,0.0)--(-0.04,-0.09)--(0.13,0.1);}
\newcommand{\popboxhead}[3]{\poppill{#1}{#2}{white}\ifx\relax#3\relax\else\hspace{6pt}{\popxbold\fontsize{10.6}{12}\selectfont\color{popink}#3}\fi\par\vspace{7pt}}

\newtcolorbox{aretenir}[1][]{popbase,colframe=popgreen,before upper={\popboxhead{À retenir}{popgreen}{#1}}}
\newtcolorbox{exemplebox}[1][]{popbase,colframe=poporange,before upper={\popboxhead{Exemple}{poporange}{#1}}}
\newtcolorbox{definition}[1][]{popbase,colframe=popblue,before upper={\popboxhead{Définition}{popblue}{#1}}}
\newtcolorbox{propriete}[1][]{popbase,colframe=poppurple,before upper={\popboxhead{Propriété}{poppurple}{#1}}}
\newtcolorbox{methode}[1][]{popbase,colframe=popgold,before upper={\popboxhead{Méthode}{popgold}{#1}}}
\newtcolorbox{attention}[1][]{popbase,colframe=poppink,before upper={\popboxhead{Attention}{poppink}{#1}}}
\newtcolorbox{experience}[1][]{popbase,colframe=popblue,colback=popblueL,before upper={\popboxhead{Expérience}{popblue}{#1}}}
% « Le savais-tu ? » : carte violette pleine (contexte, culture, Cameroun)
\newtcolorbox{savaistu}[1][]{popbase,colback=poppurple,colframe=popink,
  fontupper=\popsemi\fontsize{10.4}{14.2}\selectfont\color{white},
  before upper={\poppill{Le savais-tu\,?}{white}{poppurple}\ifx\relax#1\relax\else\hspace{6pt}{\popxbold\color{white}#1}\fi\par\vspace{7pt}}}
% Exercice résolu : énoncé en haut, \tcblower puis la solution sur fond crème
\newcounter{popexo}\newcounter{popexr}
\newtcolorbox{exoresolu}[1][]{popbase,colframe=poporange,colbacklower=popcream,
  segmentation style={popink,line width=1.2pt,dashed},
  before upper={\stepcounter{popexr}\popboxhead{Exercice résolu \thepopexr}{poporange}{#1}},
  before lower={\poppill{Solution}{popink}{popcream}\par\vspace{6pt}}}
% Exercice (non résolu en place) — la correction est dans la partie Corrigés
\newtcolorbox{exercice}[1][]{popbase,colframe=popink,
  before upper={\stepcounter{popexo}\popboxhead{Exercice \thepopexo}{popink}{#1}}}
% Corrigé
\newtcolorbox{corrige}[1][]{popbase,colframe=popgreen,colback=popgreenL,
  before upper={\popboxhead{Corrigé}{popgreen}{#1}}}
% Objectifs / prérequis / bilan
\newtcolorbox{objectifs}{popbase,colframe=popink,colback=poporangeL,
  before upper={\popboxhead{Objectifs de la leçon}{popink}{}}}
\newtcolorbox{prerequis}{popbase,colframe=popink,colback=popblueL,
  before upper={\popboxhead{Prérequis}{popblue}{}}}
\newtcolorbox{bilan}{popbase,colframe=popink,colback=white,boxrule=2.8pt,
  before upper={\popboxhead{Fiche bilan}{poporange}{L'essentiel à connaître pour l'examen}}}
% Figure encadrée avec légende
\newtcolorbox{popfigure}[1][]{enhanced,breakable=false,colback=white,colframe=popline,boxrule=1.4pt,arc=8pt,
  left=10pt,right=10pt,top=10pt,bottom=8pt,before skip=10pt,after skip=12pt,halign=center,
  lower separated=false,halign lower=center,
  fontlower=\popsemi\fontsize{9}{11.5}\selectfont\color{popmuted},
  #1}
\newcommand{\legende}[1]{\par\vspace{6pt}{\popsemi\fontsize{9}{11.5}\selectfont\color{popmuted}#1}}

% ============================================================
% Signature OnBuch+
% ============================================================
\newcommand{\poplogoinline}[1]{{\poplogo\fontsize{#1}{#1}\selectfont\color{popink}OnBuch\color{poporange}+}}
\newcommand{\signatureonbuch}{%
  \begin{tcolorbox}[enhanced,colback=popink,colframe=popink,boxrule=2.2pt,arc=10pt,
      drop shadow=poporange,left=14pt,right=14pt,top=10pt,bottom=10pt,before skip=0pt,after skip=0pt]
    \tikz[baseline=-0.7ex]\node[circle,fill=popgreen,draw=popcream,line width=1.3pt,minimum size=22pt,inner sep=0]{\popcheck{white}};%
    \hspace{9pt}\begin{minipage}[c]{0.62\linewidth}
      {\popxbold\fontsize{12}{13}\selectfont\color{popcream}Signé\ \textbullet\ {\poplogo OnBuch\color{poporange}+}}\par\vspace{2pt}
      {\raggedright\popsemi\fontsize{8}{10.5}\selectfont\color{popline}\MakeUppercase{Équipe pédagogique OnBuch+}\\\MakeUppercase{Conforme au programme officiel MINESEC}\par}
    \end{minipage}\hfill
    {\poplogo\fontsize{22}{22}\selectfont\color{popcream}OnBuch\color{poporange}+}
  \end{tcolorbox}%
}

% ============================================================
% Page de garde
% #1 titre · #2 matière · #3 niveau · #4 module · #5 n° leçon · #6 durée ·
% #7 description / objectifs (texte ou itemize)
% ============================================================
\newcommand{\pagegarde}[7]{%
  \thispagestyle{empty}
  \newgeometry{a4paper,margin=1.7cm,top=1.6cm,bottom=1.4cm}%
  \begin{tikzpicture}[remember picture,overlay]
    \fill[poporange!18] ([xshift=-3.2cm,yshift=-3.6cm]current page.north east) circle (4.6cm);
    \fill[poppurple!14] ([xshift=2.4cm,yshift=7.6cm]current page.south west) circle (3.8cm);
  \end{tikzpicture}%
  {\poplogo\fontsize{30}{30}\selectfont\color{popink}OnBuch\color{poporange}+}\hfill
  \raisebox{0.6ex}{\poppill{Cours signé \textbullet\ Premium}{popink}{popcream}}\par
  \vspace{1pt}\popsquiggle{poppurple}\par
  \vspace{8pt}
  \begin{tcolorbox}[enhanced,colback=poporange,colframe=popink,boxrule=2.8pt,arc=13pt,
      drop shadow=popink,left=18pt,right=18pt,top=12pt,bottom=13pt,after skip=10pt]
    \poppill{#2 \textbullet\ #3}{popink}{popcream}\hspace{5pt}\poppill{#4}{white}{popink}\par\vspace{8pt}
    {\popdisplay\fontsize{14}{14}\selectfont\color{popink}LEÇON #5}\par\vspace{4pt}
    {\raggedright\hyphenpenalty=10000\exhyphenpenalty=10000\popdisplay\fontsize{25}{27}\selectfont\color{popcream}\MakeUppercase{#1}\par}
    \vspace{8pt}
    {\popxbold\fontsize{9.5}{9.5}\selectfont\color{popink}\MakeUppercase{Durée conseillée : #6}}
  \end{tcolorbox}
  \begin{tcolorbox}[enhanced,colback=white,colframe=popink,boxrule=2.2pt,arc=10pt,
      drop shadow=popink,left=16pt,right=16pt,top=10pt,bottom=10pt,after skip=8pt]
    \poppill{Dans cette leçon}{poppurple}{white}\par\vspace{5pt}
    \popsemi\fontsize{9.3}{12.6}\selectfont\color{popink}#7
  \end{tcolorbox}
  \noindent\begin{minipage}[t]{0.487\linewidth}
    \begin{tcolorbox}[enhanced,colback=popgreen,colframe=popink,boxrule=2.2pt,arc=10pt,
        drop shadow=popink,left=11pt,right=11pt,top=10pt,bottom=10pt,height=3.1cm,valign=center]
      \tcbox[colback=white,colframe=popink,boxrule=1.3pt,arc=5pt,boxsep=0pt,nobeforeafter,
        left=3pt,right=3pt,top=3pt,bottom=3pt,valign=center]{\qrcode[height=1.9cm]{https://play.google.com/store/apps/details?id=com.luvvix.onbuch&hl=fr}}%
      \hspace{8pt}\begin{minipage}[c]{0.5\linewidth}
        {\popdisplay\fontsize{11}{12.5}\selectfont\color{white}\MakeUppercase{Télécharge OnBuch}}\par\vspace{4pt}
        {\raggedright\popsemi\fontsize{8.4}{11}\selectfont\color{white}Gratuit \textbullet\ Google Play\\Tuteur IA, annales, concours, résultats\par}
      \end{minipage}
    \end{tcolorbox}
  \end{minipage}\hfill
  \begin{minipage}[t]{0.487\linewidth}
    \begin{tcolorbox}[enhanced,colback=poppurple,colframe=popink,boxrule=2.2pt,arc=10pt,
        drop shadow=popink,left=11pt,right=11pt,top=10pt,bottom=10pt,height=3.1cm,valign=center]
      \tikz[baseline=-0.7ex]\node[circle,fill=white,draw=popink,line width=1.3pt,minimum size=1.6cm,inner sep=0]{\popdisplay\fontsize{24}{24}\selectfont\color{poppurple}+};%
      \hspace{8pt}\begin{minipage}[c]{0.6\linewidth}
        {\popdisplay\fontsize{11}{12.5}\selectfont\color{white}\MakeUppercase{Passe à OnBuch+}}\par\vspace{4pt}
        {\raggedright\popsemi\fontsize{8.4}{11}\selectfont\color{white}Tous les cours signés, Tuteur IA illimité, fiches PDF — onglet OnBuch+ de l'app\par}
      \end{minipage}
    \end{tcolorbox}
  \end{minipage}
  \vspace{8pt}
  \popcontenu
  \vfill
  \signatureonbuch
  \restoregeometry
  \newpage
}

% Rangée « Ce cours contient » (page de garde)
\newcommand{\poptile}[3]{% #1 couleur #2 gros texte #3 légende
  \begin{tcolorbox}[enhanced,colback=#1,colframe=popink,boxrule=2pt,arc=9pt,shadow={2.5pt}{-2.5pt}{0pt}{popink},
      left=6pt,right=6pt,top=9pt,bottom=9pt,halign=center,valign=center,height=1.75cm,width=\linewidth,nobeforeafter]
    {\popdisplay\fontsize{12.5}{13}\selectfont\color{white}\MakeUppercase{#2}}\par\vspace{3pt}
    {\centering\popsemi\fontsize{7.8}{9.5}\selectfont\color{white}#3\par}
  \end{tcolorbox}}
\newcommand{\popcontenu}{%
  \par\noindent{\popxboldls\fontsize{7.5}{7.5}\selectfont\color{popmuted}CE COURS SIGNÉ CONTIENT}\par\vspace{7pt}
  \noindent\begin{minipage}[t]{0.235\linewidth}\poptile{popblue}{Cours}{complet, illustré, pas à pas}\end{minipage}\hfill
  \begin{minipage}[t]{0.235\linewidth}\poptile{poporange}{Méthodes}{et exercices résolus}\end{minipage}\hfill
  \begin{minipage}[t]{0.235\linewidth}\poptile{poppink}{Exercices}{type examen + corrigés}\end{minipage}\hfill
  \begin{minipage}[t]{0.235\linewidth}\poptile{popgreen}{Fiche bilan}{l'essentiel à retenir}\end{minipage}\par}

% Sommaire
\newtcolorbox{sommaire}{enhanced,colback=white,colframe=popink,boxrule=2.2pt,arc=10pt,
  drop shadow=popink,left=18pt,right=18pt,top=13pt,bottom=13pt,before skip=6pt,after skip=20pt,
  before upper={\poppill{Sommaire}{poppurple}{white}\par\vspace{9pt}\popsemi\fontsize{10.6}{14.5}\selectfont\color{popink}}}

% Signature de fin de cours — poussée tout en bas de la dernière page
\newcommand{\findecours}{%
  \par\vfill\nopagebreak
  \begin{center}\popsquiggle{poporange}\end{center}\vspace{4pt}
  \signatureonbuch
}
\AtBeginDocument{\def\llbracket{\mathopen{[\mkern-3.5mu[}}\def\rrbracket{\mathclose{]\mkern-3.5mu]}}} % intervalles d'entiers (glyphe ⟦ absent de la police)
% Glyphes absents de Fira Math : redéfinis à partir de glyphes disponibles
\AtBeginDocument{%
  \def\setminus{\mathbin{\backslash}}\let\smallsetminus\setminus
  \def\vdots{\mathord{\vbox{\baselineskip3.2pt\lineskiplimit0pt\kern4pt\hbox{.}\hbox{.}\hbox{.}}}}%
  \def\ddots{\mathinner{\mkern1mu\raise6.5pt\hbox{.}\mkern2mu\raise3.6pt\hbox{.}\mkern2mu\raise0.7pt\hbox{.}\mkern1mu}}%
  \def\triangle{\mathord{\scalebox{0.9}{$\symup{\Delta}$}}}%
  \def\nabla{\mathord{\raisebox{\depth}{\scalebox{1}[-1]{$\symup{\Delta}$}}}}%
  \def\checkmark{\popcheck{popdark}}%
  \def\star{\mathbin{\ast}}\def\bigstar{\mathord{\ast}}%
  \def\diamond{\mathbin{\scalebox{0.6}{$\Diamond$}}}%
  \def\bigcup{\mathop{\mathchoice{\raisebox{-0.3em}{\scalebox{1.7}{$\cup$}}}{\scalebox{1.25}{$\cup$}}{\cup}{\cup}}\displaylimits}%
  \def\bigcap{\mathop{\mathchoice{\raisebox{-0.3em}{\scalebox{1.7}{$\cap$}}}{\scalebox{1.25}{$\cap$}}{\cap}{\cap}}\displaylimits}%
  \def\lesseqqgtr{\mathrel{\lessgtr}}\def\bigoplus{\mathop{\oplus}\displaylimits}%
}
\providecommand{\ssi}{si et seulement si} % abréviation fréquente
\AtBeginDocument{\providecommand{\wideparen}{\overparen}} % arc de cercle

%% FIGFIX-BEGIN v5 — correctifs globaux des figures (docs/onbuch-generation/tools/figfix)
\usepackage{adjustbox}\usepackage{etoolbox}\tcbuselibrary{raster}
% toute figure TikZ/pgfplots plus large que la ligne est réduite à la largeur disponible
\BeforeBeginEnvironment{tikzpicture}{\begin{adjustbox}{max width=\linewidth}}
\AfterEndEnvironment{tikzpicture}{\end{adjustbox}}
% légendes pgfplots lisibles par-dessus les courbes
\pgfplotsset{every axis/.append style={legend style={fill=white,fill opacity=0.92,text opacity=1,draw=black!15,inner sep=3pt}}}
% étiquettes d'arêtes (syntaxe edge["..."]) avec fond blanc
\tikzset{every edge quotes/.append style={fill=white,inner sep=1.2pt}}
%% FIGFIX-END
\begin{document}

\pagegarde{Équations et inéquations polynomiales}{Mathématiques}{1ère D}{Relations et opérations fondamentales dans l'ensemble des nombres réels}{N01}{12 h}{Cette leçon outille l'élève pour résoudre dans ℝ les équations et inéquations du second degré, les systèmes s'y ramenant, les équations et inéquations de degré 3 par factorisation, ainsi que les équations et inéquations irrationnelles simples. Elle consolide les acquis de Seconde et constitue le socle algébrique indispensable pour l'étude des fonctions en Terminale.
\vspace{4pt}
\begin{multicols}{2}\small
\begin{itemize}
\item À la fin de cette leçon, je sais résoudre une équation du second degré par la forme canonique et par le discriminant
\item À la fin de cette leçon, je sais discuter l'existence des solutions d'une équation du second degré dépendant d'un paramètre
\item À la fin de cette leçon, je sais factoriser un trinôme et dresser son tableau de signes pour résoudre une inéquation du second degré
\item À la fin de cette leçon, je sais résoudre un système de deux équations à deux inconnues se ramenant à une équation du second degré (somme-produit, substitution)
\item À la fin de cette leçon, je sais vérifier qu'un réel est zéro d'un polynôme et factoriser un polynôme de degré 3 par division euclidienne ou par coefficients indéterminés
\item À la fin de cette leçon, je sais résoudre des équations et inéquations de degré 3 à l'aide d'un tableau de signes
\end{itemize}\end{multicols}}

\begin{sommaire}
\begin{multicols}{2}
\begin{enumerate}
\item Équations du second degré : forme canonique et discriminant
\item Signe du trinôme et inéquations du second degré
\item Systèmes se ramenant à une équation du second degré
\item Polynômes de degré 3 : zéros et factorisation
\item Signe d'un polynôme de degré 3 et inéquations
\item Équations et inéquations irrationnelles
\item Activité d'intégration
\item Méthodes et exercices résolus
\item Exercices
\item Corrigés des exercices
\item Fiche bilan
\end{enumerate}
\end{multicols}
\end{sommaire}

\begin{objectifs}
\begin{itemize}
\item À la fin de cette leçon, je sais résoudre une équation du second degré par la forme canonique et par le discriminant
\item À la fin de cette leçon, je sais discuter l'existence des solutions d'une équation du second degré dépendant d'un paramètre
\item À la fin de cette leçon, je sais factoriser un trinôme et dresser son tableau de signes pour résoudre une inéquation du second degré
\item À la fin de cette leçon, je sais résoudre un système de deux équations à deux inconnues se ramenant à une équation du second degré (somme-produit, substitution)
\item À la fin de cette leçon, je sais vérifier qu'un réel est zéro d'un polynôme et factoriser un polynôme de degré 3 par division euclidienne ou par coefficients indéterminés
\item À la fin de cette leçon, je sais résoudre des équations et inéquations de degré 3 à l'aide d'un tableau de signes
\item À la fin de cette leçon, je sais résoudre des équations et inéquations irrationnelles du type √(ax+b) = cx+d en contrôlant les conditions de validité
\end{itemize}
\end{objectifs}

\begin{prerequis}
\begin{itemize}
\item Calcul algébrique de Seconde : développement, factorisation, identités remarquables
\item Résolution d'équations et inéquations du premier degré dans ℝ
\item Notion de racine carrée et règles de calcul sur les radicaux
\item Tableaux de signes des expressions du premier degré et des produits
\item Résolution de systèmes de deux équations linéaires à deux inconnues
\end{itemize}
\end{prerequis}

\newpage

\coursec{Équations du second degré : forme canonique et discriminant}

\textbf{Situation d'accroche.} Monsieur Ngo Bassa, menuisier à Douala, doit découper une plaque de contreplaqué rectangulaire dont le périmètre est $14$ m et l'aire $12$ m$^2$ : quelles dimensions choisir ? Si l'on note $L$ et $\ell$ les dimensions, on obtient $L+\ell=7$ et $L\ell=12$, ce qui conduit à l'équation $x^2-7x+12=0$ : une \emph{équation du second degré}. Par ailleurs, une coopérative de Bafoussam modélise son bénéfice (en milliers de FCFA) par $B(x)=-x^2+8x-7$ selon le nombre $x$ de tonnes vendues : pour quelles quantités le bénéfice est-il positif ? Il s'agit d'une \emph{inéquation du second degré}. Enfin, un réservoir cubique doit contenir exactement $2$ m$^3$ de plus qu'un cube de côté $x$ : on tombe sur l'équation $x^3+2=V$, c'est-à-dire une équation de \emph{degré 3}. Ces trois problèmes se résolvent avec les équations et inéquations polynomiales, objets de cette leçon.

\textbf{Problématique.} Comment résoudre de manière systématique une équation du type $ax^2+bx+c=0$ ? Existe-t-il un « test » permettant de savoir, avant tout calcul, combien de solutions elle possède ? C'est toute la puissance de la \cle{forme canonique} et du \cle{discriminant} que nous allons construire ensemble.

\courssub{Équations du second degré : définition et premiers exemples}

\begin{definition}[Équation du second degré]
On appelle \cle{équation du second degré} à une inconnue réelle $x$ toute équation qui peut s'écrire sous la forme
\[ ax^2+bx+c=0 \quad \text{avec } a,\ b,\ c \in \mathbb{R} \text{ et } a\neq 0. \]
L'expression $ax^2+bx+c$ est appelée \cle{trinôme du second degré}. Les réels $a$, $b$, $c$ sont les \emph{coefficients} du trinôme.
\end{definition}

\textbf{Intuition.} Pourquoi exiger $a\neq 0$ ? Tout simplement parce que si $a=0$, le terme en $x^2$ disparaît et l'équation devient $bx+c=0$, une équation du \emph{premier} degré que tu sais résoudre depuis la classe de Troisième. Le « second degré » doit son nom à la présence effective du terme en $x^2$.

\begin{attention}[La condition $a\neq 0$ est indispensable]
Avant de parler de trinôme du second degré, de forme canonique ou de discriminant, \textbf{vérifie toujours que le coefficient de $x^2$ est non nul}. C'est une erreur fréquente au Probatoire, surtout dans les équations à paramètre : pour l'équation $(m-1)x^2+2x+1=0$, le cas $m=1$ doit être traité \emph{à part}, car l'équation devient alors du premier degré. Appliquer les formules du second degré avec $a=0$ conduirait à des divisions par zéro.
\end{attention}

\begin{exemplebox}[Reconnaître une équation du second degré]
\begin{itemize}
\item $2x^2-5x+3=0$ est une équation du second degré : $a=2$, $b=-5$, $c=3$.
\item $x^2-9=0$ en est une aussi, avec $a=1$, $b=0$, $c=-9$ (on dit que le trinôme est \emph{incomplet}).
\item $3x+7=0$ n'est \textbf{pas} du second degré : le coefficient de $x^2$ est nul.
\item L'équation du menuisier $x^2-7x+12=0$ : $a=1$, $b=-7$, $c=12$.
\end{itemize}
\end{exemplebox}

\courssub{La forme canonique : complétion du carré}

\textbf{Intuition.} Résoudre $x^2=4$ est facile : $x=2$ ou $x=-2$. Résoudre $(x+3)^2=4$ l'est presque autant : $x+3=2$ ou $x+3=-2$, d'où $x=-1$ ou $x=-5$. L'idée géniale consiste donc à \emph{transformer} n'importe quel trinôme $ax^2+bx+c$ pour faire apparaître un carré parfait du type $(x+\alpha)^2$. Cette technique s'appelle la \cle{complétion du carré}.

\begin{experience}[Démonstration guidée : mise sous forme canonique]
Partons du trinôme $ax^2+bx+c$ avec $a\neq 0$ et suivons les étapes.
\begin{enumerate}
\item \textbf{Factoriser par $a$} (possible car $a\neq 0$) :
\[ ax^2+bx+c = a\left(x^2+\frac{b}{a}x\right)+c. \]
\item \textbf{Faire apparaître un carré.} On reconnaît dans $x^2+\dfrac{b}{a}x$ le début du développement de $\left(x+\dfrac{b}{2a}\right)^2$. En effet :
\[ \left(x+\frac{b}{2a}\right)^2 = x^2+\frac{b}{a}x+\frac{b^2}{4a^2}. \]
Il manque le terme $\dfrac{b^2}{4a^2}$ : on l'\emph{ajoute puis on le retranche} pour ne rien changer :
\[ x^2+\frac{b}{a}x = \left(x+\frac{b}{2a}\right)^2-\frac{b^2}{4a^2}. \]
\item \textbf{Regrouper} :
\[ ax^2+bx+c = a\left(x+\frac{b}{2a}\right)^2-\frac{b^2}{4a}+c = a\left(x+\frac{b}{2a}\right)^2-\frac{b^2-4ac}{4a}. \]
\end{enumerate}
On a donc démontré la \cle{forme canonique} :
\[ \boxed{\ ax^2+bx+c = a\left(x+\frac{b}{2a}\right)^2-\frac{b^2-4ac}{4a}\ } \]
\end{experience}

\begin{aretenir}[Forme canonique]
Tout trinôme $ax^2+bx+c$ ($a\neq 0$) s'écrit de manière unique sous la \cle{forme canonique} :
\[ ax^2+bx+c = a\left(x+\frac{b}{2a}\right)^2-\frac{b^2-4ac}{4a}. \]
Le réel $-\dfrac{b}{2a}$ est l'abscisse du sommet de la parabole d'équation $y=ax^2+bx+c$.
\end{aretenir}

\begin{exemplebox}[Mise sous forme canonique à la main]
Mettons $2x^2-5x+3$ sous forme canonique.
\begin{align*}
2x^2-5x+3 &= 2\left(x^2-\frac{5}{2}x\right)+3 = 2\left[\left(x-\frac{5}{4}\right)^2-\frac{25}{16}\right]+3\\
&= 2\left(x-\frac{5}{4}\right)^2-\frac{25}{8}+3 = 2\left(x-\frac{5}{4}\right)^2-\frac{1}{8}.
\end{align*}
Vérification : $b^2-4ac=25-24=1$ et $-\dfrac{\Delta}{4a}=-\dfrac{1}{8}$ : cohérent.
\end{exemplebox}

\courssub{Le discriminant et le théorème de résolution}

Dans la forme canonique, tout le jeu se concentre sur la quantité $b^2-4ac$ : c'est elle qui décide du nombre de solutions. Elle mérite donc un nom.

\begin{definition}[Discriminant]
Soit $ax^2+bx+c$ un trinôme du second degré ($a\neq 0$). On appelle \cle{discriminant} du trinôme le réel noté $\Delta$ (lire « delta ») défini par :
\[ \Delta = b^2-4ac. \]
\end{definition}

\begin{propriete}[Théorème de résolution de $ax^2+bx+c=0$]
Soit l'équation $ax^2+bx+c=0$ avec $a\neq 0$, de discriminant $\Delta=b^2-4ac$.
\begin{itemize}
\item Si $\Delta>0$ : l'équation admet \textbf{deux solutions réelles distinctes} :
\[ x_1=\frac{-b-\sqrt{\Delta}}{2a} \qquad \text{et} \qquad x_2=\frac{-b+\sqrt{\Delta}}{2a}. \]
\item Si $\Delta=0$ : l'équation admet \textbf{une unique solution} (dite \emph{racine double}) : $x_0=-\dfrac{b}{2a}$.
\item Si $\Delta<0$ : l'équation \textbf{n'admet aucune solution réelle} : $S=\varnothing$.
\end{itemize}
\emph{Cette démonstration est exigible au Probatoire : elle repose entièrement sur la forme canonique.}
\end{propriete}

\begin{experience}[Démonstration du théorème à partir de la forme canonique]
Grâce à la forme canonique, et comme $a\neq 0$ :
\[ ax^2+bx+c=0 \iff a\left(x+\frac{b}{2a}\right)^2=\frac{\Delta}{4a} \iff \left(x+\frac{b}{2a}\right)^2=\frac{\Delta}{4a^2}. \]
Le membre de gauche est un carré, donc positif ou nul ; le dénominateur $4a^2$ est strictement positif. Tout dépend donc du signe de $\Delta$.
\begin{enumerate}
\item \textbf{Si $\Delta>0$} : $\dfrac{\Delta}{4a^2}=\left(\dfrac{\sqrt{\Delta}}{2a}\right)^2>0$, donc
\[ x+\frac{b}{2a}=\frac{\sqrt{\Delta}}{2a} \quad \text{ou} \quad x+\frac{b}{2a}=-\frac{\sqrt{\Delta}}{2a}, \]
ce qui donne les deux solutions $x_1=\dfrac{-b-\sqrt{\Delta}}{2a}$ et $x_2=\dfrac{-b+\sqrt{\Delta}}{2a}$, distinctes car $\sqrt{\Delta}\neq 0$.
\item \textbf{Si $\Delta=0$} : l'équation devient $\left(x+\dfrac{b}{2a}\right)^2=0$, d'où l'unique solution $x_0=-\dfrac{b}{2a}$.
\item \textbf{Si $\Delta<0$} : le membre de droite $\dfrac{\Delta}{4a^2}$ est strictement négatif, alors que le membre de gauche est un carré, donc positif ou nul : c'est impossible. Il n'y a aucune solution réelle.
\end{enumerate}
\end{experience}

\textbf{Interprétation graphique.} Les solutions de $ax^2+bx+c=0$ sont les abscisses des points d'intersection de la parabole d'équation $y=ax^2+bx+c$ avec l'axe des abscisses. Le signe de $\Delta$ traduit la position de la parabole par rapport à cet axe.

\begin{popfigure}
\begin{center}
\begin{tikzpicture}[scale=0.85]
\begin{axis}[width=0.48\linewidth,axis lines=middle,xlabel=$x$,ylabel=$y$,xtick=\empty,ytick=\empty,title={$\Delta>0$ : deux intersections},title style={font=\small}]
\addplot[popcurve,domain=-1.2:3.2,samples=200]{x^2-2*x-1};
\addplot[mark=*,only marks,mark size=1.5pt,poporange] coordinates {(-0.414,0) (2.414,0)};
\end{axis}
\end{tikzpicture}\hfill
\begin{tikzpicture}[scale=0.85]
\begin{axis}[width=0.48\linewidth,axis lines=middle,xlabel=$x$,ylabel=$y$,xtick=\empty,ytick=\empty,title={$\Delta=0$ : tangente à l'axe},title style={font=\small}]
\addplot[popcurve,domain=-1.2:3.2,samples=200]{(x-1)^2};
\addplot[mark=*,only marks,mark size=1.5pt,poporange] coordinates {(1,0)};
\end{axis}
\end{tikzpicture}

\vspace{0.4cm}
\begin{tikzpicture}[scale=0.85]
\begin{axis}[width=0.48\linewidth,axis lines=middle,xlabel=$x$,ylabel=$y$,xtick=\empty,ytick=\empty,title={$\Delta<0$ : aucune intersection},title style={font=\small}]
\addplot[popcurve,domain=-1.2:3.2,samples=200]{(x-1)^2+1};
\end{axis}
\end{tikzpicture}
\end{center}
\legende{Position de la parabole $y=ax^2+bx+c$ selon le signe du discriminant $\Delta$.}
\end{popfigure}

\begin{popfigure}
\begin{center}
\begin{tikzpicture}[
  every node/.style={font=\small},
  start/.style={rectangle,rounded corners,draw=popblue,fill=popblueL,thick,text width=4.2cm,align=center},
  dec/.style={diamond,draw=poppurple,fill=poppurpleL,thick,aspect=2,align=center, inner sep=2pt},
  res/.style={rectangle,rounded corners,draw=popgreen,fill=popgreenL,thick,text width=4.2cm,align=center},
  fleche/.style={-{Stealth},thick,popdark}
]
\matrix[row sep=1cm, column sep=0.5cm] {
  & \node[start] (dep) {Équation $ax^2+bx+c=0$,\ $a\neq 0$ : calculer $\Delta=b^2-4ac$}; & \\
  & \node[dec] (test) {Signe de $\Delta$}; & \\
  \node[res] (pos) {$\Delta>0$ : deux solutions $x_{1,2}=\dfrac{-b\pm\sqrt{\Delta}}{2a}$}; &
  \node[res] (nul) {$\Delta=0$ : une solution double $x_0=-\dfrac{b}{2a}$}; &
  \node[res] (neg) {$\Delta<0$ : aucune solution réelle, $S=\varnothing$}; \\
};
\draw[fleche] (dep) -- (test);
\draw[fleche] (test) -- (pos) node[midway, above left] {$\Delta>0$};
\draw[fleche] (test) -- (nul) node[midway, right] {$\Delta=0$};
\draw[fleche] (test) -- (neg) node[midway, above right] {$\Delta<0$};
\end{tikzpicture}
\end{center}
\legende{Organigramme de résolution d'une équation du second degré selon le signe de $\Delta$.}
\end{popfigure}

\begin{exoresolu}[Résolution de $2x^2-5x+3=0$]
Résoudre dans $\mathbb{R}$ l'équation $2x^2-5x+3=0$.
\tcblower
\textbf{Étape 1 : identifier les coefficients.} $a=2$, $b=-5$, $c=3$ (et $a=2\neq 0$ : on peut utiliser le discriminant).

\textbf{Étape 2 : calculer le discriminant.}
\[ \Delta=b^2-4ac=(-5)^2-4\times 2\times 3=25-24=1. \]

\textbf{Étape 3 : conclure selon le signe de $\Delta$.} Comme $\Delta=1>0$, l'équation admet deux solutions distinctes :
\[ x_1=\frac{-(-5)-\sqrt{1}}{2\times 2}=\frac{5-1}{4}=1 \qquad ; \qquad x_2=\frac{5+1}{4}=\frac{3}{2}. \]
\textbf{Vérification} : $2(1)^2-5(1)+3=0$ et $2\left(\tfrac32\right)^2-5\left(\tfrac32\right)+3=\tfrac92-\tfrac{15}{2}+3=0$. \checkmark

\textbf{Conclusion :} $S=\left\{1\ ;\ \dfrac{3}{2}\right\}$.
\end{exoresolu}

\begin{attention}[Pièges de rédaction fréquents]
\begin{itemize}
\item \textbf{Signe de $b$.} Dans $2x^2-5x+3$, on a $b=-5$, donc $-b=+5$ et $b^2=25$ (et non $-25$). Écris $b$ entre parenthèses : $(-5)^2$.
\item \textbf{Le dénominateur est $2a$, pas $2$.} Pour $a=2$, on divise par $4$.
\item \textbf{Ne jamais « simplifier » $\sqrt{\Delta}$ à l'intérieur de la fraction sans vérifier} : si $\Delta=32$, alors $\sqrt{\Delta}=4\sqrt{2}$ et l'on peut simplifier par $2$ avec $2a$, mais uniquement terme à terme.
\item Une équation du second degré a \emph{toujours} au plus deux solutions réelles : en trouver trois est le signe d'une erreur de calcul.
\end{itemize}
\end{attention}

\courssub{Somme et produit des racines}

\textbf{Intuition.} Revenons au problème du menuisier : $L+\ell=7$ et $L\ell=12$. On cherche deux nombres connaissant leur somme et leur produit. Il existe un lien direct et remarquable entre les racines d'un trinôme et ses coefficients : c'est un outil de vérification très rapide et un classique du Probatoire.

\begin{propriete}[Somme et produit des racines]
Si l'équation $ax^2+bx+c=0$ ($a\neq 0$) admet deux racines réelles $x_1$ et $x_2$ (éventuellement égales), alors :
\[ S=x_1+x_2=-\frac{b}{a} \qquad \text{et} \qquad P=x_1\,x_2=\frac{c}{a}. \]
\emph{Démonstration exigible.}
\end{propriete}

\begin{experience}[Démonstration de la somme et du produit]
Supposons $\Delta\geq 0$ et notons $x_1=\dfrac{-b-\sqrt{\Delta}}{2a}$ et $x_2=\dfrac{-b+\sqrt{\Delta}}{2a}$ (si $\Delta=0$, les deux racines sont confondues et le calcul reste valable).

\textbf{Somme :}
\[ x_1+x_2=\frac{-b-\sqrt{\Delta}}{2a}+\frac{-b+\sqrt{\Delta}}{2a}=\frac{-2b}{2a}=-\frac{b}{a}. \]

\textbf{Produit :} en utilisant l'identité $(u-v)(u+v)=u^2-v^2$ avec $u=-b$ et $v=\sqrt{\Delta}$ :
\[ x_1\,x_2=\frac{(-b)^2-(\sqrt{\Delta})^2}{(2a)^2}=\frac{b^2-\Delta}{4a^2}=\frac{b^2-(b^2-4ac)}{4a^2}=\frac{4ac}{4a^2}=\frac{c}{a}. \]
\end{experience}

\begin{aretenir}[Deux nombres connaissant leur somme et leur produit]
Deux réels ont pour somme $S$ et pour produit $P$ si et seulement s'ils sont solutions de l'équation
\[ x^2-Sx+P=0. \]
Cette équation a des solutions réelles si et seulement si $S^2-4P\geq 0$.
\end{aretenir}

\begin{exemplebox}[Retour au problème du menuisier]
Les dimensions $L$ et $\ell$ vérifient $S=L+\ell=7$ et $P=L\ell=12$. Ce sont donc les solutions de $x^2-7x+12=0$. On calcule $\Delta=49-48=1>0$, d'où :
\[ x_1=\frac{7-1}{2}=3 \qquad ; \qquad x_2=\frac{7+1}{2}=4. \]
La plaque de Monsieur Ngo Bassa mesure donc $4$ m sur $3$ m. Vérification : périmètre $2(4+3)=14$ m et aire $4\times 3=12$ m$^2$. \checkmark
\end{exemplebox}

\courssub{Complément : le discriminant réduit}

\begin{definition}[Discriminant réduit]
Lorsque le coefficient $b$ est un nombre pair (ou plus généralement pour simplifier les calculs), on écrit $b=2b'$ et l'on définit le \cle{discriminant réduit} :
\[ \Delta'=(b')^2-ac. \]
On a alors $\Delta=4\Delta'$, donc $\Delta$ et $\Delta'$ ont le même signe, et les formules deviennent, lorsque $\Delta'\geq 0$ :
\[ x_1=\frac{-b'-\sqrt{\Delta'}}{a} \qquad \text{et} \qquad x_2=\frac{-b'+\sqrt{\Delta'}}{a}. \]
\end{definition}

\begin{exemplebox}[Le discriminant réduit en action]
Résolvons $3x^2-4x+1=0$. Ici $b=-4=2\times(-2)$, donc $b'=-2$ et :
\[ \Delta'=(-2)^2-3\times 1=4-3=1>0, \qquad x_1=\frac{2-1}{3}=\frac{1}{3},\quad x_2=\frac{2+1}{3}=1. \]
Les calculs sont plus légers qu'avec $\Delta=16$. Ce complément est utile mais non exigible : les formules avec $\Delta$ suffisent toujours.
\end{exemplebox}

\courssub{Équations du second degré à paramètre}

\textbf{Intuition.} Dans de nombreux problèmes (physique, économie), les coefficients dépendent d'une quantité inconnue $m$ appelée \cle{paramètre}. Le nombre de solutions dépend alors de la valeur de $m$ : on dit qu'on \emph{discute} l'équation suivant les valeurs du paramètre. La démarche est toujours la même : exprimer $\Delta$ en fonction de $m$, puis étudier son signe.

\begin{methode}[Discuter une équation à paramètre]
Pour discuter, selon les valeurs du paramètre $m$, le nombre de solutions de $ax^2+bx+c=0$ où un coefficient dépend de $m$ :
\begin{enumerate}
\item \textbf{Vérifier la condition $a\neq 0$.} Si $a$ dépend de $m$, traiter à part les valeurs qui annulent $a$ (l'équation devient alors du premier degré, voire impossible/indéterminée).
\item \textbf{Calculer le discriminant} $\Delta(m)$ : c'est une expression en $m$ (souvent un trinôme du second degré en $m$).
\item \textbf{Étudier le signe de $\Delta(m)$} : résoudre $\Delta(m)>0$, $\Delta(m)=0$, $\Delta(m)<0$ (souvent à l'aide d'un tableau de signes).
\item \textbf{Conclure} en rédigeant clairement les trois cas, avec les solutions explicites quand elles existent.
\end{enumerate}
\end{methode}

\begin{exoresolu}[Discussion selon $m$ de $x^2-2x+m=0$]
Discuter, selon les valeurs du réel $m$, le nombre de solutions de l'équation $x^2-2x+m=0$, puis donner les solutions lorsqu'elles existent.
\tcblower
\textbf{Étape 1.} Le coefficient de $x^2$ vaut $1\neq 0$ : c'est toujours une équation du second degré, quel que soit $m$.

\textbf{Étape 2.} Avec $a=1$, $b=-2$, $c=m$ :
\[ \Delta(m)=(-2)^2-4\times 1\times m=4-4m=4(1-m). \]

\textbf{Étape 3.} Signe de $\Delta(m)$ : $\Delta(m)>0 \iff m<1$ ; $\Delta(m)=0 \iff m=1$ ; $\Delta(m)<0 \iff m>1$.

\textbf{Étape 4.} Conclusion :
\begin{itemize}
\item Si $m<1$ : deux solutions distinctes $x_1=1-\sqrt{1-m}$ et $x_2=1+\sqrt{1-m}$.
\item Si $m=1$ : une solution double $x_0=1$ (l'équation est $(x-1)^2=0$).
\item Si $m>1$ : aucune solution réelle, $S=\varnothing$.
\end{itemize}
\end{exoresolu}

\begin{savaistu}[Al-Khwarizmi, père de la complétion du carré]
C'est le mathématicien perse Al-Khwarizmi (vers 820, Bagdad) qui a systématisé la résolution des équations du second degré par complétion du carré, avec des justifications \emph{géométriques} : il complétait littéralement des figures pour former des carrés. Son nom a donné le mot « algorithme », et le titre de son ouvrage, \emph{Al-Jabr}, a donné le mot « algèbre ». La formule avec le discriminant que tu utilises aujourd'hui est l'héritière directe de cette méthode millénaire.
\end{savaistu}

\begin{exercice}[Pour s'entraîner]
\begin{enumerate}
\item Résoudre dans $\mathbb{R}$ : a) $x^2-6x+9=0$ ; b) $3x^2+x-4=0$ ; c) $x^2+x+1=0$.
\item Mettre $-x^2+8x-7$ (le bénéfice de la coopérative de Bafoussam) sous forme canonique, puis résoudre $-x^2+8x-7=0$.
\item Trouver deux réels dont la somme est $11$ et le produit $30$.
\item Discuter selon $m$ le nombre de solutions de $x^2+mx+1=0$.
\end{enumerate}
\end{exercice}

\begin{corrige}[Exercice 1]
\textbf{1. a)} $\Delta=36-36=0$ : solution double $x_0=\dfrac{6}{2}=3$. $S=\{3\}$.

\textbf{b)} $\Delta=1+48=49>0$ : $x_1=\dfrac{-1-7}{6}=-\dfrac{4}{3}$, $x_2=\dfrac{-1+7}{6}=1$. $S=\left\{-\dfrac{4}{3}\ ;\ 1\right\}$.

\textbf{c)} $\Delta=1-4=-3<0$ : $S=\varnothing$.

\textbf{2.} $-x^2+8x-7=-\left(x^2-8x\right)-7=-\left[(x-4)^2-16\right]-7=-(x-4)^2+9$. L'équation s'écrit $(x-4)^2=9$, d'où $x-4=\pm 3$ : $S=\{1\ ;\ 7\}$. (On retrouve $\Delta=64-28=36$.)

\textbf{3.} Ce sont les solutions de $x^2-11x+30=0$ : $\Delta=121-120=1$, d'où $x_1=5$ et $x_2=6$.

\textbf{4.} $\Delta(m)=m^2-4$. Si $m<-2$ ou $m>2$ : deux solutions $\dfrac{-m\pm\sqrt{m^2-4}}{2}$. Si $m=-2$ : solution double $x_0=1$ ; si $m=2$ : solution double $x_0=-1$. Si $-2<m<2$ : aucune solution réelle.
\end{corrige}

\begin{aretenir}[L'essentiel de la section]
\begin{itemize}
\item Toute équation $ax^2+bx+c=0$ ($a\neq 0$) se résout grâce au \cle{discriminant} $\Delta=b^2-4ac$.
\item $\Delta>0$ : deux solutions $x_{1,2}=\dfrac{-b\pm\sqrt{\Delta}}{2a}$ ; $\Delta=0$ : solution double $x_0=-\dfrac{b}{2a}$ ; $\Delta<0$ : aucune solution réelle.
\item La \cle{forme canonique} $a\left(x+\dfrac{b}{2a}\right)^2-\dfrac{\Delta}{4a}$ est la clé de la démonstration et donne le sommet de la parabole.
\item Si $x_1$ et $x_2$ sont les racines : $x_1+x_2=-\dfrac{b}{a}$ et $x_1x_2=\dfrac{c}{a}$ ; réciproquement, deux nombres de somme $S$ et de produit $P$ sont solutions de $x^2-Sx+P=0$.
\item Pour une équation à paramètre : calculer $\Delta(m)$, étudier son signe, conclure cas par cas — sans oublier la condition $a\neq 0$.
\end{itemize}
\end{aretenir}

\coursec{Signe du trinôme et inéquations du second degré}

Dans la section précédente, nous avons appris à \emph{résoudre} l'équation $ax^2+bx+c=0$ : grâce à la forme canonique puis au discriminant $\Delta=b^2-4ac$, nous savons désormais trouver les racines du trinôme lorsqu'elles existent. Mais dans la vie courante, on ne demande pas toujours \og pour quelle valeur obtient-on exactement zéro ? \fg{}. On demande plutôt : \og à partir de quand le bénéfice est-il positif ? \fg{}, \og pour quelles valeurs l'aire dépasse-t-elle \SI{50}{m^2} ? \fg{}. Autrement dit, on cherche le \cle{signe} du trinôme, pas seulement ses zéros. C'est toute l'ambition de cette section : déterminer, pour chaque réel $x$, si $ax^2+bx+c$ est positif, négatif ou nul, puis résoudre les \cle{inéquations} du second degré.

\courssub{Factorisation du trinôme à l'aide de ses racines}

\begin{propriete}[Factorisation du trinôme]
Soit $f(x)=ax^2+bx+c$ un trinôme ($a\neq 0$) de discriminant $\Delta=b^2-4ac$.
\begin{itemize}
\item Si $\Delta>0$, le trinôme admet deux racines réelles $x_1$ et $x_2$, et pour tout réel $x$ :
\[ f(x)=a(x-x_1)(x-x_2). \]
\item Si $\Delta=0$, le trinôme admet une racine double $x_0=-\dfrac{b}{2a}$, et pour tout réel $x$ :
\[ f(x)=a(x-x_0)^2. \]
\item Si $\Delta<0$, le trinôme \textbf{ne se factorise pas} en produit de facteurs du premier degré dans $\mathbb{R}$.
\end{itemize}
\end{propriete}

\begin{experience}[Pourquoi cette factorisation est-elle vraie ?]
Démontrons le cas $\Delta>0$, qui est le cœur de tout ce qui suit. Partons de la forme canonique établie dans la section précédente :
\[ f(x)=a\left[\left(x+\frac{b}{2a}\right)^2-\frac{\Delta}{4a^2}\right]. \]
\textbf{Étape 1.} Comme $\Delta>0$, on peut écrire $\Delta=(\sqrt{\Delta})^2$. L'expression entre crochets devient une différence de deux carrés :
\[ \left(x+\frac{b}{2a}\right)^2-\left(\frac{\sqrt{\Delta}}{2a}\right)^2. \]
\textbf{Étape 2.} Appliquons l'identité remarquable $A^2-B^2=(A-B)(A+B)$ :
\[ f(x)=a\left(x+\frac{b}{2a}-\frac{\sqrt{\Delta}}{2a}\right)\left(x+\frac{b}{2a}+\frac{\sqrt{\Delta}}{2a}\right). \]
\textbf{Étape 3.} On reconnaît dans chaque parenthèse $x-x_1$ et $x-x_2$, puisque $x_1=\dfrac{-b-\sqrt{\Delta}}{2a}$ et $x_2=\dfrac{-b+\sqrt{\Delta}}{2a}$. D'où :
\[ f(x)=a(x-x_1)(x-x_2). \qquad \blacksquare \]
Le cas $\Delta=0$ est immédiat : la forme canonique se réduit à $f(x)=a\left(x+\dfrac{b}{2a}\right)^2=a(x-x_0)^2$.
\end{experience}

\begin{exemplebox}[Factoriser un trinôme]
Factorisons $f(x)=2x^2-10x+12$. On calcule $\Delta=(-10)^2-4\times 2\times 12=100-96=4>0$. Les racines sont :
\[ x_1=\frac{10-2}{4}=2 \qquad \text{et} \qquad x_2=\frac{10+2}{4}=3. \]
Donc $f(x)=2(x-2)(x-3)$. Vérifions en développant : $2(x-2)(x-3)=2(x^2-5x+6)=2x^2-10x+12$. \checkmark
\end{exemplebox}

\courssub{Théorème du signe du trinôme}

Voici l'intuition. Une parabole est une courbe \emph{continue} : pour passer d'une valeur positive à une valeur négative, elle est obligée de \emph{traverser} l'axe des abscisses, c'est-à-dire de s'annuler. Le signe du trinôme ne peut donc changer qu'en ses racines. Il suffit alors de savoir ce qui se passe \og loin des racines \fg{} : quand $x$ devient très grand, $ax^2$ domine, donc le trinôme a le signe de $a$. Tout le théorème découle de ces deux idées.

\begin{propriete}[Signe du trinôme — théorème fondamental]
Soit $f(x)=ax^2+bx+c$ avec $a\neq 0$.
\begin{itemize}
\item \textbf{Si $\Delta>0$} (racines $x_1<x_2$) : $f(x)$ est du \cle{signe de $a$} à l'extérieur de l'intervalle des racines, et du \cle{signe de $-a$} entre les racines. Il s'annule en $x_1$ et $x_2$.
\item \textbf{Si $\Delta=0$} (racine double $x_0$) : $f(x)$ est du signe de $a$ pour tout $x\neq x_0$, et nul en $x_0$.
\item \textbf{Si $\Delta<0$} : $f(x)$ est du signe de $a$ pour \emph{tout} réel $x$, et ne s'annule jamais.
\end{itemize}
\end{propriete}

\begin{experience}[Démonstration guidée du cas $\Delta>0$]
Supposons $\Delta>0$ et $a>0$. Grâce à la factorisation, $f(x)=a(x-x_1)(x-x_2)$ avec $x_1<x_2$. Le signe de $f(x)$ est celui du produit $(x-x_1)(x-x_2)$, puisque $a>0$. Construisons le tableau de signes de ce produit. Le facteur $x-x_1$ s'annule en $x_1$ et est positif après ; le facteur $x-x_2$ s'annule en $x_2$ et est positif après.

\begin{center}
\begin{tabular}{|c|ccccccc|}\hline
$x$ & $-\infty$ & & $x_1$ & & $x_2$ & & $+\infty$ \\ \hline
$x-x_1$ & & $-$ & $0$ & $+$ & & $+$ & \\ \hline
$x-x_2$ & & $-$ & & $-$ & $0$ & $+$ & \\ \hline
$f(x)$ & & $+$ & $0$ & $-$ & $0$ & $+$ & \\ \hline
\end{tabular}
\end{center}

Lisons la dernière ligne : sur $]-\infty\,;\,x_1[$, le produit de deux négatifs est \textbf{positif} (signe de $a$) ; sur $]x_1\,;\,x_2[$, un négatif fois un positif donne un \textbf{négatif} (signe de $-a$) ; sur $]x_2\,;\,+\infty[$, le produit de deux positifs est \textbf{positif} (signe de $a$). Si $a<0$, tous les signes de la dernière ligne sont inversés, ce qui redonne exactement l'énoncé du théorème. $\blacksquare$
\end{experience}

\begin{experience}[Démonstration du cas $\Delta<0$ via la forme canonique]
Supposons $\Delta<0$. La forme canonique s'écrit $f(x)=a\left[\left(x+\dfrac{b}{2a}\right)^2-\dfrac{\Delta}{4a^2}\right]$. Or $-\dfrac{\Delta}{4a^2}>0$ puisque $\Delta<0$. La quantité entre crochets est donc la somme d'un carré (positif ou nul) et d'un nombre \emph{strictement} positif : elle est \textbf{strictement positive} pour tout réel $x$. Par conséquent, $f(x)$ a le signe de $a$, partout, sans jamais s'annuler. $\blacksquare$

C'est un cas particulier de la propriété admise au programme : \emph{si un polynôme ne s'annule pas sur un intervalle, alors il garde un signe constant sur cet intervalle}.
\end{experience}

\begin{aretenir}[Résumé visuel à mémoriser]
\begin{center}
\begin{tabularx}{\linewidth}{|l|X|X|}\hline
\rowcolor{popblueL} \textbf{Discriminant} & \textbf{Signe du trinôme} & \textbf{Position de la parabole} \\ \hline
$\Delta>0$ & Signe de $a$ à l'extérieur de $[x_1\,;\,x_2]$, signe de $-a$ entre les racines & La parabole \emph{coupe} l'axe des abscisses en deux points \\ \hline
$\Delta=0$ & Signe de $a$ partout, nul en $x_0$ & La parabole est \emph{tangente} à l'axe en $x_0$ \\ \hline
$\Delta<0$ & Signe de $a$ pour tout $x\in\mathbb{R}$ & La parabole ne \emph{touche jamais} l'axe \\ \hline
\end{tabularx}
\end{center}
Phrase-clé : \og \textbf{le trinôme est du signe de $a$ à l'extérieur des racines} \fg{} (quand elles existent).
\end{aretenir}

\courssub{Construction du tableau de signes d'un trinôme}

\begin{exemplebox}[Tableau de signes complet]
Étudions le signe de $f(x)=-2x^2+10x-12=-2(x-2)(x-3)$ (d'après l'exemple précédent, avec $a=-2<0$, racines $2$ et $3$).

\begin{center}
\begin{tabular}{|c|ccccccc|}\hline
$x$ & $-\infty$ & & $2$ & & $3$ & & $+\infty$ \\ \hline
$f(x)$ & & $-$ & $0$ & $+$ & $0$ & $-$ & \\ \hline
\end{tabular}
\end{center}

Comme $a<0$, le trinôme est négatif à l'extérieur des racines et positif entre elles : c'est bien le signe de $-a$ sur $]2\,;\,3[$.
\end{exemplebox}

\begin{popfigure}
\begin{center}
\begin{tikzpicture}
\begin{axis}[
  width=0.85\linewidth, height=6.5cm,
  axis lines=middle, xlabel=$x$, ylabel=$y$,
  xmin=-1, xmax=5, ymin=-8, ymax=4,
  xtick={1,2,3,4}, ytick=\empty,
  legend style={at={(0.02,0.98)}, anchor=north west, font=\small}
]
\addplot[popcurve, domain=-0.5:4.5, samples=200] {-2*x^2+10*x-12};
\addlegendentry{$f(x)=-2x^2+10x-12$}
\addplot[mark=*, only marks, mark size=2pt, popdark] coordinates {(2,0) (3,0)};
\node[popgreen!60!black, font=\small\bfseries] at (axis cs:2.5,1.2) {$f(x)>0$};
\node[poporange!80!black, font=\small\bfseries] at (axis cs:0.6,-5.5) {$f(x)<0$};
\node[poporange!80!black, font=\small\bfseries] at (axis cs:4.3,-5.5) {$f(x)<0$};
\end{axis}
\end{tikzpicture}
\end{center}
\legende{Parabole de $f(x)=-2x^2+10x-12$ : elle coupe l'axe des abscisses en $x=2$ et $x=3$. Zone verte (entre les racines) : le trinôme est positif (signe de $-a$ car $a<0$). Zones orange (extérieur) : le trinôme est négatif.}
\end{popfigure}

\courssub{Résolution des inéquations du second degré}

Résoudre $ax^2+bx+c\geq 0$ (ou $>0$, $\leq 0$, $<0$) consiste simplement à \emph{lire} le tableau de signes et à écrire la réponse sous forme d'\textbf{intervalles} ou de réunion d'intervalles. La seule subtilité : les bornes sont \emph{fermées} si l'inégalité est large ($\geq$, $\leq$), \emph{ouvertes} si elle est stricte ($>$, $<$), car en une racine le trinôme vaut exactement $0$.

\begin{methode}[Résoudre une inéquation du second degré en 4 étapes]
\begin{enumerate}
\item \textbf{Calculer le discriminant} $\Delta=b^2-4ac$ du trinôme $ax^2+bx+c$.
\item \textbf{Déterminer les racines} éventuelles et les ranger dans l'ordre croissant : $x_1<x_2$.
\item \textbf{Dresser le tableau de signes} : signe de $a$ à l'extérieur des racines, signe de $-a$ entre elles (si $\Delta\leq 0$, le trinôme garde le signe de $a$ partout).
\item \textbf{Conclure en intervalles} : sélectionner les zones correspondant à l'inégalité demandée, avec des crochets ouverts ou fermés selon que l'inégalité est stricte ou large.
\end{enumerate}
\end{methode}

\begin{exoresolu}[Le bénéfice de la coopérative]
Une coopérative de producteurs de cacao de la région du Centre transforme et vend du beurre de cacao. Si $x$ désigne le nombre de tonnes produites et vendues par mois, le bénéfice mensuel (en centaines de milliers de FCFA) est modélisé par :
\[ B(x)=-x^2+8x-7. \]
Pour quelles quantités produites la coopérative réalise-t-elle un bénéfice \emph{strictement positif} ?
\tcblower
\textbf{Solution détaillée.} On cherche les valeurs de $x$ telles que $B(x)>0$, c'est-à-dire $-x^2+8x-7>0$.

\textbf{Étape 1 : discriminant.} Avec $a=-1$, $b=8$, $c=-7$ :
\[ \Delta=8^2-4\times(-1)\times(-7)=64-28=36>0. \]

\textbf{Étape 2 : racines.} $\sqrt{\Delta}=6$. Les deux racines sont :
\[ \frac{-8+6}{-2}=1 \qquad \text{et} \qquad \frac{-8-6}{-2}=7. \]
On les range dans l'ordre croissant : $x_1=1$ et $x_2=7$.

\textbf{Étape 3 : tableau de signes.} Comme $a=-1<0$, le trinôme est négatif à l'extérieur des racines et positif entre elles :
\begin{center}
\begin{tabular}{|c|ccccccc|}\hline
$x$ & $-\infty$ & & $1$ & & $7$ & & $+\infty$ \\ \hline
$B(x)$ & & $-$ & $0$ & $+$ & $0$ & $-$ & \\ \hline
\end{tabular}
\end{center}

\textbf{Étape 4 : conclusion.} L'inégalité est \emph{stricte}, donc les bornes sont exclues :
\[ \mathcal{S}=\,]1\,;\,7[. \]

\textbf{Interprétation.} La coopérative est bénéficiaire lorsqu'elle produit et vend \emph{entre 1 tonne et 7 tonnes} par mois (bornes exclues : à 1 tonne ou 7 tonnes, le bénéfice est nul, c'est le seuil de rentabilité). En dessous de 1 tonne, les coûts fixes écrasent la recette ; au-delà de 7 tonnes, les surcoûts de transformation dépassent les ventes.
\end{exoresolu}

\begin{exoresolu}[Inéquations avec discriminant nul ou négatif]
Résoudre dans $\mathbb{R}$ : \quad (a) $x^2-6x+9\geq 0$ \qquad (b) $2x^2-3x+4<0$.
\tcblower
\textbf{(a)} $\Delta=36-36=0$ : racine double $x_0=\dfrac{6}{2}=3$. Le trinôme $x^2-6x+9=(x-3)^2$ est du signe de $a=1>0$ partout, et nul en $3$. Comme l'inégalité est large, $3$ est acceptée :
\[ \mathcal{S}=\mathbb{R}. \]
\textbf{(b)} $\Delta=9-32=-23<0$ : le trinôme garde le signe de $a=2>0$ sur tout $\mathbb{R}$. Il n'est donc \emph{jamais} strictement négatif :
\[ \mathcal{S}=\varnothing. \]
\end{exoresolu}

\begin{attention}[Ne divise jamais une inéquation par une expression dont tu ignores le signe]
Face à $(x-1)(x-3)>0$, une tentation fréquente est de \og simplifier \fg{} en divisant les deux membres par $(x-1)$ pour obtenir $x-3>0$. \textbf{C'est faux et dangereux !} Diviser une inéquation par un nombre \emph{négatif} en change le sens : or le signe de $(x-1)$ dépend de $x$. En divisant sans discuter, on perd la moitié des solutions (ici $\mathcal{S}=\,]-\infty\,;\,1[\,\cup\,]3\,;\,+\infty[$, et non $]3\,;\,+\infty[$).

La bonne attitude : \textbf{tout ramener à zéro} (transposer, factoriser), puis dresser un \textbf{tableau de signes}. C'est la seule méthode qui garantit le respect du sens des inégalités.
\end{attention}

\begin{attention}[Autres pièges classiques]
\begin{itemize}
\item \textbf{Racines mal rangées} : écris toujours $x_1<x_2$ avant de tracer le tableau, surtout quand $a<0$ (la formule donne alors les racines \og dans le désordre \fg{}).
\item \textbf{Crochets} : pour $f(x)>0$ (strict), les racines sont \emph{exclues} ; pour $f(x)\geq 0$ (large), elles sont \emph{incluses}.
\item \textbf{Signe de $a$} : n'oublie pas que si $a<0$, c'est le signe de $-a$ qui règne \emph{entre} les racines. Relis ton tableau en te demandant : \og loin des racines, mon trinôme a-t-il bien le signe de $a$ ? \fg{}
\end{itemize}
\end{attention}

\courssub{Lien avec la parabole : une lecture graphique du signe}

Le tableau de signes n'est qu'une traduction algébrique d'une réalité géométrique très simple. La courbe de $f(x)=ax^2+bx+c$ est une parabole :
\begin{itemize}
\item si $\Delta>0$, elle \textbf{coupe} l'axe des abscisses aux points d'abscisses $x_1$ et $x_2$ ; les solutions de $f(x)>0$ sont les abscisses des points de la parabole situés \emph{au-dessus} de l'axe ;
\item si $\Delta=0$, elle est \textbf{tangente} à l'axe en $x_0$ et reste d'un seul côté ;
\item si $\Delta<0$, elle \textbf{flotte} entièrement au-dessus (si $a>0$) ou au-dessous (si $a<0$) de l'axe.
\end{itemize}
Résoudre une inéquation du second degré, c'est donc répondre à la question : \emph{où la parabole est-elle au-dessus (ou au-dessous) de l'axe des abscisses ?}

\begin{exercice}[À toi de jouer]
\begin{enumerate}
\item Résoudre dans $\mathbb{R}$ : $3x^2-5x-2\leq 0$ ; \quad $-x^2+4x-5>0$ ; \quad $x^2-10x+25<0$.
\item Un planteur de la plaine des Mbo clôture un champ rectangulaire avec \SI{100}{m} de grillage, contre une rivière (un seul côté sans grillage). Montrer que l'aire du champ est $A(x)=-2x^2+100x$ où $x$ est la largeur, puis déterminer pour quelles valeurs de $x$ cette aire dépasse \SI{1152}{m^2}.
\end{enumerate}
\end{exercice}

\begin{corrige}[Exercice 1]
\textbf{1.} Pour $3x^2-5x-2\leq 0$ : $\Delta=25+24=49$, racines $x_1=-\dfrac{1}{3}$ et $x_2=2$. Comme $a=3>0$, le trinôme est négatif \emph{entre} les racines ; inégalité large, donc $\mathcal{S}=\left[-\dfrac{1}{3}\,;\,2\right]$.

Pour $-x^2+4x-5>0$ : $\Delta=16-20=-4<0$, le trinôme garde le signe de $a=-1<0$ partout : $\mathcal{S}=\varnothing$.

Pour $x^2-10x+25<0$ : $\Delta=0$, $(x-5)^2$ est positif ou nul, jamais strictement négatif : $\mathcal{S}=\varnothing$.

\textbf{2.} Si $x$ est la largeur (côtés perpendiculaires à la rivière), le grillage couvre $2x$ de largeurs et une longueur $L=100-2x$. L'aire vaut $A(x)=x(100-2x)=-2x^2+100x$. On résout $-2x^2+100x>1152$, soit $-2x^2+100x-1152>0$, ou encore $x^2-50x+576<0$ en divisant par $-2$ (sens changé). $\Delta=2500-2304=196$, $\sqrt{\Delta}=14$, racines $x_1=18$ et $x_2=32$. Le trinôme est négatif entre les racines : $\mathcal{S}=\,]18\,;\,32[$. L'aire dépasse \SI{1152}{m^2} lorsque la largeur est comprise entre \SI{18}{m} et \SI{32}{m}.
\end{corrige}

\begin{savaistu}[Il y a 4000 ans, les Babyloniens complétaient déjà le carré]
Sur des tablettes d'argile babyloniennes gravées vers 1800 av. J.-C. (comme la tablette BM 13901), les scribes résolvaient des problèmes du type : \og j'ai ajouté l'aire d'un carré à son côté et j'ai obtenu $\frac{3}{4}$ ; combien vaut le côté ? \fg{} — autrement dit $x^2+x=\frac{3}{4}$. Leur méthode ? Ajouter le carré de la moitié du coefficient de $x$ des deux côtés : exactement notre \emph{forme canonique}, notre \og complétion du carré \fg{} ! Ils ne connaissaient ni les nombres négatifs ni le symbolisme algébrique, pourtant leur procédure, transmise aux Grecs puis aux mathématiciens arabes comme Al-Khwarizmi (IX\textsuperscript{e} siècle), est celle que tu utilises aujourd'hui au lycée. La prochaine fois que tu écriras $\left(x+\frac{b}{2a}\right)^2$, souviens-toi : tu appliques une idée vieille de quarante siècles.
\end{savaistu}

Tu maîtrises désormais le signe du trinôme et la résolution des inéquations du second degré, outil indispensable des problèmes de seuil. Dans la section suivante, nous verrons comment des \textbf{systèmes d'équations à deux inconnues} — somme et produit, par exemple — se ramènent précisément à une équation du second degré.

\coursec{Systèmes se ramenant à une équation du second degré}

Dans les sections précédentes, tu as appris à résoudre une équation du second degré grâce à la forme canonique et au discriminant, puis à étudier le signe d'un trinôme $ax^2+bx+c$ pour résoudre les inéquations du second degré. Nous allons maintenant utiliser ces acquis pour résoudre des problèmes à \emph{deux} inconnues. L'idée directrice est simple et puissante : beaucoup de systèmes à deux inconnues, qui semblent compliqués parce qu'ils contiennent des termes comme $xy$ ou $x^2+y^2$, peuvent être \emph{ramenés} à une seule équation du second degré. C'est exactement le savoir-faire exigé au Probatoire : transformer, résoudre, puis vérifier.

\courssub{Le problème somme-produit : une situation très concrète}

Commençons par une situation que tout le monde a déjà rencontrée. Un commerçant du marché Mokolo te dit : « J'ai deux nombres en tête. Leur somme vaut $7$ et leur produit vaut $12$. Peux-tu les retrouver ? » Après quelques essais, tu réponds : $3$ et $4$, car $3+4=7$ et $3\times 4 = 12$. Facile ! Mais que faire si la somme vaut $17$ et le produit $60$ ? Ou si la somme vaut $\sqrt{5}$ et le produit $-3$ ? Le tâtonnement ne suffit plus : il faut une méthode générale, et c'est le second degré qui va nous la fournir.

\begin{definition}[Système somme-produit]
On appelle \cle{système somme-produit} tout système de la forme
\[
\begin{cases} x + y = S \\ xy = P \end{cases}
\]
où $S$ et $P$ sont deux réels donnés et $(x\,;\,y)$ le couple d'inconnues réelles.
\end{definition}

L'observation clé est la suivante : si l'on connaît la somme et le produit de deux nombres, on peut fabriquer une équation du second degré dont ces deux nombres sont exactement les racines. C'est l'objet du théorème fondamental de cette section.

\begin{propriete}[Théorème somme-produit (démonstration exigible)]
Deux nombres réels $x$ et $y$ ont pour somme $S$ et pour produit $P$ si et seulement s'ils sont les deux racines (éventuellement confondues) de l'équation d'inconnue $X$ :
\[
X^2 - SX + P = 0.
\]
En particulier, le système $\begin{cases} x+y=S \\ xy=P \end{cases}$ admet des solutions réelles si et seulement si $S^2 - 4P \geq 0$.
\end{propriete}

\begin{experience}[Démonstration du théorème somme-produit]
Procédons en deux sens, comme il se doit pour une équivalence.

\textbf{Sens direct.} Supposons que $x$ et $y$ vérifient $x+y=S$ et $xy=P$. Alors $S = x+y$ et $P = xy$. Évaluons le trinôme $X^2 - SX + P$ en $X = x$ :
\[
x^2 - Sx + P = x^2 - (x+y)x + xy = x^2 - x^2 - xy + xy = 0.
\]
Donc $x$ est racine de $X^2 - SX + P = 0$. Le même calcul avec $y$ (les rôles de $x$ et $y$ sont symétriques) montre que $y$ est aussi racine.

\textbf{Sens réciproque.} Supposons que $x$ et $y$ soient les deux racines de $X^2 - SX + P = 0$. D'après les relations entre coefficients et racines d'un trinôme (vues en section 1), la somme des racines vaut $-\dfrac{-S}{1} = S$ et leur produit vaut $\dfrac{P}{1} = P$. Donc $x+y=S$ et $xy=P$.

\textbf{Condition d'existence.} L'équation $X^2 - SX + P = 0$ a pour discriminant $\Delta = (-S)^2 - 4P = S^2 - 4P$. Elle admet des racines réelles si et seulement si $\Delta \geq 0$, c'est-à-dire $S^2 - 4P \geq 0$. Le théorème est démontré.
\end{experience}

\begin{attention}[Pas de solutions quand $S^2 - 4P < 0$]
L'erreur la plus fréquente consiste à « inventer » des solutions lorsque $S^2 - 4P < 0$. Par exemple, chercher deux réels de somme $2$ et de produit $5$ est impossible : $S^2 - 4P = 4 - 20 = -16 < 0$. Aucun couple de nombres réels ne convient, et il faut l'écrire clairement : $\mathcal{S} = \varnothing$. Autre piège : le théorème donne les deux nombres \emph{sans ordre} ; si $x$ et $y$ jouent des rôles différents dans le problème (longueur et largeur, par exemple), il faut attribuer correctement chaque valeur, ou donner les deux couples.
\end{attention}

\begin{methode}[Résoudre un système somme-produit]
Pour résoudre $\begin{cases} x+y=S \\ xy=P \end{cases}$ :
\begin{enumerate}
\item Calculer le discriminant auxiliaire $\Delta = S^2 - 4P$.
\item Si $\Delta < 0$ : conclure qu'il n'y a \textbf{aucune solution réelle}.
\item Si $\Delta \geq 0$ : résoudre $X^2 - SX + P = 0$ ; on obtient deux racines $X_1$ et $X_2$ (éventuellement confondues).
\item Les solutions sont les couples $(X_1\,;\,X_2)$ et $(X_2\,;\,X_1)$ (un seul couple si $X_1 = X_2$).
\item \textbf{Vérifier} chaque couple dans les deux équations de départ : somme et produit.
\end{enumerate}
\end{methode}

\begin{exemplebox}[Deux nombres de somme 7 et de produit 12]
On cherche $x$ et $y$ tels que $x+y=7$ et $xy=12$. Ici $S=7$, $P=12$, donc $\Delta = 49 - 48 = 1 > 0$. L'équation auxiliaire est
\[
X^2 - 7X + 12 = 0, \qquad X = \frac{7 \pm 1}{2}, \quad \text{soit } X_1 = 3 \text{ et } X_2 = 4.
\]
Vérification : $3+4=7$ et $3\times 4 = 12$. Les solutions sont les couples $(3\,;\,4)$ et $(4\,;\,3)$.
\end{exemplebox}

\courssub{Systèmes par substitution : une équation linéaire et une équation du second degré}

Dans de nombreux problèmes, le système ne se présente pas directement sous la forme somme-produit. Un cas très classique est celui où \textbf{l'une des équations est linéaire} (par exemple $x + y = 5$ ou $2x - y = 1$) et \textbf{l'autre contient des termes du second degré} (par exemple $xy = 6$, $x^2 + y^2 = 13$). La stratégie est alors la substitution.

\begin{methode}[Résolution par substitution]
\begin{enumerate}
\item Exprimer une inconnue en fonction de l'autre grâce à l'équation linéaire (par exemple $y = 5 - x$).
\item Substituer cette expression dans l'équation du second degré : on obtient une équation du second degré en une seule inconnue.
\item Résoudre cette équation (discriminant).
\item Pour chaque valeur trouvée, calculer la valeur de l'autre inconnue grâce à la relation linéaire.
\item Vérifier chaque couple obtenu dans \emph{les deux} équations du système.
\end{enumerate}
\end{methode}

\begin{exoresolu}[Deux systèmes résolus pas à pas]
Résoudre dans $\mathbb{R}^2$ le système $\begin{cases} x + y = 5 \\ xy = 6 \end{cases}$ puis le système $\begin{cases} 2x - y = 1 \\ x^2 + xy = 6 \end{cases}$.
\tcblower
\textbf{Premier système.} C'est un système somme-produit avec $S=5$, $P=6$ : $\Delta = 25-24=1 > 0$, l'équation $X^2-5X+6=0$ donne $X = \dfrac{5 \pm 1}{2}$, soit $X \in \{2\,;\,3\}$. Vérification : $2+3=5$ et $2\times 3 = 6$. Donc $\mathcal{S} = \{(2\,;\,3)\,;\,(3\,;\,2)\}$.

\textbf{Deuxième système.} L'équation linéaire donne $y = 2x - 1$. Substituons dans la seconde équation :
\begin{align*}
x^2 + x(2x-1) &= 6 \\
x^2 + 2x^2 - x &= 6 \\
3x^2 - x - 6 &= 0
\end{align*}
Le discriminant est $\Delta = (-1)^2 - 4 \times 3 \times (-6) = 1 + 72 = 73 > 0$, d'où
\[
x_1 = \frac{1 + \sqrt{73}}{6} \quad \text{et} \quad x_2 = \frac{1 - \sqrt{73}}{6}.
\]
La relation $y = 2x - 1$ donne alors :
\[
y_1 = 2x_1 - 1 = \frac{1+\sqrt{73}}{3} - 1 = \frac{\sqrt{73}-2}{3}, \qquad
y_2 = 2x_2 - 1 = \frac{1-\sqrt{73}}{3} - 1 = \frac{-2-\sqrt{73}}{3}.
\]
Vérification pour le premier couple : $2x_1 - y_1 = \dfrac{1+\sqrt{73}}{3} - \dfrac{\sqrt{73}-2}{3} = \dfrac{3}{3} = 1$, et $x_1^2 + x_1 y_1 = x_1(x_1 + y_1) = x_1(3x_1 - 1) = 3x_1^2 - x_1 = 6$ car $x_1$ est racine de $3x^2 - x - 6 = 0$. Le même calcul vaut pour le second couple. Donc
\[
\mathcal{S} = \left\{ \left(\frac{1+\sqrt{73}}{6}\,;\,\frac{\sqrt{73}-2}{3}\right)\,;\,\left(\frac{1-\sqrt{73}}{6}\,;\,\frac{-2-\sqrt{73}}{3}\right) \right\}.
\]
\end{exoresolu}

\courssub{L'identité remarquable $x^2+y^2 = (x+y)^2 - 2xy$ : le pont vers $S$ et $P$}

Que faire quand le système contient $x^2 + y^2$ au lieu de $xy$ ? L'astuce consiste à tout exprimer en fonction de la somme $S = x+y$ et du produit $P = xy$, grâce à l'identité
\[
x^2 + y^2 = (x+y)^2 - 2xy = S^2 - 2P.
\]
Cette identité se démontre immédiatement en développant : $(x+y)^2 = x^2 + 2xy + y^2$, donc $x^2+y^2 = (x+y)^2 - 2xy$. On dit que $x^2+y^2$ est une \emph{expression symétrique} en $x$ et $y$ : elle s'exprime entièrement à l'aide de $S$ et $P$.

\begin{exemplebox}[Somme et somme des carrés]
Résolvons $\begin{cases} x + y = 5 \\ x^2 + y^2 = 13 \end{cases}$. Posons $S = x+y = 5$ et $P = xy$. Alors
\[
x^2 + y^2 = S^2 - 2P \iff 13 = 25 - 2P \iff 2P = 12 \iff P = 6.
\]
On est ramené au système somme-produit $S=5$, $P=6$, déjà résolu plus haut : les solutions sont $(2\,;\,3)$ et $(3\,;\,2)$. Vérification : $2^2+3^2 = 4+9=13$ et $2+3=5$.
\end{exemplebox}

\begin{attention}[Bien vérifier chaque couple]
Lorsqu'on passe par $S$ et $P$, on obtient les deux nombres \emph{à l'ordre près}. Il faut donc toujours tester les deux couples $(X_1\,;\,X_2)$ et $(X_2\,;\,X_1)$ dans le système \emph{initial} : si celui-ci n'est pas symétrique en $x$ et $y$, un seul des deux couples peut convenir. Oublier la vérification est une faute classique au Probatoire.
\end{attention}

\courssub{Interprétation géométrique : intersection d'une droite et d'une hyperbole}

Il est très éclairant de visualiser le système somme-produit. Dans un repère, l'équation $x + y = S$ représente une \cle{droite} de coefficient directeur $-1$, et l'équation $xy = P$ (avec $P \neq 0$) représente une \cle{hyperbole} d'équation $y = \dfrac{P}{x}$. Résoudre le système revient à chercher les points d'intersection de ces deux courbes : il y en a deux si $S^2 - 4P > 0$, un seul (point de tangence) si $S^2 - 4P = 0$, et aucun si $S^2 - 4P < 0$.

\begin{popfigure}
\begin{center}
\begin{tikzpicture}
\begin{axis}[
  width=0.85\linewidth, height=7cm,
  axis lines=middle,
  xlabel={$x$}, ylabel={$y$},
  xmin=-1, xmax=8, ymin=-1, ymax=8,
  xtick={1,2,3,4,5,6,7}, ytick={1,2,3,4,5,6,7},
  grid=both, grid style={popline!40},
  legend style={at={(0.98,0.95)}, anchor=north east, font=\small}
]
\addplot[popcurve, domain=1.6:7.5, samples=200] {12/x};
\addlegendentry{$xy=12$}
\addplot[popcurve, poporange, domain=-0.5:7.5, samples=2] {7-x};
\addlegendentry{$x+y=7$}
\addplot[only marks, mark=*, popdark, mark size=2pt] coordinates {(3,4) (4,3)};
\node[popdark, anchor=south west] at (axis cs:3,4) {$(3\,;\,4)$};
\node[popdark, anchor=north west] at (axis cs:4,3) {$(4\,;\,3)$};
\end{axis}
\end{tikzpicture}
\end{center}
\legende{Les solutions du système $x+y=7$, $xy=12$ sont les coordonnées des points d'intersection de la droite et de l'hyperbole.}
\end{popfigure}

\courssub{Application : le problème du menuisier}

Voici l'application géométrique emblématique de cette leçon, très fréquente aux examens : connaître le périmètre d'un rectangle, c'est connaître la \emph{somme} de ses dimensions ; connaître son aire, c'est connaître leur \emph{produit}.

\begin{exoresolu}[La plaque du menuisier]
Un menuisier de Douala doit découper une plaque rectangulaire de contreplaqué. Le client exige un périmètre de $14$ m et une aire de $12$ m\textsuperscript{2}. Déterminer les dimensions de la plaque.
\tcblower
Notons $x$ et $y$ les dimensions (en mètres) du rectangle, avec $x > 0$ et $y > 0$.
\begin{itemize}
\item Périmètre : $2(x+y) = 14$, donc $x + y = 7$.
\item Aire : $xy = 12$.
\end{itemize}
On reconnaît un système somme-produit avec $S = 7$ et $P = 12$. Vérifions l'existence : $S^2 - 4P = 49 - 48 = 1 \geq 0$. L'équation auxiliaire $X^2 - 7X + 12 = 0$ a pour racines
\[
X = \frac{7 \pm 1}{2}, \quad \text{soit } X_1 = 3 \text{ et } X_2 = 4.
\]
Vérification : $2(3+4) = 14$ m et $3 \times 4 = 12$ m\textsuperscript{2}. Les deux valeurs étant strictement positives, elles conviennent comme dimensions. La plaque doit donc mesurer $3$ m sur $4$ m.
\end{exoresolu}

\begin{popfigure}
\begin{center}
\begin{tikzpicture}[scale=0.85]
\draw[fill=popblueL, draw=popblue, thick] (0,0) rectangle (4,3);
\draw[<->, popdark, thick] (0,-0.5) -- (4,-0.5) node[midway, below] {$x$};
\draw[<->, popdark, thick] (-0.5,0) -- (-0.5,3) node[midway, left] {$y$};
\node[popdark] at (2,1.5) {Aire $= xy = 12$ m\textsuperscript{2}};
\node[popdark, anchor=south] at (2,3.15) {Périmètre $= 2(x+y) = 14$ m};
\draw[poporange, thick] (0,0) -- (4,0) -- (4,3) -- (0,3) -- cycle;
\end{tikzpicture}
\end{center}
\legende{Le rectangle du menuisier : connaître le périmètre donne la somme $S$, connaître l'aire donne le produit $P$.}
\end{popfigure}

\begin{savaistu}[Somme, produit et identités remarquables au quotidien]
Les relations entre somme et produit remontent aux mathématiques babyloniennes, il y a près de 4000 ans : les scribes résolvaient déjà des problèmes de partage de champs rectangulaires dont on connaissait le périmètre et l'aire, exactement comme notre menuisier ! Aujourd'hui, ces mêmes idées interviennent partout : en physique (association de résistances), en économie, et même dans les codes correcteurs d'erreurs utilisés par les réseaux de téléphonie mobile comme MTN et Orange, qui reposent sur des propriétés des racines de polynômes.
\end{savaistu}

\begin{aretenir}[L'essentiel de la section]
\begin{itemize}
\item Deux nombres de somme $S$ et de produit $P$ sont les racines de $X^2 - SX + P = 0$.
\item Le système somme-produit a des solutions réelles si et seulement si $S^2 - 4P \geq 0$.
\item Une équation linéaire permet la substitution ; l'identité $x^2+y^2 = S^2 - 2P$ ramène les sommes de carrés à un système somme-produit.
\item Toujours \textbf{vérifier} les couples solutions dans le système de départ.
\item Périmètre et aire d'un rectangle fournissent directement $S$ et $P$.
\end{itemize}
\end{aretenir}

\begin{exercice}[À toi de jouer]
\begin{enumerate}
\item Résoudre dans $\mathbb{R}^2$ : $\begin{cases} x + y = 9 \\ xy = 20 \end{cases}$ puis $\begin{cases} x + y = 3 \\ x^2 + y^2 = 29 \end{cases}$.
\item Un champ rectangulaire a un périmètre de $100$ m et une aire de $600$ m\textsuperscript{2}. Quelles sont ses dimensions ?
\item Existe-t-il deux réels dont la somme vaut $6$ et le produit $10$ ? Justifier.
\end{enumerate}
\end{exercice}

\begin{corrige}[Exercice : À toi de jouer]
\textbf{1.} Premier système : $S=9$, $P=20$, $\Delta = 81 - 80 = 1 > 0$ ; $X^2 - 9X + 20 = 0$ donne $X = \dfrac{9 \pm 1}{2} \in \{4\,;\,5\}$. Solutions : $(4\,;\,5)$ et $(5\,;\,4)$. Second système : $S = 3$ et $S^2 - 2P = 29$ donne $9 - 2P = 29$, soit $P = -10$. Alors $\Delta = S^2 - 4P = 9 + 40 = 49 > 0$ ; $X^2 - 3X - 10 = 0$ donne $X = \dfrac{3 \pm 7}{2} \in \{-2\,;\,5\}$. Solutions : $(-2\,;\,5)$ et $(5\,;\,-2)$. Vérification : $(-2)^2 + 5^2 = 4 + 25 = 29$ et $-2 + 5 = 3$.

\textbf{2.} $2(x+y) = 100$ donne $S = 50$ ; $P = 600$. $\Delta = 2500 - 2400 = 100 > 0$ ; $X = \dfrac{50 \pm 10}{2} \in \{20\,;\,30\}$. Le champ mesure $20$ m sur $30$ m. Vérification : $2(20+30) = 100$ m et $20 \times 30 = 600$ m\textsuperscript{2}.

\textbf{3.} $S^2 - 4P = 36 - 40 = -4 < 0$ : d'après le théorème somme-produit, aucune solution réelle n'existe.
\end{corrige}

\coursec{Polynômes de degré 3 : zéros et factorisation}

Dans la section précédente, nous avons résolu des systèmes à deux inconnues en nous \emph{ramenant} à une équation du second degré. Cette idée — transformer un problème difficile en un problème déjà maîtrisé — est l'une des plus puissantes de toutes les mathématiques. Nous allons maintenant l'appliquer aux équations de degré 3 : grâce à la \cle{factorisation}, résoudre une équation de degré 3 consistera à se ramener à une équation du second degré, que tu sais déjà traiter avec le discriminant.

\courssub{Polynômes de degré 3 : définition et vocabulaire}

\begin{definition}[Polynôme de degré 3]
On appelle \cle{polynôme de degré 3} (ou \cle{polynôme du troisième degré}) toute fonction $P$ définie sur $\mathbb{R}$ par
\[ P(x) = ax^3 + bx^2 + cx + d, \]
où $a$, $b$, $c$ et $d$ sont des nombres réels fixés, appelés les \cle{coefficients} du polynôme, avec $a \neq 0$. Le coefficient $a$ est appelé \cle{coefficient dominant}, et $d$ est le \cle{terme constant}.
\end{definition}

\begin{attention}[La condition $a \neq 0$ est essentielle]
C'est la présence effective du terme en $x^3$ qui fait le degré. Si $a = 0$, alors $P(x) = bx^2 + cx + d$ est un polynôme de degré au plus 2. Dire « $P$ est de degré 3 » sans vérifier que $a \neq 0$ est une erreur de rigueur fréquente, surtout dans les exercices avec paramètre.
\end{attention}

\begin{exemplebox}[Reconnaître les coefficients]
Pour $P(x) = 2x^3 - 5x^2 + x - 7$, on a $a = 2$, $b = -5$, $c = 1$, $d = -7$. Pour $Q(x) = x^3 + 4x$, on a $a = 1$, $b = 0$, $c = 4$, $d = 0$ : un coefficient nul est parfaitement autorisé, sauf pour $a$.
\end{exemplebox}

\courssub{Zéro (ou racine) d'un polynôme}

Voici l'intuition fondatrice de toute cette section. Graphiquement, une équation $P(x) = 0$ revient à chercher les abscisses des points où la courbe de $P$ \emph{coupe} l'axe des abscisses. Ces abscisses portent un nom.

\begin{definition}[Zéro d'un polynôme]
Soit $P$ un polynôme et $\alpha$ un nombre réel. On dit que $\alpha$ est un \cle{zéro} (ou une \cle{racine}) de $P$ lorsque
\[ P(\alpha) = 0. \]
\end{definition}

Vérifier qu'un nombre est zéro d'un polynôme est un simple calcul de substitution — mais c'est un calcul qui doit être rédigé proprement.

\begin{exemplebox}[Vérifier qu'un nombre est un zéro]
Soit $P(x) = x^3 - 4x^2 + x + 6$. Vérifions que $2$ est un zéro de $P$ :
\[ P(2) = 2^3 - 4 \times 2^2 + 2 + 6 = 8 - 16 + 2 + 6 = 0. \]
Donc $2$ est un zéro de $P$. En revanche, $P(1) = 1 - 4 + 1 + 6 = 4 \neq 0$, donc $1$ n'est pas un zéro de $P$.
\end{exemplebox}

\begin{methode}[Chercher un zéro évident]
Pour trouver un premier zéro d'un polynôme $P(x) = ax^3 + bx^2 + cx + d$ \textbf{à coefficients entiers}, on teste dans l'ordre :
\begin{enumerate}
\item les valeurs simples : $0$, $1$, $-1$, $2$, $-2$, $3$, $-3$ ;
\item si besoin, les diviseurs entiers du terme constant $d$ (lorsque $a = 1$), ou plus généralement les fractions $\dfrac{p}{q}$ où $p$ divise $d$ et $q$ divise $a$ (théorème des racines rationnelles).
\end{enumerate}
Dès qu'une valeur $\alpha$ vérifie $P(\alpha) = 0$, on s'arrête : un seul zéro suffit pour amorcer la factorisation.
\end{methode}

\courssub{Le théorème fondamental : factorisation par $(x - \alpha)$}

\begin{propriete}[Théorème de factorisation (admis)]
Soit $P$ un polynôme et $\alpha$ un nombre réel. Alors :
\[ P(\alpha) = 0 \quad \Longleftrightarrow \quad \text{il existe un polynôme } Q \text{ tel que } P(x) = (x - \alpha)\,Q(x). \]
Autrement dit, $P$ est \cle{factorisable par $(x - \alpha)$} (on dit aussi : $P$ est \cle{divisible par $(x - \alpha)$}) si et seulement si $\alpha$ est un zéro de $P$. De plus, si $P$ est de degré $n$, alors $Q$ est de degré $n - 1$.
\end{propriete}

Ce théorème est admis au programme de Première, mais l'implication facile mérite d'être comprise : si $P(x) = (x - \alpha)Q(x)$, alors en remplaçant $x$ par $\alpha$, on obtient $P(\alpha) = (\alpha - \alpha)Q(\alpha) = 0$. Le sens réciproque — le plus utile — est celui que nous exploiterons : \textbf{chaque zéro trouvé offre un facteur $(x - \alpha)$}.

Conséquence capitale pour un polynôme de degré 3 : si $P(\alpha) = 0$, alors
\[ P(x) = (x - \alpha)\bigl(a x^2 + b'x + c'\bigr), \]
et résoudre $P(x) = 0$ revient à résoudre $(x - \alpha) = 0$ ou une équation du \emph{second degré}. C'est exactement la stratégie annoncée en introduction.

\courssub{Méthode 1 : la division euclidienne en potence}

\begin{methode}[Factoriser par division euclidienne]
Connaissant un zéro $\alpha$ de $P$, on divise $P(x)$ par $(x - \alpha)$ comme on divise des nombres entiers :
\begin{enumerate}
\item on divise le terme dominant de $P$ par $x$ pour obtenir le premier terme du quotient ;
\item on multiplie ce terme par $(x - \alpha)$ et on soustrait le résultat de $P$ ;
\item on recommence avec le reste obtenu, jusqu'à obtenir un reste nul (il est \emph{toujours} nul quand $\alpha$ est un zéro) ;
\item on conclut : $P(x) = (x - \alpha) \times (\text{quotient})$.
\end{enumerate}
\end{methode}

\begin{exemplebox}[Division de $x^3 - 4x^2 + x + 6$ par $(x - 2)$]
On sait que $P(2) = 0$. La division en potence donne :
\begin{itemize}
\item $x^3 \div x = x^2$ ; on calcule $x^2(x-2) = x^3 - 2x^2$, reste : $-2x^2 + x + 6$ ;
\item $-2x^2 \div x = -2x$ ; on calcule $-2x(x-2) = -2x^2 + 4x$, reste : $-3x + 6$ ;
\item $-3x \div x = -3$ ; on calcule $-3(x-2) = -3x + 6$, reste : $0$.
\end{itemize}
Le quotient est $x^2 - 2x - 3$ et le reste est nul, donc
\[ x^3 - 4x^2 + x + 6 = (x - 2)(x^2 - 2x - 3). \]
\end{exemplebox}

\begin{popfigure}
\begin{center}
\begin{tikzpicture}[scale=1, every node/.style={font=\normalsize}]
% Dividende et quotient
\node[anchor=west] (A) at (0,0) {$x^3 - 4x^2 + x + 6$};
\draw[popink, thick] (3.4,0.35) -- (3.4,-2.6);
\node[anchor=west] at (3.55,0) {$x - 2$};
\draw[popink, thick] (3.4,-0.45) -- (5.6,-0.45);
\node[anchor=west] at (3.55,-0.95) {$x^2 - 2x - 3$};
% Étapes de soustraction
\node[anchor=west, popmuted] at (0.35,-0.95) {$-\,(x^3 - 2x^2)$};
\draw[popline] (0.35,-1.25) -- (3.3,-1.25);
\node[anchor=west] at (1.15,-1.7) {$-2x^2 + x + 6$};
\node[anchor=west, popmuted] at (1.15,-2.25) {$-\,(-2x^2 + 4x)$};
\draw[popline] (1.15,-2.55) -- (3.3,-2.55);
\node[anchor=west] at (2.15,-3.0) {$-3x + 6$};
\node[anchor=west, popmuted] at (2.15,-3.55) {$-\,(-3x + 6)$};
\draw[popline] (2.15,-3.85) -- (3.3,-3.85);
\node[anchor=west, poporange] at (2.75,-4.3) {$0$};
\end{tikzpicture}
\end{center}
\legende{Division euclidienne de $x^3 - 4x^2 + x + 6$ par $x - 2$ : le reste est nul car $2$ est un zéro.}
\end{popfigure}

\courssub{Méthode 2 : les coefficients indéterminés}

\begin{methode}[Factoriser par coefficients indéterminés]
Connaissant un zéro $\alpha$ de $P(x) = ax^3 + bx^2 + cx + d$ :
\begin{enumerate}
\item on écrit $P(x) = (x - \alpha)(ax^2 + b'x + c')$, où $b'$ et $c'$ sont inconnus — \textbf{le coefficient dominant du quotient est le même $a$ que celui de $P$} ;
\item on développe le produit et on ordonne selon les puissances de $x$ ;
\item on \cle{identifie} les coefficients : deux polynômes sont égaux si et seulement si leurs coefficients de même degré sont égaux ;
\item on résout le petit système obtenu pour trouver $b'$ et $c'$, puis on vérifie.
\end{enumerate}
\end{methode}

\begin{exemplebox}[Factorisation par identification]
Factorisons $P(x) = x^3 - 4x^2 + x + 6$ sachant que $P(2) = 0$. On écrit :
\[ P(x) = (x - 2)(x^2 + b'x + c'). \]
En développant :
\begin{align*}
(x - 2)(x^2 + b'x + c') &= x^3 + b'x^2 + c'x - 2x^2 - 2b'x - 2c' \\
&= x^3 + (b' - 2)x^2 + (c' - 2b')x - 2c'.
\end{align*}
Par identification avec $x^3 - 4x^2 + x + 6$ :
\[ \begin{cases} b' - 2 = -4 \\ c' - 2b' = 1 \\ -2c' = 6 \end{cases} \quad\Longrightarrow\quad b' = -2, \quad c' = -3, \]
et la deuxième équation confirme : $-3 - 2(-2) = 1$. D'où
\[ P(x) = (x - 2)(x^2 - 2x - 3). \]
\end{exemplebox}

\begin{attention}[Le coefficient dominant du quotient est $a$, pas $1$]
Erreur classique : écrire $P(x) = (x - \alpha)(x^2 + b'x + c')$ alors que $P(x) = 2x^3 + \cdots$ En développant, le terme en $x^3$ provient du produit $x \times (\text{terme en } x^2)$ ; pour obtenir $2x^3$, il faut donc écrire $P(x) = (x - \alpha)(2x^2 + b'x + c')$. De même, le terme constant vérifie $d = -\alpha \times c'$, ce qui donne un contrôle rapide : $c' = -\dfrac{d}{\alpha}$ (si $\alpha \neq 0$).
\end{attention}

\courssub{Factorisation complète et résolution des équations de degré 3}

Une fois le quotient du second degré obtenu, on l'étudie grâce au discriminant. Selon le signe de $\Delta$, trois cas se présentent pour la factorisation \textbf{dans $\mathbb{R}$} :

\begin{itemize}
\item $\Delta > 0$ : le quotient admet deux racines distinctes $x_1, x_2$ ; factorisation complète en trois facteurs du premier degré :
      $P(x) = a(x - \alpha)(x - x_1)(x - x_2)$.
\item $\Delta = 0$ : le quotient admet une racine double $x_0$ ; factorisation :
      $P(x) = a(x - \alpha)(x - x_0)^2$.
\item $\Delta < 0$ : le quotient ne s'annule pas sur $\mathbb{R}$ ; il garde un signe constant (celui de $a$). La factorisation dans $\mathbb{R}$ s'arrête au produit d'un facteur du premier degré et d'un trinôme irréductible :
      $P(x) = (x - \alpha)(ax^2 + b'x + c')$.
\end{itemize}

Résoudre $P(x) = 0$ devient alors immédiat : un produit de facteurs est nul si et seulement si l'un de ses facteurs est nul (règle du produit nul).

\begin{exoresolu}[Factoriser puis résoudre]
Énoncé. Soit $P(x) = x^3 - 4x^2 + x + 6$.
\begin{enumerate}
\item Vérifier que $2$ est un zéro de $P$.
\item Factoriser complètement $P(x)$ dans $\mathbb{R}$.
\item Résoudre dans $\mathbb{R}$ l'équation $P(x) = 0$.
\end{enumerate}
\tcblower
Solution détaillée.
\begin{enumerate}
\item $P(2) = 8 - 16 + 2 + 6 = 0$, donc $2$ est un zéro de $P$.
\item D'après le théorème de factorisation, $P(x) = (x - 2)(x^2 - 2x - 3)$ (voir les deux méthodes ci-dessus). Factorisons le quotient : $\Delta = (-2)^2 - 4 \times 1 \times (-3) = 4 + 12 = 16 > 0$, d'où
\[ x_1 = \frac{2 - 4}{2} = -1, \qquad x_2 = \frac{2 + 4}{2} = 3. \]
La factorisation complète dans $\mathbb{R}$ est donc
\[ P(x) = (x - 2)(x - 3)(x + 1). \]
\item $P(x) = 0 \iff (x - 2)(x - 3)(x + 1) = 0 \iff x = 2$ ou $x = 3$ ou $x = -1$. L'ensemble des solutions est
\[ \mathcal{S} = \{-1 \,;\, 2 \,;\, 3\}. \]
\end{enumerate}
\end{exoresolu}

La figure ci-dessous confirme graphiquement ce résultat : la courbe de $P$ traverse l'axe des abscisses exactement aux points d'abscisses $-1$, $2$ et $3$.

\begin{popfigure}
\begin{center}
\begin{tikzpicture}
\begin{axis}[
  width=0.85\linewidth, height=6.5cm,
  axis lines=middle, axis line style={->, popdark},
  xlabel={$x$}, ylabel={$y$},
  xlabel style={anchor=west}, ylabel style={anchor=south},
  xmin=-2.2, xmax=4.2, ymin=-10, ymax=15,
  xtick={-2,-1,1,2,3,4}, ytick=\empty,
  grid=both, grid style={popline!40, very thin},
  tick label style={font=\small}
]
\addplot[popcurve, domain=-2:4, samples=200]{x^3 - 4*x^2 + x + 6};
\addplot[only marks, mark=*, mark size=2pt, poporange] coordinates {(-1,0) (2,0) (3,0)};
\node[poporange, font=\small, anchor=south] at (axis cs:-1,0.3) {$-1$};
\node[poporange, font=\small, anchor=north west] at (axis cs:2,-0.2) {$2$};
\node[poporange, font=\small, anchor=south west] at (axis cs:3,0.2) {$3$};
\end{axis}
\end{tikzpicture}
\end{center}
\legende{Courbe de $P(x) = x^3 - 4x^2 + x + 6 = (x-2)(x-3)(x+1)$ : les zéros $-1$, $2$, $3$ sont les abscisses des intersections avec l'axe des abscisses.}
\end{popfigure}

\begin{savaistu}[Pourquoi « au plus trois » racines ?]
Un polynôme de degré 3 possède \textbf{au plus trois} zéros réels : chaque zéro fournit un facteur $(x - \alpha)$, et un produit de plus de trois facteurs de degré 1 serait de degré supérieur à 3. Le mathématicien italien Niccolò Tartaglia découvrit au XVI\textsuperscript{e} siècle une formule générale résolvant les équations de degré 3 — formule publiée par Cardan et source d'une célèbre querelle de priorité. En Première, la méthode du zéro évident suivi de la factorisation suffit pour tous les exercices du Probatoire.
\end{savaistu}

\begin{exercice}[À toi de jouer]
Soit $P(x) = 2x^3 - 3x^2 - 3x + 2$.
\begin{enumerate}
\item Vérifier que $-1$ est un zéro de $P$.
\item En utilisant la méthode des coefficients indéterminés, factoriser $P(x)$ complètement dans $\mathbb{R}$.
\item Résoudre dans $\mathbb{R}$ l'équation $P(x) = 0$.
\end{enumerate}
\end{exercice}

\begin{corrige}[Exercice : polynôme $2x^3 - 3x^2 - 3x + 2$]
\begin{enumerate}
\item $P(-1) = 2(-1)^3 - 3(-1)^2 - 3(-1) + 2 = -2 - 3 + 3 + 2 = 0$, donc $-1$ est un zéro de $P$.
\item On écrit $P(x) = (x + 1)(2x^2 + b'x + c')$ — attention, le coefficient dominant du quotient est $2$. En développant :
\[ (x+1)(2x^2 + b'x + c') = 2x^3 + (b' + 2)x^2 + (c' + b')x + c'. \]
Par identification : $b' + 2 = -3$ donne $b' = -5$ ; $c' = 2$ ; vérification : $c' + b' = 2 - 5 = -3$, conforme au coefficient de $x$. Donc
\[ P(x) = (x + 1)(2x^2 - 5x + 2). \]
Pour le quotient : $\Delta = 25 - 16 = 9$, d'où $x_1 = \dfrac{5 - 3}{4} = \dfrac{1}{2}$ et $x_2 = \dfrac{5 + 3}{4} = 2$. Ainsi, la factorisation complète est
\[ P(x) = 2(x + 1)\left(x - \tfrac{1}{2}\right)(x - 2) \quad \text{ou encore} \quad (x+1)(2x-1)(x-2). \]
\item $\mathcal{S} = \left\{-1 \,;\, \dfrac{1}{2} \,;\, 2\right\}$.
\end{enumerate}
\end{corrige}

\begin{aretenir}[L'essentiel de la section]
\begin{itemize}
\item $\alpha$ est un \cle{zéro} de $P$ si et seulement si $P(\alpha) = 0$.
\item \textbf{Théorème (admis)} : $P(\alpha) = 0 \iff P(x) = (x - \alpha)Q(x)$, avec $\deg Q = \deg P - 1$.
\item Pour factoriser un polynôme de degré 3 : chercher un zéro évident (tester $0$, $\pm 1$, $\pm 2$, $\pm 3$, diviseurs de $d$), puis déterminer le quotient par \cle{division euclidienne} ou par \cle{coefficients indéterminés} — le coefficient dominant du quotient est le même $a$ que celui de $P$.
\item Factorisation complète dans $\mathbb{R}$ : on factorise le quotient du second degré selon le signe de son discriminant $\Delta$.
      \begin{itemize}
      \item $\Delta > 0$ : $P(x) = a(x - \alpha)(x - x_1)(x - x_2)$ (trois zéros distincts).
      \item $\Delta = 0$ : $P(x) = a(x - \alpha)(x - x_0)^2$ (un zéro simple, un zéro double).
      \item $\Delta < 0$ : $P(x) = (x - \alpha)(ax^2 + b'x + c')$ (un seul zéro réel).
      \end{itemize}
\item Résoudre $P(x) = 0$ revient alors à appliquer la règle du produit nul. Un polynôme de degré 3 a au plus trois zéros réels.
\end{itemize}
\end{aretenir}

\coursec{Signe d'un polynôme de degré 3 et inéquations}

Dans la section précédente, tu as appris à trouver les zéros d'un polynôme de degré 3 et à le factoriser complètement : par exemple, $P(x) = x^3 - 4x^2 + x + 6$ s'écrit $P(x) = (x+1)(x-2)(x-3)$. Mais à quoi cela sert-il, concrètement ? Imagine un entrepreneur de Bafoussam qui modélise son bénéfice mensuel (en centaines de milliers de FCFA) par un polynôme $B(x)$, où $x$ est le nombre de tonnes de maïs vendues. La question qui l'intéresse n'est pas seulement « quand le bénéfice est-il nul ? » mais surtout : « pour quelles quantités le bénéfice est-il \emph{positif} ? ». Autrement dit, il faut résoudre $B(x) \geq 0$. C'est exactement l'objet de cette section : étudier le \cle{signe} d'un polynôme de degré 3 et résoudre les \cle{inéquations} correspondantes.

\courssub{La propriété fondatrice : constance du signe}

\textbf{Intuition : un polynôme ne peut pas « sauter »}\par

La courbe d'un polynôme est un tracé \emph{continu} : on peut la dessiner sans lever le crayon. Or, pour passer d'une valeur positive à une valeur négative sans lever le crayon, il faut forcément \emph{traverser l'axe des abscisses}, c'est-à-dire s'annuler. Un polynôme ne peut donc changer de signe qu'en passant par l'un de ses zéros. Cette idée intuitive est formalisée par la propriété suivante, que le programme admet.

\begin{propriete}[Constance du signe (admise)]
Si un polynôme $P$ ne s'annule pas sur un intervalle $I$, alors $P$ garde un \cle{signe constant} sur $I$ : soit $P(x) > 0$ pour tout $x$ de $I$, soit $P(x) < 0$ pour tout $x$ de $I$.
\end{propriete}

\begin{aretenir}[Conséquence fondamentale]
Le signe d'un polynôme ne peut changer \textbf{qu'en ses zéros}. Par conséquent, pour connaître le signe d'un polynôme sur $\mathbb{R}$, il suffit de :
\begin{enumerate}
\item trouver tous ses zéros, qui découpent $\mathbb{R}$ en intervalles ;
\item déterminer le signe du polynôme en \emph{un seul} point de chaque intervalle (ou utiliser les signes des facteurs).
\end{enumerate}
\end{aretenir}

C'est exactement la même philosophie que pour le trinôme du second degré : les racines sont les « frontières » où le signe peut basculer. La différence, c'est qu'un polynôme de degré 3 peut avoir jusqu'à trois zéros, donc jusqu'à quatre intervalles à étudier.

\courssub{Tableau de signes d'un polynôme de degré 3 factorisé}

\textbf{Principe de construction}\par

Lorsque le polynôme est entièrement factorisé en produit de facteurs du premier degré, $P(x) = a(x - \alpha)(x - \beta)(x - \gamma)$, le plus sûr est de dresser un tableau de signes \emph{ligne par ligne} : une ligne pour chaque facteur, puis une ligne pour le produit, obtenue par la règle des signes. Si le coefficient $a$ est négatif, on ajoute une ligne pour son signe (constant) ou on l'intègre dans la conclusion.

\begin{methode}[Dresser le tableau de signes d'un polynôme de degré 3]
\begin{enumerate}
\item \textbf{Factoriser} complètement le polynôme (zéro évident, puis factorisation du quotient par division euclidienne ou coefficients indéterminés).
\item \textbf{Ranger les zéros} dans l'ordre croissant : $\alpha < \beta < \gamma$.
\item \textbf{Construire le tableau} : une ligne par facteur $(x - \alpha)$, $(x - \beta)$, $(x - \gamma)$, chacun s'annulant en son zéro, négatif avant, positif après.
\item \textbf{Multiplier les signes} colonne par colonne pour obtenir la ligne de $P(x)$, en plaçant des $0$ sous chaque zéro.
\end{enumerate}
\end{methode}

\begin{exemplebox}[Tableau de signes de $P(x) = (x+1)(x-2)(x-3)$]
Les zéros, rangés dans l'ordre croissant, sont $-1$, $2$ et $3$. On obtient :

\begin{center}
\begin{tabular}{|c|ccccccc|}
\hline
$x$ & $-\infty$ & & $-1$ & & $2$ & & $3$ \\
\hline
$x+1$ & & $-$ & $0$ & $+$ & $|$ & $+$ & $|$ \\
\hline
$x-2$ & & $-$ & $|$ & $-$ & $0$ & $+$ & $|$ \\
\hline
$x-3$ & & $-$ & $|$ & $-$ & $|$ & $-$ & $0$ \\
\hline
$P(x)$ & & $-$ & $0$ & $+$ & $0$ & $-$ & $0$ \\
\hline
\end{tabular}
\end{center}

Vérifions deux colonnes. Pour $x > 3$, les trois facteurs sont positifs, donc $P(x) > 0$ ; en effet $P(4) = 5 \times 2 \times 1 = 10 > 0$. Sur $]2\,;\,3[$, deux facteurs positifs et un négatif donnent un produit négatif : $P(2{,}5) = 3{,}5 \times 0{,}5 \times (-0{,}5) = -0{,}875 < 0$. Tout est cohérent.
\end{exemplebox}

\begin{attention}[Un zéro double ne change pas le signe]
Si la factorisation contient un facteur $(x - \alpha)^2$, ce facteur est \textbf{toujours positif ou nul} : il s'annule en $\alpha$ mais ne change jamais de signe. Le polynôme garde donc le \emph{même signe} de part et d'autre d'un zéro double. Par exemple, pour $Q(x) = (x-1)^2(x+2)$ : le signe de $Q$ est celui de $(x+2)$, sauf en $x = 1$ où $Q$ s'annule sans changer de signe. Oublier cette règle et « faire alterner » les signes mécaniquement est l'erreur la plus fréquente au Probatoire.
\end{attention}

\courssub{Lecture graphique du signe}

Le tableau de signes se lit directement sur la courbe : le polynôme est positif là où sa courbe est \emph{au-dessus} de l'axe des abscisses, négatif là où elle est \emph{en dessous}.

\begin{popfigure}
\begin{center}
\begin{tikzpicture}
\begin{axis}[
  width=0.9\linewidth, height=7cm,
  axis lines=middle, xlabel={$x$}, ylabel={$y$},
  xmin=-2.2, xmax=4.2, ymin=-5, ymax=6,
  xtick={-1,2,3}, ytick=\empty,
  domain=-1.8:3.8, samples=200, clip=false
]
\addplot[popcurve]{(x-2)*(x-3)*(x+1)};
\addplot[poporange, very thick, domain=-1:2, samples=100]{(x-2)*(x-3)*(x+1)};
\addplot[popgreen, very thick, domain=-1.8:-1, samples=50]{(x-2)*(x-3)*(x+1)};
\addplot[popgreen, very thick, domain=2:3, samples=50]{(x-2)*(x-3)*(x+1)};
\addplot[poporange, very thick, domain=3:3.8, samples=50]{(x-2)*(x-3)*(x+1)};
\addplot[only marks, mark=*, popdark] coordinates {(-1,0) (2,0) (3,0)};
\node[popgreen, font=\small] at (axis cs:-1.5,-2.5) {$P<0$};
\node[poporange, font=\small] at (axis cs:0.5,3) {$P>0$};
\node[popgreen, font=\small] at (axis cs:2.5,-1.8) {$P<0$};
\node[poporange, font=\small] at (axis cs:3.6,4) {$P>0$};
\end{axis}
\end{tikzpicture}
\end{center}
\legende{Courbe de $P(x) = (x+1)(x-2)(x-3)$ : en orange les intervalles où $P(x) > 0$, en vert ceux où $P(x) < 0$. Les zéros $-1$, $2$, $3$ sont les seuls points où le signe peut changer.}
\end{popfigure}

\courssub{Résolution d'inéquations de degré 3}

Résoudre $P(x) \geq 0$ ou $P(x) < 0$ consiste simplement à \emph{lire la conclusion} sur le tableau de signes. La solution s'écrit comme une \textbf{réunion d'intervalles}.

\begin{methode}[Résoudre une inéquation polynomiale]
\begin{enumerate}
\item \textbf{Transposer} pour se ramener à la forme $P(x) \geq 0$ (ou $>$, $\leq$, $<$).
\item \textbf{Factoriser} $P$ complètement et trouver tous ses zéros.
\item \textbf{Dresser le tableau de signes} de $P$ sur $\mathbb{R}$.
\item \textbf{Conclure en intervalles} : retenir les intervalles où le signe convient ; \emph{fermer} les bornes si l'inégalité est large ($\geq$ ou $\leq$), les \emph{ouvrir} si elle est stricte ($>$ ou $<$).
\end{enumerate}
\end{methode}

\begin{exoresolu}[Résoudre $x^3 - 4x^2 + x + 6 \leq 0$]
Résoudre dans $\mathbb{R}$ l'inéquation $x^3 - 4x^2 + x + 6 \leq 0$.
\tcblower
\textbf{Étape 1 : factorisation.} Posons $P(x) = x^3 - 4x^2 + x + 6$. On teste les zéros évidents : $P(-1) = -1 - 4 - 1 + 6 = 0$. Donc $-1$ est un zéro et $P$ est divisible par $(x+1)$. Par identification, $P(x) = (x+1)(x^2 - 5x + 6)$. Le trinôme $x^2 - 5x + 6$ a pour discriminant $\Delta = 25 - 24 = 1$, d'où les racines $2$ et $3$. Ainsi :
\[ P(x) = (x+1)(x-2)(x-3). \]

\textbf{Étape 2 : tableau de signes.} C'est le tableau construit plus haut : $P$ est strictement négatif sur $]-\infty\,;\,-1[$ et sur $]2\,;\,3[$, strictement positif sur $]-1\,;\,2[$ et sur $]3\,;\,+\infty[$, et nul en $-1$, $2$, $3$.

\textbf{Étape 3 : conclusion.} L'inégalité est large, donc on inclut les zéros :
\[ \mathcal{S} = \left]-\infty\,;\,-1\right] \cup \left[2\,;\,3\right]. \]
\end{exoresolu}

\begin{exemplebox}[Inéquation par transposition : $x^3 \geq 3x + 2$]
Il faut d'abord \emph{tout ramener dans un même membre} :
\[ x^3 \geq 3x + 2 \iff x^3 - 3x - 2 \geq 0. \]
Posons $P(x) = x^3 - 3x - 2$. Zéro évident : $P(-1) = -1 + 3 - 2 = 0$. On factorise : $P(x) = (x+1)(x^2 - x - 2)$. Le trinôme $x^2 - x - 2$ a pour discriminant $\Delta = 1 + 8 = 9$, donc pour racines $-1$ et $2$. Donc :
\[ P(x) = (x+1)^2(x-2). \]
Le facteur $(x+1)^2$ est toujours positif ou nul : le signe de $P$ est celui de $(x-2)$, avec annulation en $-1$ (zéro double, sans changement de signe) et en $2$. Donc $P(x) \geq 0$ si et seulement si $x = -1$ ou $x \geq 2$ :
\[ \mathcal{S} = \{-1\} \cup \left[2\,;\,+\infty\right[. \]
Remarque le point isolé $x = -1$ : il est solution car l'inégalité est large, même si le signe ne change pas autour de lui.
\end{exemplebox}

\begin{attention}[Pièges de rédaction]
\begin{itemize}
\item Ne jamais « diviser par $(x - a)$ » dans une inéquation sans discuter le signe de $(x-a)$ : on perd ou on fausse des solutions. La bonne méthode est toujours : transposer, factoriser, tableau de signes.
\item Les solutions d'une inéquation s'écrivent en \textbf{intervalles}, jamais sous la forme « $x \leq -1$ ou $2 \leq x \leq 3$ » laissée en l'état sans conclusion ensembliste.
\item Avec une inégalité stricte, les zéros sont \emph{exclus} : pour $P(x) < 0$ avec $P(x) = (x+1)(x-2)(x-3)$, on obtient $]-\infty\,;\,-1[ \cup ]2\,;\,3[$.
\end{itemize}
\end{attention}

\courssub{Extension : signe d'un produit ou d'un quotient}

La même technique s'applique à toute expression dont on connaît tous les zéros : produits de polynômes, mais aussi \textbf{quotients}. Pour un quotient, il y a une règle supplémentaire : les zéros du \emph{dénominateur} sont des \cle{valeurs interdites}, marquées par une double barre dans le tableau, et toujours exclues des solutions.

\begin{exemplebox}[Signe d'un quotient]
Étudions le signe de $f(x) = \dfrac{(x-1)(x-3)}{x+2}$. Les zéros du numérateur sont $1$ et $3$ ; la valeur interdite est $-2$.

\begin{center}
\begin{tabular}{|c|ccccccc|}
\hline
$x$ & $-\infty$ & & $-2$ & & $1$ & & $3$ \\
\hline
$(x-1)(x-3)$ & & $+$ & $|$ & $+$ & $0$ & $-$ & $0$ \\
\hline
$x+2$ & & $-$ & $0$ & $+$ & $|$ & $+$ & $|$ \\
\hline
$f(x)$ & & $-$ & $\|$ & $+$ & $0$ & $-$ & $0$ \\
\hline
\end{tabular}
\end{center}

Ainsi $f(x) \geq 0$ sur $]-2\,;\,1] \cup [3\,;\,+\infty[$ : la borne $-2$ reste \emph{ouverte} car elle est interdite, même avec une inégalité large.
\end{exemplebox}

\begin{savaistu}[Cardan et la plus célèbre querelle des mathématiques]
Au XVI\textsuperscript{e} siècle, les mathématiciens italiens se lançaient des défis publics de résolution d'équations. Niccolò Tartaglia découvrit comment résoudre les équations de degré 3 et confia son secret, sous serment, à Jérôme Cardan. Celui-ci publia la méthode en 1545 dans son \emph{Ars Magna} — en citant Tartaglia, mais en brisant sa promesse. S'ensuivit l'une des querelles les plus violentes de l'histoire des sciences. Ironie du sort : la « formule de Cardan » oblige parfois à manipuler des racines carrées de nombres négatifs pour obtenir des solutions pourtant bien réelles — ce qui poussa les mathématiciens à inventer les nombres complexes, que tu étudieras en Terminale !
\end{savaistu}

\begin{exercice}[À toi de jouer]
\begin{enumerate}
\item Dresser le tableau de signes de $P(x) = x^3 - 2x^2 - 5x + 6$ (on vérifiera que $1$ est un zéro), puis résoudre $P(x) > 0$.
\item Résoudre dans $\mathbb{R}$ l'inéquation $x^3 \leq 4x$.
\item Résoudre $\dfrac{x^2 - 4}{x - 1} \geq 0$.
\end{enumerate}
\end{exercice}

\begin{corrige}[Exercice]
\textbf{1.} $P(1) = 1 - 2 - 5 + 6 = 0$, donc $P(x) = (x-1)(x^2 - x - 6) = (x-1)(x+2)(x-3)$. Les zéros rangés sont $-2$, $1$, $3$. Le tableau donne $P$ strictement négatif sur $]-\infty\,;\,-2[$ et $]1\,;\,3[$, strictement positif sur $]-2\,;\,1[$ et $]3\,;\,+\infty[$. L'inégalité étant stricte : $\mathcal{S} = ]-2\,;\,1[ \cup ]3\,;\,+\infty[$.

\textbf{2.} $x^3 \leq 4x \iff x^3 - 4x \leq 0 \iff x(x-2)(x+2) \leq 0$. Les zéros sont $-2$, $0$, $2$. Le produit de ces trois facteurs du premier degré est négatif sur $]-\infty\,;\,-2[$ et $]0\,;\,2[$, positif sur $]-2\,;\,0[$ et $]2\,;\,+\infty[$. Inégalité large : $\mathcal{S} = ]-\infty\,;\,-2] \cup [0\,;\,2]$.

\textbf{3.} $x^2 - 4 = (x-2)(x+2)$ s'annule en $-2$ et $2$ ; valeur interdite : $1$. Le tableau de signes du quotient donne un signe positif ou nul sur $[-2\,;\,1[$ et $[2\,;\,+\infty[$ (la borne $1$ est ouverte car interdite). $\mathcal{S} = [-2\,;\,1[ \cup [2\,;\,+\infty[$.
\end{corrige}

Tu maîtrises désormais l'étude du signe des polynômes et la résolution des inéquations polynomiales. Dans la dernière section, nous verrons comment ces mêmes techniques — transposition, factorisation, tableau de signes — permettent de dompter les équations et inéquations contenant des racines carrées, dites \emph{irrationnelles}.

\coursec{Équations et inéquations irrationnelles}

Dans la section précédente, tu as appris à étudier le signe d'un polynôme de degré $3$ et à résoudre des inéquations du type $P(x) \geq 0$. Toutes les expressions rencontrées jusqu'ici étaient \emph{rationnelles} : l'inconnue n'apparaissait que dans des sommes, des produits et des quotients de polynômes. Nous franchissons maintenant une dernière étape : celle où l'inconnue se cache \emph{sous un radical}. C'est une situation très concrète : quand un technicien de CAMTEL calcule la longueur $d$ d'un câble tendu entre deux poteaux à l'aide de la formule $d = \sqrt{x^2 + h^2}$, l'inconnue $x$ (distance horizontale) figure sous une racine carrée. Comment résoudre alors une équation ? C'est l'objet de cette section, qui clôt notre leçon.

\courssub{Notion d'équation irrationnelle et conditions d'existence}

\begin{definition}[Équation et inéquation irrationnelles]
On appelle \cle{équation irrationnelle} (respectivement \cle{inéquation irrationnelle}) toute équation (respectivement inéquation) dans laquelle l'\emph{inconnue figure sous un signe radical}, par exemple sous une racine carrée.
\end{definition}

\begin{exemplebox}[Reconnaître une équation irrationnelle]
\begin{itemize}
\item $\sqrt{x+2} = x$ est une équation irrationnelle : $x$ est sous le radical.
\item $\sqrt{2x+3} = x$ et $\sqrt{x+2} \leq x$ sont respectivement une équation et une inéquation irrationnelles.
\item $x^2 - 3x + 2 = \sqrt{2}$ n'est \emph{pas} une équation irrationnelle : le radical porte sur un nombre, pas sur l'inconnue. C'est une simple équation du second degré.
\end{itemize}
\end{exemplebox}

Avant toute chose, rappelons un point fondamental sur la racine carrée, qui gouvernera toute cette section.

\begin{aretenir}[Les deux règles d'or de la racine carrée]
Pour tout réel $u$ :
\begin{enumerate}
\item \textbf{Existence :} $\sqrt{u}$ n'existe (dans $\mathbb{R}$) que si $u \geq 0$ ;
\item \textbf{Positivité :} quand elle existe, $\sqrt{u}$ est un nombre \emph{positif ou nul} : $\sqrt{u} \geq 0$.
\end{enumerate}
\end{aretenir}

Il en découle immédiatement la première étape de toute résolution.

\begin{definition}[Condition d'existence]
Pour que l'expression $\sqrt{ax+b}$ ait un sens, il faut et il suffit que
\[ ax + b \geq 0. \]
Cette condition définit le \cle{domaine de résolution} de l'équation : toute solution cherchée devra obligatoirement lui appartenir.
\end{definition}

\begin{exemplebox}[Déterminer un domaine]
Pour l'équation $\sqrt{2x+3} = x$, la condition d'existence est $2x+3 \geq 0$, c'est-à-dire $x \geq -\dfrac{3}{2}$. Le domaine de résolution est donc $\mathcal{D} = \left[-\dfrac{3}{2} \, ; +\infty\right[$. Aucun nombre inférieur à $-\dfrac{3}{2}$ ne peut prétendre être solution.
\end{exemplebox}

\courssub{Résolution de l'équation $\sqrt{ax+b} = cx+d$}

\textbf{Intuition.} Une racine carrée est toujours positive. Donc si $\sqrt{ax+b} = cx+d$, le membre de droite $cx+d$ est forcément positif ou nul. Sous cette condition, les deux membres sont positifs, et on peut alors élever au carré : deux nombres \emph{positifs} sont égaux si et seulement si leurs carrés sont égaux. Toute la difficulté est là : l'élévation au carré n'est légitime que si l'on a d'abord garanti la positivité des deux membres.

\begin{propriete}[Résolution de $\sqrt{ax+b} = cx+d$]
Résoudre l'équation $(E) : \sqrt{ax+b} = cx+d$ revient à résoudre le \cle{système} :
\[
\begin{cases}
cx + d \geq 0 & \text{(condition de positivité)} \\
ax + b = (cx+d)^2 & \text{(équation obtenue en élevant au carré)}
\end{cases}
\]
La condition $ax+b \geq 0$ est automatiquement satisfaite dès que $ax+b = (cx+d)^2$, car un carré est toujours positif ou nul.
\end{propriete}

\begin{experience}[Pourquoi ce système est-il équivalent à l'équation ?]
Raisonnons en deux sens.
\begin{enumerate}
\item \textbf{Sens direct.} Supposons que $x$ vérifie $\sqrt{ax+b} = cx+d$. Comme une racine carrée est positive, on a $cx+d \geq 0$ : la première ligne du système est vraie. En élevant les deux membres (égaux) au carré, on obtient $ax+b = (cx+d)^2$ : la deuxième ligne est vraie.
\item \textbf{Sens réciproque.} Supposons que $x$ vérifie le système. De $ax+b = (cx+d)^2$ on tire, en prenant la racine carrée (licite car $ax+b = (cx+d)^2 \geq 0$) : $\sqrt{ax+b} = |cx+d|$. Or la condition $cx+d \geq 0$ donne $|cx+d| = cx+d$, donc $\sqrt{ax+b} = cx+d$.
\end{enumerate}
Le système est donc bien \emph{équivalent} à l'équation $(E)$. Remarque le rôle décisif de la condition $cx+d \geq 0$ : c'est elle qui transforme $|cx+d|$ en $cx+d$.
\end{experience}

\begin{attention}[L'élévation au carré n'est pas une équivalence !]
Si l'on se contente d'élever au carré \emph{sans aucune précaution}, on résout en réalité $ax+b = (cx+d)^2$, c'est-à-dire $\sqrt{ax+b} = |cx+d|$, ce qui regroupe \emph{deux} équations : $\sqrt{ax+b} = cx+d$ \textbf{et} $\sqrt{ax+b} = -(cx+d)$. On risque donc de récolter des \cle{solutions étrangères}, qui ne vérifient pas l'équation de départ. Deux réflexes obligatoires :
\begin{itemize}
\item ne jamais oublier la condition $cx+d \geq 0$, car une racine carrée est toujours positive ;
\item en cas de doute, \emph{vérifier} chaque solution en la reportant dans l'équation initiale.
\end{itemize}
\end{attention}

\begin{methode}[Résoudre $\sqrt{ax+b} = cx+d$]
\begin{enumerate}
\item \textbf{Domaine :} écrire la condition d'existence $ax+b \geq 0$ et la condition de positivité $cx+d \geq 0$.
\item \textbf{Élévation au carré :} résoudre l'équation $ax+b = (cx+d)^2$, qui est une équation du second degré.
\item \textbf{Filtrage :} ne retenir que les solutions qui appartiennent au domaine \emph{et} vérifient $cx+d \geq 0$ ; le plus sûr est de les reporter dans l'équation initiale.
\end{enumerate}
\end{methode}

\begin{exoresolu}[Résoudre $\sqrt{2x+3} = x$]
Résoudre dans $\mathbb{R}$ l'équation $\sqrt{2x+3} = x$.
\tcblower
\textbf{Étape 1 : conditions.} Existence : $2x+3 \geq 0$, soit $x \geq -\dfrac{3}{2}$. Positivité du membre de droite : $x \geq 0$. On cherchera donc les solutions dans $[0 \, ; +\infty[$.

\textbf{Étape 2 : élévation au carré.}
\begin{align*}
2x+3 &= x^2 \\
x^2 - 2x - 3 &= 0
\end{align*}
Discriminant : $\Delta = (-2)^2 - 4(1)(-3) = 4 + 12 = 16 > 0$. Deux racines :
\[ x_1 = \frac{2 - 4}{2} = -1 \qquad \text{et} \qquad x_2 = \frac{2+4}{2} = 3. \]

\textbf{Étape 3 : filtrage.} La condition $x \geq 0$ élimine $x_1 = -1$. Vérification de $x_2 = 3$ dans l'équation initiale : $\sqrt{2(3)+3} = \sqrt{9} = 3$ : c'est bien une solution.

\textbf{Conclusion :} $\mathcal{S} = \{3\}$.
\end{exoresolu}

\begin{attention}[Une solution étrangère classique]
Résolvons $\sqrt{x+2} = x$. L'élévation au carré donne $x+2 = x^2$, soit $x^2 - x - 2 = 0$, de racines $x = 2$ et $x = -1$. Si l'on s'arrête là, on écrit $\mathcal{S} = \{-1 \, ; 2\}$ : c'est \textbf{faux} ! En effet, pour $x = -1$ : $\sqrt{-1+2} = \sqrt{1} = 1 \neq -1$. Le nombre $-1$ est une \emph{solution étrangère} introduite par l'élévation au carré. La condition $x \geq 0$ l'aurait éliminée d'office. Conclusion : $\mathcal{S} = \{2\}$.
\end{attention}

La figure suivante permet de \emph{voir} pourquoi il n'y a qu'une seule solution : les courbes d'équations $y = \sqrt{x+2}$ et $y = x$ ne se coupent qu'en un point, d'abscisse $2$.

\begin{popfigure}
\begin{center}
\begin{tikzpicture}
\begin{axis}[
  width=0.85\linewidth, height=6.5cm,
  axis lines=middle,
  xlabel={$x$}, ylabel={$y$},
  xmin=-2.5, xmax=4.5, ymin=-1.5, ymax=3.5,
  xtick={-2,-1,1,2,3,4}, ytick={1,2,3},
  grid=both, grid style={popline!40},
  legend style={at={(0.02,0.98)}, anchor=north west, font=\small}
]
\addplot[popcurve, popblue, domain=-2:4.2, samples=200]{sqrt(max(0,x+2))};
\addlegendentry{$y=\sqrt{x+2}$}
\addplot[popcurve, poporange, domain=-1.5:3.5, samples=2]{x};
\addlegendentry{$y=x$}
\addplot[only marks, mark=*, mark size=2.5pt, popdark] coordinates {(2,2)};
\node[popdark, anchor=south west, font=\small] at (axis cs:2,2) {$(2\,;\,2)$};
\addplot[only marks, mark=o, mark size=2.5pt, poppink] coordinates {(-1,1)};
\node[poppink, anchor=south east, font=\small] at (axis cs:-1,1) {pas d'intersection en $-1$};
\end{axis}
\end{tikzpicture}
\end{center}
\legende{Les courbes $y=\sqrt{x+2}$ et $y=x$ ne se coupent qu'au point $(2\,;\,2)$ : l'équation $\sqrt{x+2}=x$ a pour unique solution $x=2$. Le point $(-1\,;\,1)$ appartient à la courbe de la racine mais pas à la droite : $-1$ est une solution étrangère.}
\end{popfigure}

\courssub{Résolution de l'inéquation $\sqrt{ax+b} \leq cx+d$}

\textbf{Intuition.} On compare un nombre positif ou nul (la racine) à $cx+d$. Si $cx+d$ était strictement négatif, l'inégalité $\sqrt{ax+b} \leq cx+d$ serait impossible, car un nombre positif ou nul ne peut pas être inférieur ou égal à un nombre strictement négatif. Il faut donc d'abord exiger $cx+d \geq 0$. Une fois les deux membres positifs, la fonction carré étant \emph{strictement croissante} sur $[0\,;\,+\infty[$, on peut élever au carré \emph{sans changer le sens de l'inégalité}.

\begin{propriete}[Résolution de $\sqrt{ax+b} \leq cx+d$]
L'inéquation $\sqrt{ax+b} \leq cx+d$ équivaut au système :
\[
\begin{cases}
ax + b \geq 0 \\
cx + d \geq 0 \\
ax + b \leq (cx+d)^2
\end{cases}
\]
La dernière ligne est une inéquation du second degré, que l'on résout par un tableau de signes (section 2).
\end{propriete}

\begin{exoresolu}[Résoudre $\sqrt{x+2} \leq x$]
Résoudre dans $\mathbb{R}$ l'inéquation $\sqrt{x+2} \leq x$.
\tcblower
\textbf{Conditions :} $x+2 \geq 0$ (soit $x \geq -2$) et $x \geq 0$. On travaille donc sur $[0\,;\,+\infty[$.

\textbf{Élévation au carré :} $x+2 \leq x^2$, soit $x^2 - x - 2 \geq 0$. Le trinôme $x^2 - x - 2$ a pour racines $-1$ et $2$ ; il est positif à l'extérieur des racines (coefficient de $x^2$ positif) :
\[ x^2 - x - 2 \geq 0 \iff x \in \left]-\infty \, ; -1\right] \cup \left[2 \, ; +\infty\right[. \]

\textbf{Intersection avec le domaine $[0\,;\,+\infty[$ :}
\[ \mathcal{S} = \left[2 \, ; +\infty\right[. \]
Graphiquement, cela correspond exactement à la zone où la courbe $y=\sqrt{x+2}$ est \emph{sous} la droite $y=x$ (voir la figure précédente) : à partir de $x=2$.
\end{exoresolu}

\courssub{Résolution de l'inéquation $\sqrt{ax+b} \geq cx+d$ : disjonction de cas}

Ici, la situation est plus subtile, car le signe de $cx+d$ change tout.

\begin{itemize}
\item \textbf{Si $cx+d < 0$ :} l'inégalité $\sqrt{ax+b} \geq cx+d$ est \emph{automatiquement vraie} dès que la racine existe, puisqu'un nombre positif ou nul est toujours supérieur à un nombre strictement négatif. Il suffit donc d'exiger $ax+b \geq 0$.
\item \textbf{Si $cx+d \geq 0$ :} les deux membres sont positifs, on peut élever au carré en conservant le sens : $ax+b \geq (cx+d)^2$.
\end{itemize}

\begin{propriete}[Résolution de $\sqrt{ax+b} \geq cx+d$]
L'ensemble des solutions de $\sqrt{ax+b} \geq cx+d$ est la \emph{réunion} des ensembles de solutions des deux systèmes suivants :
\[
(S_1) : \begin{cases} ax+b \geq 0 \\ cx+d < 0 \end{cases}
\qquad \text{ou} \qquad
(S_2) : \begin{cases} cx+d \geq 0 \\ ax+b \geq (cx+d)^2 \end{cases}
\]
\end{propriete}

\begin{exemplebox}[Résoudre $\sqrt{x+2} \geq x$]
\textbf{Système $(S_1)$ :} $x+2 \geq 0$ et $x < 0$, soit $x \in [-2 \, ; 0[$.

\textbf{Système $(S_2)$ :} $x \geq 0$ et $x+2 \geq x^2$, c'est-à-dire $x^2 - x - 2 \leq 0$, soit $x \in [-1 \, ; 2]$. Intersection avec $x \geq 0$ : $x \in [0 \, ; 2]$.

\textbf{Réunion :} $\mathcal{S} = [-2\,;\,0[ \cup [0\,;\,2] = [-2 \, ; 2]$. C'est cohérent avec la figure : la courbe $y=\sqrt{x+2}$ est au-dessus de la droite $y=x$ entre $x=-2$ et $x=2$.
\end{exemplebox}

\begin{attention}[Les deux pièges des inéquations irrationnelles]
\begin{itemize}
\item Pour $\sqrt{ax+b} \leq cx+d$, ne jamais oublier la condition $cx+d \geq 0$ : sans elle, on récolte des solutions où le membre de droite est négatif, ce qui est absurde.
\item Pour $\sqrt{ax+b} \geq cx+d$, ne pas oublier le \emph{premier cas} ($cx+d < 0$) : tous les réels du domaine avec $cx+d < 0$ sont solutions, et les oublier ampute l'ensemble des solutions.
\end{itemize}
\end{attention}

\courssub{Synthèse : la méthode générale en trois temps}

\begin{aretenir}[Méthode générale : domaine, résolution, filtrage]
Face à toute équation ou inéquation irrationnelle :
\begin{enumerate}
\item \textbf{Domaine :} écrire les conditions d'existence des radicaux ($ax+b \geq 0$) et les conditions de signe imposées par la positivité de la racine ;
\item \textbf{Résolution :} élever au carré (licite quand les deux membres sont positifs) pour se ramener à une équation ou inéquation du second degré, que l'on résout avec le discriminant et un tableau de signes ;
\item \textbf{Filtrage :} ne retenir que les solutions compatibles avec toutes les conditions, ou vérifier directement chaque solution dans la relation initiale.
\end{enumerate}
\end{aretenir}

\begin{popfigure}
\begin{center}
\begin{tikzpicture}[
  node distance=0.55cm,
  etape/.style={rectangle, rounded corners, draw=popblue, fill=popblueL, thick, text width=6.2cm, align=center, minimum height=1cm, font=\small},
  decision/.style={diamond, draw=poporange, fill=poporangeL, thick, aspect=2.2, align=center, font=\small, inner sep=1pt},
  fleche/.style={-{Stealth[length=3mm]}, thick, popdark}
]
\node[etape, align=center] (dom) {\textbf{1. Domaine} \\ $ax+b \geq 0$ \; et \; $cx+d \geq 0$ (si nécessaire)};
\node[etape, below=of dom, align=center] (carre) {\textbf{2. Élévation au carré} \\ $ax+b = (cx+d)^2$ : équation du second degré};
\node[etape, below=of carre, align=center] (reso) {\textbf{3. Résolution} \\ discriminant $\Delta$, racines $x_1, x_2$};
\node[decision, below=of reso, align=center] (test) {Chaque racine\\ vérifie-t-elle\\ les conditions ?};
\node[etape, below left=0.6cm and -1.2cm of test, fill=popgreenL, draw=popgreen, align=center] (oui) {\textbf{Solution retenue} \\ elle appartient à $\mathcal{S}$};
\node[etape, below right=0.6cm and -1.2cm of test, fill=poppinkL, draw=poppink, align=center] (non) {\textbf{Solution étrangère} \\ elle est rejetée};
\draw[fleche] (dom) -- (carre);
\draw[fleche] (carre) -- (reso);
\draw[fleche] (reso) -- (test);
\draw[fleche] (test) -| node[pos=0.25, above, font=\small]{oui} (oui);
\draw[fleche] (test) -| node[pos=0.25, above, font=\small]{non} (non);
\end{tikzpicture}
\end{center}
\legende{Organigramme de résolution d'une équation irrationnelle : le filtrage final est l'étape qui distingue une copie correcte d'une copie fautive.}
\end{popfigure}

\begin{savaistu}[Pourquoi « irrationnelle » ?]
Le mot vient du latin \emph{ratio} (raison, rapport) : les Grecs appelaient « irrationnels » les nombres comme $\sqrt{2}$, qui ne sont pas le rapport de deux entiers. La découverte de l'irrationalité de $\sqrt{2}$, attribuée à l'école de Pythagore, fut un véritable scandale mathématique ! Aujourd'hui, les équations irrationnelles apparaissent dès qu'on manipule des distances : longueur d'un câble électrique entre deux poteaux, portée d'une antenne MTN ($d = \sqrt{2Rh}$ où $h$ est la hauteur de l'antenne), vitesse d'écoulement de l'eau dans un canal. Savoir les résoudre proprement, c'est savoir filtrer les solutions mathématiques pour ne garder que celles qui ont un sens physique.
\end{savaistu}

\begin{exercice}[Pour t'entraîner]
\begin{enumerate}
\item Résoudre dans $\mathbb{R}$ : $\sqrt{3x-2} = x-2$.
\item Résoudre dans $\mathbb{R}$ : $\sqrt{x+7} \leq x+1$.
\item Résoudre dans $\mathbb{R}$ : $\sqrt{2x+5} \geq x+1$.
\end{enumerate}
\end{exercice}

\begin{corrige}[Exercice 1]
\textbf{1.} Conditions : $3x-2 \geq 0$ et $x-2 \geq 0$, soit $x \geq 2$. Élévation au carré : $3x-2 = (x-2)^2 = x^2 - 4x + 4$, soit $x^2 - 7x + 6 = 0$. $\Delta = 49 - 24 = 25$ ; racines $x = 1$ et $x = 6$. Filtrage ($x \geq 2$) : on rejette $1$. Vérification : $\sqrt{18-2} = 4 = 6-2$. $\mathcal{S} = \{6\}$.

\textbf{2.} Conditions : $x+7 \geq 0$ et $x+1 \geq 0$, soit $x \geq -1$. Élévation au carré : $x+7 \leq x^2+2x+1$, soit $x^2+x-6 \geq 0$. Racines $-3$ et $2$ ; le trinôme est positif à l'extérieur : $x \leq -3$ ou $x \geq 2$. Intersection avec $x \geq -1$ : $\mathcal{S} = [2 \, ; +\infty[$.

\textbf{3.} Cas 1 : $2x+5 \geq 0$ et $x+1 < 0$, soit $x \in \left[-\dfrac{5}{2} \, ; -1\right[$. Cas 2 : $x+1 \geq 0$ et $2x+5 \geq x^2+2x+1$, soit $x^2 \leq 4$, c'est-à-dire $x \in [-2\,;\,2]$, intersecté avec $x \geq -1$ : $x \in [-1\,;\,2]$. Réunion : $\mathcal{S} = \left[-\dfrac{5}{2} \, ; 2\right]$.
\end{corrige}

\newpage

\coursec{Activité d'intégration}

\courssub{Objectifs de l'activité}

À l'issue de cette activité d'intégration, tu seras capable de mobiliser, dans une situation économique authentique, l'ensemble des savoir-faire de la leçon :

\begin{itemize}
\item déterminer et factoriser un \cle{polynôme de degré 3} en utilisant un \cle{zéro} (méthode de la division euclidienne ou des coefficients indéterminés) ;
\item résoudre une \cle{équation du second degré} à l'aide du discriminant ;
\item dresser le \cle{tableau de signes} d'un polynôme et résoudre une inéquation de degré 3 ;
\item résoudre une \cle{équation irrationnelle} du type $\sqrt{ax+b}=cx+d$ en contrôlant la validité des solutions ;
\item rédiger un argumentaire mathématique destiné à un décideur non mathématicien (le président d'une coopérative).
\end{itemize}

\courssub{Situation}

La coopérative agricole \emph{Nkouam et Frères} de Mbouda, dans la région de l'Ouest, pratique le maraîchage (carottes, choux, pommes de terre). La saison sèche menace les rendements : les membres décident de forer un puits équipé d'une pompe.

Le bureau d'études hydrauliques de Bafoussam a établi les modèles suivants, où $x$ désigne la profondeur du forage exprimée en \emph{dizaines de mètres} (ainsi $x=3$ correspond à \SI{30}{m}) :

\begin{itemize}
\item le \cle{coût total} du forage (terrassement, tubage, équipement), exprimé en \emph{centaines de milliers de FCFA}, est
\[ C(x)=x^{3}-6x^{2}+9x+18 \ ; \]
\item la \cle{recette annuelle} attendue (vente des récoltes irriguées), exprimée dans la même unité, est
\[ R(x)=10x \ .\]
\end{itemize}

Le \cle{bénéfice} de l'opération est $B(x)=R(x)-C(x)$. Le forage est dit \emph{rentable} lorsque $B(x)\geq 0$.

Par ailleurs, la pompe immergée doit remonter l'eau d'une profondeur $h$ (en mètres) qui, d'après les caractéristiques hydrauliques du terrain, vérifie la relation
\[ \sqrt{2h+1}=h-1 \ .\]

Le président de la coopérative, peu à l'aise avec les mathématiques, demande aux élèves de la classe de Première du lycée de Mbouda de répondre à trois questions concrètes :

\begin{enumerate}
\item À quelles profondeurs le forage est-il rentable ?
\item Quelle est la profondeur maximale rentable, arrondie au mètre près ?
\item Quelle profondeur $h$ la pompe doit-elle effectivement remonter ?
\end{enumerate}

\courssub{Consignes}

Travail en groupes de quatre élèves. Chaque groupe rédigera un \emph{rapport écrit} d'une page destiné au président de la coopérative.

\begin{enumerate}
\item \textbf{Mise en place du bénéfice.} Exprimer $B(x)=R(x)-C(x)$ et montrer que
\[ B(x)=-\bigl(x^{3}-6x^{2}-x+18\bigr)\ .\]
On pose $P(x)=x^{3}-6x^{2}-x+18$. Vérifier que $x=2$ est un zéro de $P$.
\item \textbf{Factorisation.} En déduire que $P$ est divisible par $(x-2)$, puis factoriser complètement $P(x)$ en produit de facteurs du premier degré, par la méthode de la division euclidienne \emph{ou} par la méthode des coefficients indéterminés (au choix du groupe ; les deux seront comparées en restitution).
\item \textbf{Seuils de rentabilité.} Résoudre dans $\mathbb{R}$ l'équation du second degré issue du quotient, et donner les valeurs exactes puis approchées (à $10^{-2}$ près) des trois zéros de $P$.
\item \textbf{Étude de signe.} Dresser le tableau de signes de $P(x)$ pour $x\in[0\,;+\infty[$, en déduire celui de $B(x)$, puis résoudre l'inéquation $B(x)\geq 0$. Interpréter le résultat en mètres et en FCFA pour le président.
\item \textbf{Bonus — la pompe.} Résoudre l'équation $\sqrt{2h+1}=h-1$ en élevant au carré, puis \emph{choisir} parmi les solutions obtenues celles qui conviennent. Conclure sur la profondeur $h$.
\item \textbf{Rédaction du rapport.} Rédiger la réponse au président : recommandation chiffrée, seuil maximal de profondeur rentable, et mise en garde éventuelle.
\end{enumerate}

\courssub{Ressources à mobiliser}

\begin{aretenir}[Boîte à outils de la leçon]
\begin{itemize}
\item \textbf{Discriminant} : pour $ax^{2}+bx+c=0$ ($a\neq 0$), $\Delta=b^{2}-4ac$. Si $\Delta>0$, deux racines $x_{1}=\dfrac{-b-\sqrt{\Delta}}{2a}$ et $x_{2}=\dfrac{-b+\sqrt{\Delta}}{2a}$, et $ax^{2}+bx+c=a(x-x_{1})(x-x_{2})$.
\item \textbf{Zéro et factorisation} : $a$ est un zéro du polynôme $P$ si et seulement si $P$ est divisible par $(x-a)$ : $P(x)=(x-a)Q(x)$ avec $\deg Q=\deg P -1$.
\item \textbf{Signe constant (propriété admise)} : si un polynôme ne s'annule pas sur un intervalle, il garde un signe constant sur cet intervalle. Le signe d'un polynôme se lit donc dans un tableau à partir de ses zéros rangés dans l'ordre croissant.
\item \textbf{Équation irrationnelle} $\sqrt{A}=B$ : on résout $A=B^{2}$ \emph{avec la condition} $B\geq 0$ (et $A\geq 0$), puis on vérifie chaque solution dans l'équation de départ.
\end{itemize}
\end{aretenir}

\courssub{Démarche et corrigé commenté}

\begin{exoresolu}[Le forage rentable — corrigé complet]
Énoncé : avec $C(x)=x^{3}-6x^{2}+9x+18$ et $R(x)=10x$, étudier la rentabilité $B(x)\geq 0$ et résoudre $\sqrt{2h+1}=h-1$.
\tcblower
\textbf{Étape 1 : expression du bénéfice.}
\begin{align*}
B(x)&=R(x)-C(x)=10x-\bigl(x^{3}-6x^{2}+9x+18\bigr)\\
&=-x^{3}+6x^{2}+x-18=-\bigl(x^{3}-6x^{2}-x+18\bigr).
\end{align*}
Posons $P(x)=x^{3}-6x^{2}-x+18$. Vérifions que $2$ est un zéro de $P$ :
\[ P(2)=2^{3}-6\times 2^{2}-2+18=8-24-2+18=0\ .\]
Donc $2$ est bien un zéro de $P$ : $P$ est divisible par $(x-2)$.

\textbf{Étape 2 : factorisation de $P$.}

\emph{Méthode 1 : division euclidienne.} On effectue la division de $x^{3}-6x^{2}-x+18$ par $x-2$ :
\begin{itemize}
\item $x^{3}\div x = x^{2}$ ; $x^{2}(x-2)=x^{3}-2x^{2}$ ; reste : $-4x^{2}-x+18$ ;
\item $-4x^{2}\div x = -4x$ ; $-4x(x-2)=-4x^{2}+8x$ ; reste : $-9x+18$ ;
\item $-9x\div x = -9$ ; $-9(x-2)=-9x+18$ ; reste : $0$.
\end{itemize}
Le reste est nul (ce qui confirme que $2$ est un zéro) et le quotient est $x^{2}-4x-9$. Donc
\[ P(x)=(x-2)\bigl(x^{2}-4x-9\bigr)\ .\]

\emph{Méthode 2 : coefficients indéterminés.} On cherche $Q(x)=ax^{2}+bx+c$ tel que $P(x)=(x-2)(ax^{2}+bx+c)$. En développant :
\[ (x-2)(ax^{2}+bx+c)=ax^{3}+(b-2a)x^{2}+(c-2b)x-2c\ .\]
Par identification avec $x^{3}-6x^{2}-x+18$ :
\[ a=1\ ;\quad b-2a=-6\ ;\quad c-2b=-1\ ;\quad -2c=18\ .\]
On obtient $a=1$, $b=-4$, $c=-9$ (les quatre équations sont compatibles), d'où le même résultat : $P(x)=(x-2)(x^{2}-4x-9)$.

\textbf{Étape 3 : résolution de l'équation du second degré.}

Résolvons $x^{2}-4x-9=0$. Le discriminant est
\[ \Delta=(-4)^{2}-4\times 1\times(-9)=16+36=52>0\ .\]
Il y a deux racines réelles :
\[ x_{1}=\frac{4-\sqrt{52}}{2}=\frac{4-2\sqrt{13}}{2}=2-\sqrt{13}\ ;\qquad x_{2}=\frac{4+\sqrt{52}}{2}=2+\sqrt{13}\ .\]
La factorisation complète de $P$ est donc
\[ P(x)=(x-2)\bigl(x-(2-\sqrt{13})\bigr)\bigl(x-(2+\sqrt{13})\bigr)\ .\]
Les trois zéros de $P$, rangés dans l'ordre croissant, sont :
\[ 2-\sqrt{13}\approx -1{,}61\ ;\qquad 2\ ;\qquad 2+\sqrt{13}\approx 5{,}61\ .\]

\textbf{Étape 4 : tableau de signes et rentabilité.}

Le coefficient dominant de $P$ est $1>0$ : $P$ est du signe de $x^{3}$ pour les grandes valeurs, donc positif à droite du plus grand zéro, et les signes alternent en traversant chaque zéro simple. Sur $[0\,;+\infty[$, seuls les zéros $2$ et $2+\sqrt{13}$ interviennent (le zéro $2-\sqrt{13}$ est négatif, hors du domaine d'étude) :

\begin{center}
\begin{tabular}{|c|ccccc|}\hline
$x$ & $0$ & & $2$ & & $2+\sqrt{13}$ \\ \hline
$P(x)$ & & $+$ & $0$ & $-$ & $0$ \\ \hline
$B(x)=-P(x)$ & & $-$ & $0$ & $+$ & $0$ \\ \hline
\end{tabular}
\end{center}

(à droite de $2+\sqrt{13}$, $P$ redevient positif et $B$ négatif : le tableau ci-dessus s'arrête à la dernière valeur remarquable, mais l'étude se poursuit jusqu'en $+\infty$ avec $B(x)<0$ pour $x>2+\sqrt{13}$).

Contrôle par une valeur test : $B(3)=-27+54+3-18=12>0$, conforme au tableau.

L'inéquation $B(x)\geq 0$ a pour ensemble de solutions
\[ \mathcal{S}=\bigl[\,2\ ;\ 2+\sqrt{13}\,\bigr]\ .\]

\textbf{Interprétation.} Le forage est rentable pour une profondeur comprise entre $2$ dizaines de mètres et $(2+\sqrt{13})$ dizaines de mètres, c'est-à-dire entre \SI{20}{m} et environ \SI{56,06}{m}. La profondeur maximale rentable, arrondie au mètre près, est donc \SI{56}{m}. Le bénéfice maximal est atteint aux environs de $x\approx 3{,}9$ (environ \SI{39}{m}), où $B$ vaut environ $12{,}3$ centaines de milliers de FCFA, soit environ \num{1230000} FCFA.

\textbf{Étape 5 : la pompe (équation irrationnelle).}

L'équation $\sqrt{2h+1}=h-1$ n'a de sens que si $2h+1\geq 0$, et une racine carrée étant positive, il faut aussi $h-1\geq 0$, c'est-à-dire $h\geq 1$. Sous cette condition, on élève au carré :
\begin{align*}
2h+1=(h-1)^{2}&\iff 2h+1=h^{2}-2h+1\\
&\iff h^{2}-4h=0\\
&\iff h(h-4)=0\iff h=0\ \text{ou}\ h=4\ .
\end{align*}
\emph{Choix des solutions} : $h=0$ ne vérifie pas la condition $h\geq 1$ (et effectivement $\sqrt{1}=1\neq -1$) : c'est une solution parasite introduite par l'élévation au carré, on la rejette. Pour $h=4$ : $\sqrt{2\times 4+1}=\sqrt{9}=3$ et $4-1=3$ : l'égalité est vérifiée. La pompe doit donc remonter l'eau d'une profondeur $h=\SI{4}{m}$.

\textbf{Étape 6 : éléments de rédaction du rapport.}

Le rapport au président doit contenir : la plage de rentabilité $[20\ \text{m}\ ;\ 56\ \text{m}]$, la profondeur maximale rentable (\SI{56}{m}), la recommandation de viser environ \SI{39}{m} pour un bénéfice maximal voisin de \num{1230000} FCFA, et la mise en garde : au-delà de \SI{56}{m}, le coût du forage dépasse la recette attendue et l'opération devient déficitaire.
\end{exoresolu}

\begin{attention}[Pièges à éviter dans cette activité]
\begin{itemize}
\item \textbf{Ne pas oublier de vérifier le zéro annoncé} : avant toute factorisation, calculer $P(2)$. Si le résultat n'est pas nul, toute la suite serait fausse.
\item \textbf{Élévation au carré non réversible} : résoudre $A=B^{2}$ donne les solutions de $\sqrt{A}=B$ \emph{et} celles de $\sqrt{A}=-B$. La condition $B\geq 0$ (ici $h\geq 1$) et la vérification finale sont obligatoires ; $h=0$ est une solution parasite à rejeter explicitement.
\item \textbf{Changement de signe} : $B(x)=-P(x)$, donc $B\geq 0$ équivaut à $P\leq 0$. Inverser les signes dans le tableau est l'erreur la plus fréquente.
\item \textbf{Unités} : $x$ est en dizaines de mètres et $B$ en centaines de milliers de FCFA ; la réponse au président doit être convertie en mètres et en FCFA.
\end{itemize}
\end{attention}

\begin{exercice}[Pour aller plus loin]
Le bureau d'études propose une seconde coopérative, à Bafang, avec $C(x)=x^{3}-6x^{2}+11x-6$ et $R(x)=0$ (étude du coût seul). Montrer que $2$ est un zéro de $C$, factoriser complètement $C(x)$, et déterminer pour quelles profondeurs le coût est inférieur ou égal à $6$ centaines de milliers de FCFA.
\end{exercice}

\begin{corrige}[Exercice « Pour aller plus loin »]
$C(2)=8-24+22-6=0$, donc $C$ est divisible par $(x-2)$. Par division euclidienne ou coefficients indéterminés, $C(x)=(x-2)(x^{2}-4x+3)$. Le trinôme $x^{2}-4x+3$ a pour discriminant $\Delta=16-12=4$, racines $1$ et $3$. Donc $C(x)=(x-1)(x-2)(x-3)$. L'inéquation $C(x)\leq 6$ s'écrit $x^{3}-6x^{2}+11x-12\leq 0$ ; on vérifie que $4$ est un zéro de $Q(x)=x^{3}-6x^{2}+11x-12$ ($64-96+44-12=0$), d'où $Q(x)=(x-4)(x^{2}-2x+3)$. Le trinôme $x^{2}-2x+3$ a un discriminant $4-12=-8<0$, il est donc strictement positif partout. Le signe de $Q$ est celui de $(x-4)$ : $Q(x)\leq 0\iff x\leq 4$. Le coût reste sous le seuil pour une profondeur inférieure ou égale à \SI{40}{m}.
\end{corrige}

\newpage

\coursec{Méthodes et exercices résolus}

\courssub{Fiche 1 — Résoudre une équation du second degré}

\begin{methode}[Résoudre $ax^2+bx+c=0$ ($a\neq 0$) : forme canonique puis discriminant]
\begin{enumerate}
\item \textbf{Vérifier} que $a\neq 0$ et identifier les coefficients $a$, $b$, $c$.
\item \textbf{Mise sous forme canonique} (utile quand on demande la factorisation ou les variations) :
\[ ax^2+bx+c = a\left[\left(x+\frac{b}{2a}\right)^2 - \frac{b^2-4ac}{4a^2}\right] = a\left[\left(x+\frac{b}{2a}\right)^2 - \frac{\Delta}{4a^2}\right]. \]
\item \textbf{Calculer le discriminant} $\Delta = b^2-4ac$.
\item \textbf{Conclure selon le signe de $\Delta$} :
\begin{itemize}
\item si $\Delta > 0$ : deux solutions réelles distinctes $x_1=\dfrac{-b-\sqrt{\Delta}}{2a}$ et $x_2=\dfrac{-b+\sqrt{\Delta}}{2a}$ ;
\item si $\Delta = 0$ : une solution double $x_0=-\dfrac{b}{2a}$ ;
\item si $\Delta < 0$ : aucune solution réelle, $S=\emptyset$.
\end{itemize}
\item \textbf{Vérifier} mentalement avec la somme $x_1+x_2=-\dfrac{b}{a}$ et le produit $x_1x_2=\dfrac{c}{a}$.
\end{enumerate}
\textbf{Réflexes} : si $b$ est pair, utiliser le discriminant réduit $\Delta'=\left(\dfrac{b}{2}\right)^2-ac$, les solutions étant alors $x=\dfrac{-\frac{b}{2}\pm\sqrt{\Delta'}}{a}$ ; si $a+b+c=0$, alors $x_1=1$ et $x_2=\dfrac{c}{a}$ ; si $a-b+c=0$, alors $x_1=-1$ et $x_2=-\dfrac{c}{a}$.
\end{methode}

\begin{exoresolu}[Équation du second degré — type Probatoire]
Résoudre dans $\mathbb{R}$ l'équation $(E)$ : $2x^2-5x-3=0$.
\tcblower
\textbf{Étape 1 : identification.} Ici $a=2$, $b=-5$, $c=-3$, avec $a=2\neq 0$ : c'est bien une équation du second degré.

\textbf{Étape 2 : discriminant.}
\[ \Delta = b^2-4ac = (-5)^2 - 4\times 2 \times (-3) = 25+24 = 49. \]

\textbf{Étape 3 : discussion.} $\Delta = 49 > 0$, donc $(E)$ admet deux solutions réelles distinctes :
\[ x_1 = \frac{-b-\sqrt{\Delta}}{2a} = \frac{5-\sqrt{49}}{4} = \frac{5-7}{4} = \frac{-2}{4} = -\frac{1}{2}, \qquad x_2 = \frac{5+7}{4} = \frac{12}{4} = 3. \]

\textbf{Étape 4 : vérification.} Somme : $-\dfrac{1}{2}+3 = \dfrac{5}{2} = -\dfrac{b}{a}$. Produit : $-\dfrac{1}{2}\times 3 = -\dfrac{3}{2} = \dfrac{c}{a}$. Les deux contrôles sont conformes.

\textbf{Conclusion :} $S = \left\{-\dfrac{1}{2}\,;\,3\right\}$. Par conséquent, $2x^2-5x-3 = 2\left(x+\dfrac{1}{2}\right)(x-3) = (2x+1)(x-3)$.
\end{exoresolu}

\courssub{Fiche 2 — Équation du second degré avec paramètre}

\begin{methode}[Résoudre $ax^2+bx+c=0$ où un coefficient dépend d'un paramètre $m$]
\begin{enumerate}
\item \textbf{Repérer si le coefficient de $x^2$ dépend de $m$} : si oui, traiter à part les valeurs de $m$ qui l'annulent (l'équation devient alors du premier degré et se résout directement).
\item \textbf{Calculer $\Delta(m)$}, qui est lui-même un polynôme en $m$.
\item \textbf{Étudier le signe de $\Delta(m)$} : factoriser $\Delta(m)$, trouver ses racines, dresser un tableau de signes en $m$.
\item \textbf{Conclure par cas} : pour chaque intervalle de valeurs de $m$, préciser le nombre de solutions et, si on le demande, les exprimer en fonction de $m$.
\item \textbf{Cas particuliers fréquents} : « l'équation admet une racine double » $\Leftrightarrow \Delta(m)=0$ ; « deux racines distinctes » $\Leftrightarrow \Delta(m)>0$ (avec $a\neq 0$).
\end{enumerate}
\end{methode}

\begin{exoresolu}[Équation paramétrique — type Probatoire]
Soit $m$ un réel. On considère l'équation $(E_m)$ : $x^2-2(m-1)x+m=0$.
Déterminer, suivant les valeurs de $m$, le nombre de solutions de $(E_m)$, puis résoudre $(E_2)$.
\tcblower
\textbf{Étape 1.} Le coefficient de $x^2$ est $1$, indépendant de $m$ et non nul : $(E_m)$ est toujours une équation du second degré, quel que soit $m$.

\textbf{Étape 2 : discriminant.} Avec $a=1$, $b=-2(m-1)$, $c=m$, on utilise le discriminant réduit (car $b$ est pair) :
\[ \Delta' = (m-1)^2 - m = m^2 - 2m + 1 - m = m^2 - 3m + 1. \]

\textbf{Étape 3 : signe de $\Delta'$.} Résolvons $m^2-3m+1=0$ : son discriminant est $\delta = (-3)^2-4\times 1\times 1 = 9-4 = 5 > 0$, d'où deux racines
\[ m_1 = \frac{3-\sqrt{5}}{2}, \qquad m_2 = \frac{3+\sqrt{5}}{2}. \]
$\Delta'$ est un trinôme en $m$ de coefficient dominant positif : il est \textbf{positif à l'extérieur des racines} et \textbf{négatif entre les racines}.

\textbf{Étape 4 : conclusion par cas.}
\begin{itemize}
\item Si $m < \dfrac{3-\sqrt{5}}{2}$ ou $m > \dfrac{3+\sqrt{5}}{2}$ : $\Delta'>0$, $(E_m)$ a \textbf{deux solutions distinctes} $x = (m-1) \pm \sqrt{m^2-3m+1}$.
\item Si $m = \dfrac{3-\sqrt{5}}{2}$ ou $m = \dfrac{3+\sqrt{5}}{2}$ : $\Delta'=0$, $(E_m)$ a \textbf{une solution double} $x_0 = m-1$.
\item Si $\dfrac{3-\sqrt{5}}{2} < m < \dfrac{3+\sqrt{5}}{2}$ : $\Delta'<0$, $(E_m)$ \textbf{n'a pas de solution réelle}.
\end{itemize}

\textbf{Résolution de $(E_2)$.} Pour $m=2$, l'équation s'écrit $x^2 - 2x + 2 = 0$. Deux façons de conclure :
\begin{itemize}
\item \emph{Par la discussion précédente} : $\dfrac{3+\sqrt{5}}{2} \approx 2{,}62$, donc $m=2$ appartient à l'intervalle $\left]\dfrac{3-\sqrt{5}}{2}\,;\,\dfrac{3+\sqrt{5}}{2}\right[$ : on est dans le troisième cas, pas de solution réelle.
\item \emph{Par vérification directe} : $\Delta' = 2^2 - 3\times 2 + 1 = 4-6+1 = -1 < 0$.
\end{itemize}
\textbf{Conclusion :} $S = \emptyset$.
\end{exoresolu}

\courssub{Fiche 3 — Signe du trinôme et inéquations du second degré}

\begin{methode}[Dresser le tableau de signes de $ax^2+bx+c$ et résoudre une inéquation]
\begin{enumerate}
\item \textbf{Calculer $\Delta$} et déterminer les racines éventuelles $x_1 \leq x_2$.
\item \textbf{Appliquer la règle du signe du trinôme} :
\begin{itemize}
\item $\Delta > 0$ : le trinôme est du \textbf{signe de $a$ à l'extérieur des racines}, du signe contraire entre les racines ;
\item $\Delta = 0$ : signe de $a$ partout, nul en $x_0$ ;
\item $\Delta < 0$ : signe de $a$ pour tout réel $x$.
\end{itemize}
\item \textbf{Dresser le tableau de signes} en ordonnant les racines sur la première ligne.
\item \textbf{Lire la solution de l'inéquation} : pour $\geq 0$ ou $\leq 0$, \textbf{fermer les crochets} aux racines ; pour $>$ ou $<$, les laisser \textbf{ouverts}.
\end{enumerate}
\end{methode}

\begin{exoresolu}[Inéquation du second degré — type Probatoire]
Un transporteur estime que la recette journalière (en milliers de FCFA) d'une ligne, en fonction du nombre $x$ de rotations, est modélisée par $R(x) = -x^2 + 8x - 7$. Pour quels nombres de rotations la recette est-elle strictement supérieure à \num{5000} FCFA ? Autrement dit, résoudre dans $\mathbb{R}$ : $-x^2+8x-7 > 5$.
\tcblower
\textbf{Étape 1 : se ramener à un trinôme comparé à 0.}
\[ -x^2 + 8x - 7 > 5 \iff -x^2 + 8x - 12 > 0. \]

\textbf{Étape 2 : discriminant} de $-x^2+8x-12$ ($a=-1$, $b=8$, $c=-12$) :
\[ \Delta = 64 - 4\times(-1)\times(-12) = 64 - 48 = 16 > 0. \]
Racines : $x = \dfrac{-8-4}{-2} = 6$ et $x = \dfrac{-8+4}{-2} = 2$. On ordonne : $x_1 = 2$, $x_2 = 6$.

\textbf{Étape 3 : tableau de signes.} Comme $a = -1 < 0$, le trinôme est négatif à l'extérieur des racines, positif entre :
\begin{center}
\begin{tabular}{|c|ccccc|}\hline
$x$ & $-\infty$ & & $2$ & & $+\infty$ \\ \hline
$-x^2+8x-12$ & & $-$ & $0$ & $+$ & \\ \hline
\end{tabular}
\end{center}

\textbf{Étape 4 : lecture.} L'inéquation est stricte ($>$), donc les racines sont exclues :
\[ S = \left]2\,;\,6\right[. \]
\textbf{Conclusion concrète :} la recette dépasse \num{5000} FCFA pour un nombre de rotations strictement compris entre 2 et 6, c'est-à-dire pour 3, 4 ou 5 rotations.
\end{exoresolu}

\courssub{Fiche 4 — Systèmes se ramenant à une équation du second degré}

\begin{methode}[Résoudre un système du type $\begin{cases} x+y = S \\ xy = P \end{cases}$ ou s'y ramenant]
\begin{enumerate}
\item \textbf{Identifier la somme $S$ et le produit $P$} des deux inconnues (parfois après une transformation : poser $X=x^2$, $Y=y^2$, ou utiliser $x^2+y^2 = S^2-2P$).
\item \textbf{Utiliser le théorème} : $x$ et $y$ sont les solutions de l'équation du second degré
\[ t^2 - St + P = 0. \]
\item \textbf{Condition d'existence} : des solutions réelles existent si et seulement si $S^2 - 4P \geq 0$.
\item \textbf{Résoudre} cette équation en $t$, puis \textbf{former les couples} $(x\,;\,y)$ : attention, si $t_1 \neq t_2$, il y a \textbf{deux couples symétriques} $(t_1\,;\,t_2)$ et $(t_2\,;\,t_1)$.
\item \textbf{Vérifier} chaque couple dans le système de départ.
\end{enumerate}
\end{methode}

\begin{exoresolu}[Système somme-produit — type Probatoire]
Un champ rectangulaire de la région de l'Ouest a un périmètre de $46$ m et une aire de $120$ m$^2$. Déterminer ses dimensions $x$ et $y$ (en mètres), c'est-à-dire résoudre :
\[ \begin{cases} x + y = 23 \\ xy = 120 \end{cases} \]
\tcblower
\textbf{Étape 1.} Le demi-périmètre donne $x+y = 23$ et l'aire donne $xy = 120$. On a donc $S = 23$ et $P = 120$.

\textbf{Étape 2 : condition d'existence.} $S^2 - 4P = 529 - 480 = 49 \geq 0$ : il existe des solutions réelles.

\textbf{Étape 3 : équation auxiliaire.} $x$ et $y$ sont les racines de
\[ t^2 - 23t + 120 = 0. \]
$\Delta = 23^2 - 4\times 120 = 529 - 480 = 49 > 0$, d'où
\[ t_1 = \frac{23 - 7}{2} = 8, \qquad t_2 = \frac{23 + 7}{2} = 15. \]

\textbf{Étape 4 : couples solutions.} Les inconnues jouant des rôles symétriques, on obtient deux couples :
\[ (x\,;\,y) = (8\,;\,15) \quad \text{ou} \quad (x\,;\,y) = (15\,;\,8). \]

\textbf{Étape 5 : vérification.} $8+15 = 23$ et $8 \times 15 = 120$ : conforme.

\textbf{Conclusion :} le champ mesure $15$ m de long sur $8$ m de large ; $S = \{(8\,;\,15)\,;\,(15\,;\,8)\}$.
\end{exoresolu}

\courssub{Fiche 5 — Zéros d'un polynôme de degré 3 et factorisation}

\begin{methode}[Factoriser $P(x) = ax^3+bx^2+cx+d$ connaissant un zéro $\alpha$]
\begin{enumerate}
\item \textbf{Chercher un zéro évident} parmi $\pm 1$, $\pm 2$, $\pm 3$, les diviseurs de $d$ : calculer $P(1)$, $P(-1)$, $P(2)$\dots{} Vérifier que $P(\alpha) = 0$ si $\alpha$ est donné.
\item \textbf{Écrire la factorisation annoncée} : $P(x) = (x-\alpha)(ax^2 + b'x + c')$.
\item \textbf{Déterminer $b'$ et $c'$} par l'une des deux méthodes :
\begin{itemize}
\item \emph{Coefficients indéterminés} : développer $(x-\alpha)(ax^2+b'x+c')$, identifier les coefficients avec ceux de $P$, résoudre le petit système obtenu ;
\item \emph{Division euclidienne} : diviser $P(x)$ par $x-\alpha$ (reste nul puisque $\alpha$ est zéro).
\end{itemize}
\item \textbf{Factoriser le quotient} $ax^2+b'x+c'$ (discriminant) s'il admet des racines.
\item \textbf{Conclure} par la factorisation complète : $P(x) = a(x-\alpha)(x-\beta)(x-\gamma)$.
\end{enumerate}
\textbf{Réflexe} : le coefficient de $x^2$ dans le quotient est toujours $a$ (même coefficient dominant), et le terme constant vérifie $-\alpha \times c' = d$.
\end{methode}

\begin{exoresolu}[Factorisation d'un polynôme de degré 3 — type Probatoire]
Soit $P(x) = 2x^3 - 3x^2 - 3x + 2$.
\begin{enumerate}
\item Vérifier que $-1$ est un zéro de $P$.
\item Factoriser $P(x)$ en produit de trois facteurs du premier degré.
\item Résoudre dans $\mathbb{R}$ l'équation $P(x) = 0$.
\end{enumerate}
\tcblower
\textbf{1. Vérification.} $P(-1) = 2(-1)^3 - 3(-1)^2 - 3(-1) + 2 = -2 - 3 + 3 + 2 = 0$. Donc $-1$ est bien un zéro de $P$.

\textbf{2. Factorisation par coefficients indéterminés.} Puisque $P(-1)=0$, $P$ est divisible par $(x+1)$ : il existe $b'$ et $c'$ tels que
\[ P(x) = (x+1)(2x^2 + b'x + c'). \]
Développons : $(x+1)(2x^2 + b'x + c') = 2x^3 + (b'+2)x^2 + (c'+b')x + c'$.

Identification avec $2x^3 - 3x^2 - 3x + 2$ :
\[ \begin{cases} b' + 2 = -3 \\ c' + b' = -3 \\ c' = 2 \end{cases} \iff \begin{cases} b' = -5 \\ c' = 2 \end{cases} \]
(la deuxième équation donne $2 - 5 = -3$, ce qui est bien cohérent). Donc $P(x) = (x+1)(2x^2 - 5x + 2)$.

Factorisons $2x^2 - 5x + 2$ : $\Delta = 25 - 16 = 9 > 0$, racines $x = \dfrac{5-3}{4} = \dfrac{1}{2}$ et $x = \dfrac{5+3}{4} = 2$. Donc $2x^2-5x+2 = 2\left(x-\dfrac{1}{2}\right)(x-2)$.

\textbf{Factorisation complète :} $P(x) = 2(x+1)\left(x-\dfrac{1}{2}\right)(x-2)$.

\textbf{3. Résolution de $P(x)=0$.} Un produit de facteurs est nul si et seulement si l'un des facteurs est nul :
\[ S = \left\{-1\,;\,\frac{1}{2}\,;\,2\right\}. \]
\end{exoresolu}

\courssub{Fiche 6 — Signe d'un polynôme de degré 3 et inéquations}

\begin{methode}[Dresser le tableau de signes de $P(x)$ de degré 3 et résoudre $P(x) \geq 0$]
\begin{enumerate}
\item \textbf{Factoriser complètement} $P(x)$ (fiche 5) : $P(x) = a(x-\alpha)(x-\beta)(x-\gamma)$ avec $\alpha < \beta < \gamma$.
\item \textbf{Construire le tableau de signes} : une ligne par facteur $(x-\alpha)$, $(x-\beta)$, $(x-\gamma)$ (chacun s'annule en sa racine, négatif avant, positif après), puis une ligne produit obtenue par la \textbf{règle des signes}.
\item \textbf{Contrôle rapide} : pour $x$ très grand, $P(x)$ a le signe de $a$ (terme dominant $ax^3$) ; le signe \textbf{change à chaque racine simple}.
\item \textbf{Lire les solutions} de l'inéquation sur la dernière ligne, en fermant les crochets si l'inégalité est large.
\end{enumerate}
\textbf{Propriété admise à citer} : un polynôme qui ne s'annule pas sur un intervalle y garde un signe constant — c'est ce qui justifie le tableau.
\end{methode}

\begin{exoresolu}[Inéquation de degré 3 — type Probatoire]
Résoudre dans $\mathbb{R}$ l'inéquation $2x^3 - 3x^2 - 3x + 2 \leq 0$ (on utilisera la factorisation de la fiche 5).
\tcblower
\textbf{Étape 1 : factorisation.} D'après l'exercice précédent :
\[ P(x) = 2(x+1)\left(x-\frac{1}{2}\right)(x-2). \]
Les zéros, rangés par ordre croissant, sont : $-1$, $\dfrac{1}{2}$, $2$.

\textbf{Étape 2 : tableau de signes.} Le facteur constant $2$ est strictement positif, il n'influence pas le signe.
\begin{center}
\begin{tabular}{|c|ccccccccc|}\hline
$x$ & $-\infty$ & & $-1$ & & $\frac{1}{2}$ & & $2$ & & $+\infty$ \\ \hline
$x+1$ & & $-$ & $0$ & $+$ & $+$ & $+$ & $+$ & $+$ & \\ \hline
$x-\frac{1}{2}$ & & $-$ & $-$ & $-$ & $0$ & $+$ & $+$ & $+$ & \\ \hline
$x-2$ & & $-$ & $-$ & $-$ & $-$ & $-$ & $0$ & $+$ & \\ \hline
$P(x)$ & & $-$ & $0$ & $+$ & $0$ & $-$ & $0$ & $+$ & \\ \hline
\end{tabular}
\end{center}
Détail de la ligne produit : avant $-1$, trois facteurs négatifs donnent un produit négatif ; entre $-1$ et $\tfrac12$, un seul facteur positif parmi trois : produit positif ; entre $\tfrac12$ et $2$, deux positifs et un négatif : produit négatif ; après $2$, tout est positif. Contrôle : pour $x$ grand, $P(x)$ a le signe de $2x^3$, donc est positif : cohérent avec la dernière colonne.

\textbf{Étape 3 : lecture.} On cherche $P(x) \leq 0$ (inégalité large : on ferme les crochets aux zéros) :
\[ S = \left]-\infty\,;\,-1\right] \cup \left[\frac{1}{2}\,;\,2\right]. \]
\end{exoresolu}

\courssub{Fiche 7 — Équations et inéquations irrationnelles}

\begin{methode}[Résoudre $\sqrt{ax+b} = cx+d$]
\begin{enumerate}
\item \textbf{Poser les conditions} : une racine carrée exige $ax+b \geq 0$ ; de plus, une racine carrée étant \textbf{positive}, l'égalité impose $cx+d \geq 0$. (On peut aussi reporter cette vérification à la fin.)
\item \textbf{Élever au carré} : $ax+b = (cx+d)^2$, ce qui donne une équation du second degré.
\item \textbf{Résoudre} cette équation (2).
\item \textbf{Filtrer} : ne retenir que les solutions de (2) qui vérifient la condition $cx+d \geq 0$ (ou, plus sûr, qui vérifient l'équation (1) de départ). \textbf{Toute solution non vérifiée est une solution parasite à rejeter.}
\item \textbf{Pour une inéquation} $\sqrt{ax+b} \leq cx+d$ : résoudre le système $\begin{cases} ax+b \geq 0 \\ cx+d \geq 0 \\ ax+b \leq (cx+d)^2 \end{cases}$ ; pour $\sqrt{ax+b} \geq cx+d$ : réunion des deux cas $\begin{cases} cx+d < 0 \\ ax+b \geq 0 \end{cases}$ ou $\begin{cases} cx+d \geq 0 \\ ax+b \geq (cx+d)^2 \end{cases}$.
\end{enumerate}
\end{methode}

\begin{exoresolu}[Équation irrationnelle — type Probatoire]
Résoudre dans $\mathbb{R}$ l'équation $(1)$ : $\sqrt{3x+1} = x - 1$.
\tcblower
\textbf{Étape 1 : conditions.}
\begin{itemize}
\item Existence de la racine : $3x+1 \geq 0 \iff x \geq -\dfrac{1}{3}$.
\item Le membre de gauche est positif ou nul, donc il faut $x - 1 \geq 0 \iff x \geq 1$.
\end{itemize}
La condition la plus forte est $x \geq 1$ : toute solution devra appartenir à $[1\,;\,+\infty[$.

\textbf{Étape 2 : élévation au carré.} Sous ces conditions, les deux membres sont positifs, donc
\[ (1) \implies 3x+1 = (x-1)^2 \iff 3x+1 = x^2 - 2x + 1 \iff x^2 - 5x = 0 \iff x(x-5) = 0. \]
L'équation (2) admet pour solutions $x = 0$ et $x = 5$.

\textbf{Étape 3 : filtrage.} La condition $x \geq 1$ élimine $x = 0$. Vérifions $x = 5$ dans $(1)$ : $\sqrt{3\times 5 + 1} = \sqrt{16} = 4$ et $5 - 1 = 4$. L'égalité est vraie.

Vérification du rejet de $0$ : $\sqrt{3\times 0 + 1} = 1$ mais $0 - 1 = -1$ : $1 \neq -1$, donc $0$ est bien une \textbf{solution parasite} introduite par l'élévation au carré.

\textbf{Conclusion :} $S = \{5\}$.
\end{exoresolu}

\begin{attention}[Les 5 erreurs qui coûtent des points au Probatoire]
\begin{enumerate}
\item \textbf{Oublier la condition $cx+d \geq 0$ dans $\sqrt{ax+b}=cx+d$.} Élever au carré n'est pas une équivalence : il faut filtrer les solutions. Écrire « $S=\{0\,;\,5\}$ » dans l'exercice de la fiche 7 fait perdre la moitié des points.
\item \textbf{Se tromper dans la règle des signes du trinôme.} C'est le signe de $a$ \emph{à l'extérieur} des racines, le signe \emph{opposé} entre les racines. Beaucoup d'élèves l'inversent quand $a<0$ : réflexe de survie — multiplier par $-1$ en changeant le sens de l'inégalité pour se ramener à $a>0$.
\item \textbf{Oublier les crochets fermés pour les inégalités larges.} Pour $\leq$ ou $\geq$, les racines sont solutions : $S = [x_1\,;\,x_2]$, pas $]x_1\,;\,x_2[$. La rédaction du tableau de signes doit montrer les zéros.
\item \textbf{Factoriser un polynôme de degré 3 sans vérifier le zéro annoncé.} Avant d'écrire $P(x)=(x-\alpha)Q(x)$, il faut avoir \emph{calculé} $P(\alpha)$ et constaté qu'il vaut $0$ : c'est la justification exigée. De même, ne pas oublier le coefficient dominant $a$ devant la factorisation : $2x^2-5x+2 = 2\left(x-\tfrac12\right)(x-2)$ et non $\left(x-\tfrac12\right)(x-2)$.
\item \textbf{Donner un seul couple pour un système symétrique.} Si $x$ et $y$ jouent des rôles symétriques et que $t_1 \neq t_2$, il y a \emph{deux} couples solutions $(t_1\,;\,t_2)$ et $(t_2\,;\,t_1)$. Et dans une équation paramétrique, ne jamais oublier de traiter à part les valeurs du paramètre qui annulent le coefficient de $x^2$.
\end{enumerate}
\end{attention}

\newpage

\coursec{Exercices}

\courssub{Je vérifie mes connaissances}

\begin{exercice}[QCM : Discriminant et nombre de solutions]
Parmi les affirmations suivantes, une seule est exacte. Laquelle ? Justifier brièvement.
\begin{enumerate}
    \item L'équation $2x^2 - 3x + 4 = 0$ admet deux solutions réelles distinctes.
    \item L'équation $x^2 - 4x + 4 = 0$ n'admet aucune solution réelle.
    \item L'équation $-x^2 + 2x - 5 = 0$ admet une unique solution réelle.
    \item L'équation $3x^2 + 5x - 2 = 0$ admet deux solutions réelles distinctes.
\end{enumerate}
\end{exercice}

\begin{exercice}[Vrai ou Faux : Signe d'un trinôme]
Dire si les affirmations suivantes sont vraies ou fausses. Justifier chaque réponse par un calcul ou une propriété.
\begin{enumerate}
    \item Le trinôme $f(x) = -2x^2 + 4x - 3$ est strictement négatif sur $\mathbb{R}$.
    \item L'ensemble des solutions de l'inéquation $x^2 - 6x + 9 \le 0$ est $\mathbb{R}$.
    \item Si $\Delta < 0$ et $a > 0$, alors $ax^2+bx+c > 0$ pour tout $x \in \mathbb{R}$.
    \item L'inéquation $x^2 + 1 < 0$ admet au moins une solution dans $\mathbb{R}$.
\end{enumerate}
\end{exercice}

\begin{exercice}[Calcul direct : Forme canonique et racines]
Pour chacun des trinômes suivants, donner sa forme canonique, son discriminant $\Delta$, et ses racines réelles s'il en existe.
\begin{enumerate}
    \item $P_1(x) = x^2 - 5x + 6$
    \item $P_2(x) = 2x^2 + 4x + 2$
    \item $P_3(x) = -x^2 + 3x - 5$
\end{enumerate}
\end{exercice}

\begin{exercice}[Définitions et propriété : Polynôme de degré 3]
Soit $P(x) = 2x^3 - 5x^2 - x + 6$.
\begin{enumerate}
    \item Vérifier que $2$ est un zéro de $P$.
    \item En déduire un facteur de $P(x)$.
    \item Écrire $P(x)$ sous la forme $(x-2)Q(x)$ où $Q$ est un polynôme de degré 2 (on ne demande pas de développer $Q$).
    \item Rappeler la propriété admise liant le signe d'un polynôme et ses zéros sur un intervalle.
\end{enumerate}
\end{exercice}

\courssub{Je m'entraîne}

\begin{exercice}[Équations du second degré à paramètre]
Soit l'équation $(E_m) : (m-1)x^2 + 2mx + (m+2) = 0$, où $m \in \mathbb{R}$.
\begin{enumerate}
    \item Résoudre $(E_m)$ pour $m = 1$.
    \item Pour $m \neq 1$, calculer le discriminant $\Delta$ de $(E_m)$ en fonction de $m$.
    \item Déterminer les valeurs de $m$ pour lesquelles $(E_m)$ admet :
    \begin{enumerate}
        \item une racine double ;
        \item deux racines réelles distinctes ;
        \item aucune solution réelle.
    \end{enumerate}
    \item Pour $m = -2$, résoudre $(E_{-2})$.
\end{enumerate}
\end{exercice}

\begin{exercice}[Inéquations du second degré et tableaux de signes]
Résoudre dans $\mathbb{R}$ les inéquations suivantes. On présentera le tableau de signes du trinôme associé.
\begin{enumerate}
    \item $-2x^2 + 5x + 3 \ge 0$
    \item $x^2 - 6x + 9 < 0$
    \item $3x^2 + 2x - 1 > 0$
\end{enumerate}
\end{exercice}

\begin{exercice}[Problème géométrique : Rectangle]
Un agriculteur souhaite clôturer un champ rectangulaire d'aire $28\,\text{m}^2$ avec $22\,\text{m}$ de grillage (périmètre).
\begin{enumerate}
    \item Soit $x$ la longueur (en mètres) et $y$ la largeur. Exprimer $y$ en fonction de $x$ à partir du périmètre.
    \item En déduire une équation du second degré vérifiée par $x$.
    \item Résoudre cette équation et donner les dimensions du rectangle.
    \item Y a-t-il une autre solution possible si on échange longueur et largeur ?
\end{enumerate}
\end{exercice}

\begin{exercice}[Polynôme de degré 3 : Factorisation, équation et inéquation]
Soit le polynôme $P(x) = 2x^3 - 3x^2 - 3x + 2$.
\begin{enumerate}
    \item Vérifier que $-1$ est un zéro de $P$.
    \item Factoriser $P(x)$ en utilisant la division euclidienne par $x+1$ (ou la méthode des coefficients indéterminés).
    \item Résoudre l'équation $P(x) = 0$ dans $\mathbb{R}$.
    \item Dresser le tableau de signes de $P$ sur $\mathbb{R}$.
    \item En déduire l'ensemble des solutions de l'inéquation $P(x) > 0$.
\end{enumerate}
\end{exercice}

\courssub{Je me prépare au Probatoire}

\begin{exercice}[Bénéfice d'un commerçant — Synthèse 2nd degré]
Un commerçant de la place du Marché Central à Yaoundé vend des sacs de riz. Son bénéfice net (en dizaines de milliers de FCFA) en fonction du nombre $x$ de sacs vendus (en centaines) est modélisé par la fonction :
\[ B(x) = -x^2 + 10x - 16 \quad \text{avec } x \in [0\,;\, 10]. \]
\begin{enumerate}
    \item Déterminer les nombres de sacs vendus pour lesquels le bénéfice est nul (seuil de rentabilité).
    \item Dresser le tableau de signes de $B(x)$ sur $[0\,;\, 10]$.
    \item Pour quelles quantités vendues le bénéfice est-il \textbf{strictement positif} ?
    \item Déterminer le nombre de sacs à vendre pour obtenir le bénéfice maximal et calculer ce bénéfice maximal (en FCFA).
    \item Le commerçant doit payer une taxe fixe de $50\,000$ FCFA. Écrire la nouvelle fonction de bénéfice $B_1(x)$ et déterminer si le bénéfice maximal permet de couvrir cette taxe.
\end{enumerate}
\end{exercice}

\begin{exercice}[Polynôme de degré 4 et équations irrationnelles]
\emph{Partie A : Étude d'un polynôme de degré 4}
Soit la fonction $f$ définie sur $\mathbb{R}$ par $f(x) = (x^2 - 4)(x^2 - 2x + 5)$.
\begin{enumerate}
    \item Factoriser $x^2 - 4$. Déterminer le discriminant de $x^2 - 2x + 5$ et en déduire son signe sur $\mathbb{R}$.
    \item Dresser le tableau de signes de $f$ sur $\mathbb{R}$.
    \item Résoudre l'inéquation $f(x) \le 0$.
\end{enumerate}

\emph{Partie B : Équation irrationnelle}
Résoudre dans $\mathbb{R}$ l'équation :
\[ \sqrt{3x - 2} = x - 2 \]
\begin{enumerate}
    \item Préciser l'ensemble de définition $D$ de l'équation.
    \item Élever les deux membres au carré et résoudre l'équation du second degré obtenue.
    \item Vérifier lesquelles des solutions obtenues appartiennent à $D$ et vérifient l'équation initiale.
    \item Conclure.
\end{enumerate}
\end{exercice}

\begin{exercice}[Inéquation irrationnelle et problème de dimensionnement]
\emph{Partie A : Inéquation irrationnelle}
Résoudre dans $\mathbb{R}$ l'inéquation :
\[ \sqrt{x + 6} \ge x \]
\begin{enumerate}
    \item Déterminer l'ensemble de définition $D$ de l'inéquation.
    \item Étudier le signe du second membre $x$ sur $D$ pour distinguer les cas.
    \item Résoudre l'inéquation en distinguant les cas $x < 0$ et $x \ge 0$ (pour $x \ge 0$, on peut élever au carré).
    \item Donner l'ensemble final des solutions.
\end{enumerate}

\emph{Partie B : Application — Câblage électrique}
Pour l'électrification d'un village, un poteau CAMTEL de hauteur $h$ mètres est haubané par un câble tendu fixé au sol à $3$ mètres du pied du poteau. La longueur du câble est $\sqrt{h^2 + 9}$ mètres.
On souhaite que la longueur du câble soit \textbf{inférieure ou égale} à la hauteur du poteau augmentée de $1$ mètre.
\begin{enumerate}
    \item Traduire cette condition par une inéquation en $h$.
    \item Montrer que cette inéquation se ramène à $\sqrt{h^2 + 9} \le h + 1$.
    \item Résoudre cette inéquation (on supposera $h > 0$).
    \item Interpréter le résultat pour le choix de la hauteur du poteau.
\end{enumerate}
\end{exercice}

\newpage

\coursec{Corrigés des exercices}

\coursec{Équations et inéquations polynomiales}

\courssub{Je vérifie mes connaissances}

\begin{corrige}[Exercice 1 — QCM : Discriminant et nombre de solutions]
On calcule le discriminant $\Delta = b^2 - 4ac$ de chaque équation.
\begin{enumerate}
    \item $2x^2 - 3x + 4 = 0$ : $\Delta = (-3)^2 - 4 \times 2 \times 4 = 9 - 32 = -23 < 0$. Aucune solution réelle. Affirmation \textbf{fausse}.
    \item $x^2 - 4x + 4 = 0$ : $\Delta = (-4)^2 - 4 \times 1 \times 4 = 16 - 16 = 0$. L'équation admet une racine double $x_0 = 2$. Affirmation \textbf{fausse}.
    \item $-x^2 + 2x - 5 = 0$ : $\Delta = 2^2 - 4 \times (-1) \times (-5) = 4 - 20 = -16 < 0$. Aucune solution réelle. Affirmation \textbf{fausse}.
    \item $3x^2 + 5x - 2 = 0$ : $\Delta = 5^2 - 4 \times 3 \times (-2) = 25 + 24 = 49 > 0$. L'équation admet deux solutions réelles distinctes : $x_1 = \dfrac{-5-7}{6} = -2$ et $x_2 = \dfrac{-5+7}{6} = \dfrac{1}{3}$. Affirmation \textbf{exacte}.
\end{enumerate}
La bonne réponse est donc la réponse \textbf{d)}.
\end{corrige}

\begin{corrige}[Exercice 2 — Vrai ou Faux : Signe d'un trinôme]
\begin{enumerate}
    \item \textbf{Vrai.} Pour $f(x) = -2x^2 + 4x - 3$ : $\Delta = 4^2 - 4 \times (-2) \times (-3) = 16 - 24 = -8 < 0$. Le trinôme ne s'annule pas et garde un signe constant sur $\mathbb{R}$, celui de $a = -2 < 0$. Donc $f(x) < 0$ pour tout $x \in \mathbb{R}$.
    \item \textbf{Faux.} $x^2 - 6x + 9 = (x-3)^2$. Un carré est toujours positif ou nul : $(x-3)^2 \le 0$ n'est vérifié que pour $x = 3$. L'ensemble des solutions est $S = \{3\}$, et non $\mathbb{R}$.
    \item \textbf{Vrai.} C'est une propriété du cours : si $\Delta < 0$, le trinôme $ax^2+bx+c$ est du signe de $a$ pour tout réel $x$. Ici $a > 0$, donc $ax^2+bx+c > 0$ sur $\mathbb{R}$.
    \item \textbf{Faux.} Pour tout $x \in \mathbb{R}$, $x^2 \ge 0$, donc $x^2 + 1 \ge 1 > 0$. L'inéquation $x^2 + 1 < 0$ n'admet aucune solution : $S = \emptyset$.
\end{enumerate}
\end{corrige}

\begin{corrige}[Exercice 3 — Calcul direct : Forme canonique et racines]
\begin{enumerate}
    \item $P_1(x) = x^2 - 5x + 6$. On a $a = 1$, $b = -5$, $c = 6$, donc $-\dfrac{b}{2a} = \dfrac{5}{2}$.
    \[ P_1(x) = \left(x - \tfrac{5}{2}\right)^2 - \tfrac{25}{4} + 6 = \left(x - \tfrac{5}{2}\right)^2 - \tfrac{1}{4}. \]
    $\Delta = 25 - 24 = 1 > 0$ : deux racines $x_1 = \dfrac{5-1}{2} = 2$ et $x_2 = \dfrac{5+1}{2} = 3$.
    \item $P_2(x) = 2x^2 + 4x + 2 = 2\left(x^2 + 2x + 1\right) = 2(x+1)^2$.
    $\Delta = 16 - 16 = 0$ : racine double $x_0 = -1$.
    \item $P_3(x) = -x^2 + 3x - 5$. On a $-\dfrac{b}{2a} = \dfrac{3}{2}$.
    \[ P_3(x) = -\left[\left(x - \tfrac{3}{2}\right)^2 - \tfrac{9}{4} + 5\right] = -\left(x - \tfrac{3}{2}\right)^2 - \tfrac{11}{4}. \]
    $\Delta = 9 - 20 = -11 < 0$ : aucune racine réelle. (On remarque d'ailleurs que $P_3(x) \le -\tfrac{11}{4} < 0$ pour tout $x$.)
\end{enumerate}
\end{corrige}

\begin{corrige}[Exercice 4 — Définitions et propriété : Polynôme de degré 3]
\begin{enumerate}
    \item $P(2) = 2 \times 2^3 - 5 \times 2^2 - 2 + 6 = 16 - 20 - 2 + 6 = 0$. Donc $2$ est bien un \cle{zéro} de $P$.
    \item D'après la propriété de divisibilité, si $a$ est un zéro de $P$, alors $P$ est divisible par $x - a$. Donc $x - 2$ est un facteur de $P(x)$.
    \item On cherche $Q(x) = 2x^2 + bx + c$ tel que $P(x) = (x-2)(2x^2 + bx + c)$. Par identification (le terme constant : $-2c = 6$ donne $c = -3$ ; le coefficient de $x$ : $-2b + c = -1$ donne $b = -1$) :
    \[ P(x) = (x-2)\left(2x^2 - x - 3\right). \]
    Vérification rapide : $(x-2)(2x^2 - x - 3) = 2x^3 - x^2 - 3x - 4x^2 + 2x + 6 = 2x^3 - 5x^2 - x + 6$. \checkmark
    \item \textbf{Propriété admise :} si un polynôme ne s'annule pas sur un intervalle $I$, alors il garde un \cle{signe constant} sur $I$.
\end{enumerate}
\end{corrige}

\courssub{Je m'entraîne}

\begin{corrige}[Exercice 5 — Équations du second degré à paramètre]
\begin{enumerate}
    \item Pour $m = 1$, l'équation devient $0 \cdot x^2 + 2x + 3 = 0$, c'est-à-dire $2x + 3 = 0$. C'est une équation du premier degré : $S = \left\{-\dfrac{3}{2}\right\}$.
    \item Pour $m \neq 1$, $(E_m)$ est du second degré avec $a = m-1$, $b = 2m$, $c = m+2$ :
    \begin{align*}
    \Delta &= (2m)^2 - 4(m-1)(m+2) = 4m^2 - 4\left(m^2 + m - 2\right) \\
    &= 4m^2 - 4m^2 - 4m + 8 = -4m + 8 = 4(2 - m).
    \end{align*}
    \item Le signe de $\Delta$ dépend de $2 - m$ :
    \begin{enumerate}
        \item \textbf{Racine double} : $\Delta = 0 \iff m = 2$ (et $2 \neq 1$, donc valide). La racine double est $x_0 = \dfrac{-2m}{2(m-1)} = \dfrac{-4}{2} = -2$.
        \item \textbf{Deux racines distinctes} : $\Delta > 0 \iff m < 2$, avec $m \neq 1$. Donc $m \in ]-\infty\,;\,1[ \cup ]1\,;\,2[$.
        \item \textbf{Aucune solution réelle} : $\Delta < 0 \iff m > 2$, c'est-à-dire $m \in ]2\,;\,+\infty[$.
    \end{enumerate}
    \item Pour $m = -2$ : $(E_{-2}) : -3x^2 - 4x + 0 = 0$, soit $-x(3x + 4) = 0$. Donc $x = 0$ ou $x = -\dfrac{4}{3}$ : $S = \left\{-\dfrac{4}{3}\,;\,0\right\}$. (Cohérent avec $\Delta = 4(2-(-2)) = 16 > 0$.)
\end{enumerate}
\textbf{Point de vigilance :} le cas $m = 1$ doit être traité \emph{à part}, car l'équation cesse d'être du second degré ; appliquer la formule du discriminant serait une erreur.
\end{corrige}

\begin{corrige}[Exercice 6 — Inéquations du second degré et tableaux de signes]
\begin{enumerate}
    \item $-2x^2 + 5x + 3 \ge 0$. Discriminant : $\Delta = 25 - 4 \times (-2) \times 3 = 25 + 24 = 49$. Racines : $x_1 = \dfrac{-5 - 7}{-4} = 3$ et $x_2 = \dfrac{-5 + 7}{-4} = -\dfrac{1}{2}$. Comme $a = -2 < 0$, le trinôme est positif \emph{entre} les racines :
    \begin{center}
    \begin{tabular}{|c|ccccccc|}\hline
    $x$ & $-\infty$ & & $-\tfrac{1}{2}$ & & $3$ & & $+\infty$ \\ \hline
    $-2x^2+5x+3$ & $-$ & & $0$ & $+$ & & $0$ & $-$ \\ \hline
    \end{tabular}
    \end{center}
    L'inéquation est au sens large :
    \[ S = \left[-\tfrac{1}{2}\,;\,3\right]. \]
    \item $x^2 - 6x + 9 < 0 \iff (x-3)^2 < 0$. Un carré n'est jamais strictement négatif : $S = \emptyset$.
    \item $3x^2 + 2x - 1 > 0$. Discriminant : $\Delta = 4 + 12 = 16$. Racines : $x_1 = \dfrac{-2-4}{6} = -1$ et $x_2 = \dfrac{-2+4}{6} = \dfrac{1}{3}$. Comme $a = 3 > 0$, le trinôme est strictement positif \emph{à l'extérieur} des racines :
    \[ S = \left]-\infty\,;\,-1\right[ \cup \left]\tfrac{1}{3}\,;\,+\infty\right[. \]
\end{enumerate}
\end{corrige}

\begin{corrige}[Exercice 7 — Problème géométrique : Rectangle]
\begin{enumerate}
    \item Le périmètre vaut $2x + 2y = 22$, donc $x + y = 11$, d'où $y = 11 - x$.
    \item L'aire vaut $xy = 28$, soit $x(11 - x) = 28$, c'est-à-dire :
    \[ x^2 - 11x + 28 = 0. \]
    \item $\Delta = 121 - 112 = 9 > 0$. Racines : $x_1 = \dfrac{11 - 3}{2} = 4$ et $x_2 = \dfrac{11 + 3}{2} = 7$.
    Si $x = 7$, alors $y = 4$ ; si $x = 4$, alors $y = 7$. Comme $x$ désigne la longueur (la plus grande dimension), on retient :
    \[ \text{Longueur} = 7\ \text{m}, \qquad \text{largeur} = 4\ \text{m}. \]
    Vérification : périmètre $= 2(7+4) = 22$ m \checkmark ; aire $= 7 \times 4 = 28\ \text{m}^2$ \checkmark.
    \item Échanger longueur et largeur donne le même rectangle : il n'y a donc qu'\textbf{une seule} configuration géométrique, décrite par les deux solutions de l'équation.
\end{enumerate}
\end{corrige}

\begin{corrige}[Exercice 8 — Polynôme de degré 3 : Factorisation, équation et inéquation]
\begin{enumerate}
    \item $P(-1) = 2(-1)^3 - 3(-1)^2 - 3(-1) + 2 = -2 - 3 + 3 + 2 = 0$. Donc $-1$ est un zéro de $P$.
    \item \textbf{Méthode des coefficients indéterminés :} on écrit $P(x) = (x+1)(2x^2 + bx + c)$. En développant : $(x+1)(2x^2+bx+c) = 2x^3 + (b+2)x^2 + (c+b)x + c$. Par identification avec $2x^3 - 3x^2 - 3x + 2$ :
    \[ \begin{cases} b + 2 = -3 \\ c = 2 \end{cases} \iff \begin{cases} b = -5 \\ c = 2 \end{cases} \]
    (vérification : $c + b = 2 - 5 = -3$ \checkmark). Donc :
    \[ P(x) = (x+1)\left(2x^2 - 5x + 2\right). \]
    \item $P(x) = 0 \iff x + 1 = 0$ ou $2x^2 - 5x + 2 = 0$. Pour le trinôme : $\Delta = 25 - 16 = 9$, racines $x = \dfrac{5-3}{4} = \dfrac{1}{2}$ et $x = \dfrac{5+3}{4} = 2$. Donc :
    \[ S = \left\{-1\,;\,\tfrac{1}{2}\,;\,2\right\}. \]
    \item On factorise complètement : $P(x) = 2(x+1)\left(x - \tfrac{1}{2}\right)(x-2)$. Tableau de signes :
    \begin{center}
    \begin{tabular}{|c|ccccccccc|}\hline
    $x$ & $-\infty$ & & $-1$ & & $\tfrac{1}{2}$ & & $2$ & & $+\infty$ \\ \hline
    $P(x)$ & $-$ & & $0$ & $+$ & & $0$ & $-$ & & $+$ \\ \hline
    \end{tabular}
    \end{center}
    Plus précisément : $P$ est négatif sur $]-\infty;-1[$, positif sur $]-1;\tfrac12[$, négatif sur $]\tfrac12;2[$, positif sur $]2;+\infty[$ (le signe change à chaque racine simple, et le signe pour $x > 2$ est celui du terme dominant $2x^3$, positif).
    \item $P(x) > 0$ :
    \[ S = \left]-1\,;\,\tfrac{1}{2}\right[ \cup \left]2\,;\,+\infty\right[. \]
\end{enumerate}
\end{corrige}

\courssub{Je me prépare au Probatoire}

\begin{corrige}[Exercice 9 — Bénéfice d'un commerçant — Synthèse 2nd degré]
\begin{enumerate}
    \item $B(x) = 0 \iff -x^2 + 10x - 16 = 0$. $\Delta = 100 - 64 = 36$. Racines : $x_1 = \dfrac{-10 + 6}{-2} = 2$ et $x_2 = \dfrac{-10 - 6}{-2} = 8$.
    Le bénéfice est nul pour $x = 2$ et $x = 8$, soit \textbf{200 sacs} et \textbf{800 sacs} vendus (seuils de rentabilité).
    \item $B$ est un trinôme avec $a = -1 < 0$ : il est positif entre ses racines, négatif à l'extérieur. Sur $[0\,;\,10]$ :
    \begin{center}
    \begin{tabular}{|c|ccccccc|}\hline
    $x$ & $0$ & & $2$ & & $8$ & & $10$ \\ \hline
    $B(x)$ & $-$ & $-$ & $0$ & $+$ & $0$ & $-$ & $-$ \\ \hline
    \end{tabular}
    \end{center}
    (précisément : $B < 0$ sur $[0;2[$, $B = 0$ en $2$ et $8$, $B > 0$ sur $]2;8[$, $B < 0$ sur $]8;10]$.)
    \item Le bénéfice est strictement positif pour $x \in ]2\,;\,8[$, c'est-à-dire pour un nombre de sacs vendus \textbf{strictement compris entre 200 et 800}.
    \item Le maximum d'un trinôme ($a < 0$) est atteint en $x = -\dfrac{b}{2a} = \dfrac{-10}{-2} = 5$.
    \[ B(5) = -25 + 50 - 16 = 9. \]
    Il faut vendre \textbf{500 sacs} ; le bénéfice maximal est $9$ dizaines de milliers de FCFA, soit $\mathbf{90\,000}$ \textbf{FCFA}.
    \item La taxe de $50\,000$ FCFA correspond à $5$ dizaines de milliers : $B_1(x) = -x^2 + 10x - 16 - 5 = -x^2 + 10x - 21$.
    Le maximum de $B_1$ est atteint au même point $x = 5$ : $B_1(5) = 9 - 5 = 4$, soit $40\,000$ FCFA. Comme $4 > 0$, le bénéfice maximal \textbf{couvre la taxe} ; il reste $40\,000$ FCFA de bénéfice net après taxe.
\end{enumerate}
\textbf{Barème indicatif (sur 5 pts) :} 1) 1 pt ; 2) 1 pt ; 3) 0,5 pt ; 4) 1 pt ; 5) 1,5 pt.
\textbf{Points de vigilance :} bien convertir les unités (centaines de sacs, dizaines de milliers de FCFA) ; justifier le maximum par $x = -\frac{b}{2a}$ avec $a < 0$ ; donner les réponses concrètes en sacs et en FCFA, pas seulement en $x$.
\end{corrige}

\begin{corrige}[Exercice 10 — Polynôme de degré 4 et équations irrationnelles]
\textbf{Partie A.}
\begin{enumerate}
    \item $x^2 - 4 = (x-2)(x+2)$ (identité remarquable). Pour $x^2 - 2x + 5$ : $\Delta = (-2)^2 - 4 \times 5 = 4 - 20 = -16 < 0$. Comme $a = 1 > 0$, ce trinôme est \textbf{strictement positif pour tout} $x \in \mathbb{R}$.
    \item Le signe de $f(x) = (x-2)(x+2)(x^2-2x+5)$ est donc celui du produit $(x-2)(x+2)$ :
    \begin{center}
    \begin{tabular}{|c|ccccccc|}\hline
    $x$ & $-\infty$ & & $-2$ & & $2$ & & $+\infty$ \\ \hline
    $x+2$ & $-$ & & $0$ & $+$ & & $+$ & \\ \hline
    $x-2$ & $-$ & & $-$ & & $0$ & $+$ & \\ \hline
    $x^2-2x+5$ & $+$ & & $+$ & & $+$ & & \\ \hline
    $f(x)$ & $+$ & & $0$ & $-$ & $0$ & $+$ & \\ \hline
    \end{tabular}
    \end{center}
    \item $f(x) \le 0$ : $S = [-2\,;\,2]$.
\end{enumerate}
\textbf{Partie B.}
\begin{enumerate}
    \item La racine carrée exige $3x - 2 \ge 0$, soit $x \ge \dfrac{2}{3}$ : $D = \left[\dfrac{2}{3}\,;\,+\infty\right[$.
    \item En élevant au carré : $3x - 2 = (x-2)^2 = x^2 - 4x + 4$, soit :
    \[ x^2 - 7x + 6 = 0. \]
    $\Delta = 49 - 24 = 25$ ; racines : $x_1 = \dfrac{7-5}{2} = 1$ et $x_2 = \dfrac{7+5}{2} = 6$.
    \item \textbf{Vérification (étape obligatoire !)} : l'élévation au carré peut introduire des solutions étrangères. De plus, le membre de droite $x - 2$ doit être positif ou nul (il est égal à une racine carrée), donc $x \ge 2$.
    \begin{itemize}
        \item $x = 1$ : $1 \in D$ mais $x - 2 = -1 < 0$ alors que $\sqrt{3 \times 1 - 2} = 1 \ge 0$. \textbf{Rejetée.}
        \item $x = 6$ : $\sqrt{3 \times 6 - 2} = \sqrt{16} = 4$ et $6 - 2 = 4$. \textbf{Acceptée.}
    \end{itemize}
    \item Conclusion : $S = \{6\}$.
\end{enumerate}
\textbf{Barème indicatif (sur 5 pts) :} Partie A : 2 pts ; Partie B : 3 pts dont 1 pt pour la vérification.
\textbf{Points de vigilance :} citer la condition $x - 2 \ge 0$ (le membre de droite doit être positif) ; la vérification des candidats dans l'équation initiale est \emph{exigée} au Probatoire ; ne pas oublier que $x^2 - 2x + 5$, sans racine, garde un signe constant.
\end{corrige}

\begin{corrige}[Exercice 11 — Inéquation irrationnelle et problème de dimensionnement]
\textbf{Partie A.}
\begin{enumerate}
    \item La racine carrée exige $x + 6 \ge 0$, soit $D = [-6\,;\,+\infty[$.
    \item Sur $D$, le membre de gauche $\sqrt{x+6}$ est toujours positif ou nul. Le signe de $x$ partage $D$ en deux zones : $x < 0$ sur $[-6\,;\,0[$ et $x \ge 0$ sur $[0\,;\,+\infty[$.
    \item \textbf{Cas 1 : $x \in [-6\,;\,0[$.} Alors $\sqrt{x+6} \ge 0 > x$ : l'inéquation est automatiquement vérifiée. Tout l'intervalle $[-6\,;\,0[$ est solution.

    \textbf{Cas 2 : $x \ge 0$.} Les deux membres sont positifs ou nuls : on peut élever au carré (la fonction carré est croissante sur $[0;+\infty[$, l'ordre est conservé) :
    \[ x + 6 \ge x^2 \iff x^2 - x - 6 \le 0. \]
    $\Delta = 1 + 24 = 25$ ; racines : $x_1 = \dfrac{1-5}{2} = -2$ et $x_2 = \dfrac{1+5}{2} = 3$. Le trinôme ($a = 1 > 0$) est négatif entre ses racines : $-2 \le x \le 3$. En intersectant avec le cas $x \ge 0$ : $x \in [0\,;\,3]$.
    \item Réunion des deux cas : $S = [-6\,;\,0[ \cup [0\,;\,3] = \mathbf{[-6\,;\,3]}$.
\end{enumerate}
\textbf{Partie B.}
\begin{enumerate}
    \item La condition « longueur du câble inférieure ou égale à la hauteur augmentée de 1 m » s'écrit :
    \[ \sqrt{h^2 + 9} \le h + 1. \]
    \item C'est exactement l'inéquation demandée : le câble mesure $\sqrt{h^2+9}$ m (théorème de Pythagore : le câble est l'hypoténuse d'un triangle rectangle de côtés $h$ et $3$) et la hauteur augmentée de 1 m est $h + 1$.
    \item Comme $h > 0$, on a $h + 1 > 0$ : les deux membres sont positifs, on peut élever au carré :
    \[ h^2 + 9 \le (h+1)^2 = h^2 + 2h + 1 \iff 9 \le 2h + 1 \iff 2h \ge 8 \iff h \ge 4. \]
    Avec la condition $h > 0$ : $S = [4\,;\,+\infty[$.
    \item \textbf{Interprétation :} le poteau doit mesurer \textbf{au moins 4 mètres} de haut pour que le câble, fixé à 3 m du pied, ne dépasse pas la hauteur du poteau augmentée de 1 m. En dessous de 4 m, le câble serait « trop long » relativement à la hauteur, ce qui traduirait un mauvais dimensionnement du haubanage.
\end{enumerate}
\textbf{Barème indicatif (sur 5 pts) :} Partie A : 3 pts (dont 1 pt pour la discussion des cas) ; Partie B : 2 pts.
\textbf{Points de vigilance :} pour une inéquation irrationnelle, on n'élève au carré que lorsque \emph{les deux membres sont positifs} — la disjonction de cas est la justification attendue ; ne pas oublier l'ensemble de définition ; conclure par une phrase d'interprétation concrète dans les problèmes.
\end{corrige}

\newpage

\coursec{Fiche bilan}

\begin{popfigure}
\begin{tikzpicture}[mindmap, grow cyclic, every node/.style={concept, text=white, font=\small\bfseries},
root concept/.append style={concept color=popdark, font=\large\bfseries},
level 1/.append style={level distance=3.6cm, sibling angle=60},
level 2/.append style={level distance=2.2cm, sibling angle=45, font=\scriptsize}]
\node[root concept]{Équations et inéquations polynomiales}
  child[concept color=popblue]{ node {Second degré}
    child { node {Forme canonique} }
    child { node {Discriminant $\Delta$} }
  }
  child[concept color=poporange]{ node {Signe du trinôme}
    child { node {Inéquations degré 2} }
  }
  child[concept color=popgreen]{ node {Systèmes}
    child { node {Somme et produit} }
  }
  child[concept color=poppurple]{ node {Polynômes degré 3}
    child { node {Racine évidente} }
    child { node {Factorisation} }
  }
  child[concept color=poppink]{ node {Signe degré 3}
    child { node {Tableau de signes} }
  }
  child[concept color=popgold]{ node {Irrationnelles}
    child { node {Conditions} }
    child { node {Vérification} }
  };
\end{tikzpicture}
\legende{Carte mentale de la leçon : les six grandes compétences à maîtriser.}
\end{popfigure}

\begin{aretenir}[L'essentiel de la leçon : définitions clés]
\begin{itemize}
\item \cle{Trinôme du second degré} : expression $P(x)=ax^2+bx+c$ avec $a\neq 0$.
\item \cle{Forme canonique} : $P(x)=a\left(x+\dfrac{b}{2a}\right)^2-\dfrac{\Delta}{4a}$ où $\Delta=b^2-4ac$.
\item \cle{Zéro (ou racine) d'un polynôme} $P$ : tout réel $\alpha$ tel que $P(\alpha)=0$.
\item \cle{Équation irrationnelle} : équation où l'inconnue apparaît sous un radical, par exemple $\sqrt{ax+b}=cx+d$.
\end{itemize}
\end{aretenir}

\begin{propriete}[Formules et théorèmes à connaître par cœur]
\begin{center}
\begin{tabularx}{\linewidth}{|l|X|}
\hline
\rowcolor{popblueL} \textbf{Notion} & \textbf{Résultat} \\
\hline
Discriminant & $\Delta=b^2-4ac$ \\
\hline
$\Delta>0$ & Deux racines : $x_1=\dfrac{-b-\sqrt{\Delta}}{2a}$ et $x_2=\dfrac{-b+\sqrt{\Delta}}{2a}$ \\
\hline
$\Delta=0$ & Racine double : $x_0=-\dfrac{b}{2a}$ \\
\hline
$\Delta<0$ & Aucune racine réelle ; $P(x)$ garde le signe de $a$ sur $\mathbb{R}$ \\
\hline
Factorisation & Si $\Delta>0$ : $P(x)=a(x-x_1)(x-x_2)$ ; si $\Delta=0$ : $P(x)=a(x-x_0)^2$ \\
\hline
Somme et produit & $S=x_1+x_2=-\dfrac{b}{a}$ ; $P=x_1x_2=\dfrac{c}{a}$ \\
\hline
Signe du trinôme & $P(x)$ est du signe de $a$ à l'extérieur des racines, du signe contraire entre les racines \\
\hline
Divisibilité & $\alpha$ racine de $P$ $\iff$ $P(x)$ est divisible par $(x-\alpha)$ : $P(x)=(x-\alpha)Q(x)$ \\
\hline
Signe constant & (Admis) Un polynôme qui ne s'annule pas sur un intervalle y garde un signe constant \\
\hline
\end{tabularx}
\end{center}
\end{propriete}

\begin{methode}[Méthodes express]
\begin{enumerate}
\item \cle{Résoudre} $ax^2+bx+c=0$ : calculer $\Delta=b^2-4ac$, puis conclure selon son signe (deux racines, racine double, ou aucune racine réelle).
\item \cle{Factoriser un polynôme de degré 3} : chercher une racine évidente $\alpha$ parmi les diviseurs du terme constant (tester $-2$ ; $-1$ ; $1$ ; $2$ ; \dots), puis écrire $P(x)=(x-\alpha)(ax^2+bx+c)$ par division euclidienne ou par coefficients indéterminés, puis résoudre le trinôme obtenu.
\item \cle{Tableau de signes} : placer toutes les racines en ordre croissant, mettre le signe de $a$ (ou du facteur dominant) à droite de la plus grande racine, puis alterner les signes pour des racines simples.
\item \cle{Système somme-produit} : si $x+y=S$ et $xy=P$, alors $x$ et $y$ sont les racines de l'équation $X^2-SX+P=0$ (condition d'existence de solutions réelles : $S^2-4P\geq 0$).
\item \cle{Équation irrationnelle} $\sqrt{A}=B$ : poser les conditions $A\geq 0$ et $B\geq 0$, élever au carré, résoudre l'équation obtenue, puis \textbf{vérifier} chaque solution dans l'équation de départ.
\end{enumerate}
\end{methode}

\begin{attention}[Pièges fréquents au Probatoire]
\begin{itemize}
\item Ne pas oublier le facteur $a$ dans la factorisation : $P(x)=a(x-x_1)(x-x_2)$ et non $(x-x_1)(x-x_2)$.
\item Pour une équation avec paramètre $m$, discuter d'abord le cas où le coefficient de $x^2$ s'annule (l'équation n'est alors plus du second degré).
\item L'élévation au carré n'est pas une équivalence : elle peut introduire des solutions étrangères. La vérification finale est obligatoire et doit figurer dans la rédaction.
\item Pour $\sqrt{A}=B$, la condition $B\geq 0$ est indispensable : une racine carrée est toujours positive.
\item Dans un tableau de signes, une racine double ne fait pas changer le signe du facteur correspondant.
\end{itemize}
\end{attention}

\begin{aretenir}[Checklist avant le Probatoire]
\begin{itemize}
\item[$\square$] Je sais passer de la forme développée à la forme canonique sans erreur de signe.
\item[$\square$] Je calcule $\Delta=b^2-4ac$ et je conclus correctement selon son signe.
\item[$\square$] Je factorise un trinôme à l'aide de ses racines : $a(x-x_1)(x-x_2)$.
\item[$\square$] Je résous une équation du second degré avec paramètre en discutant le signe de $\Delta(m)$.
\item[$\square$] Je dresse le tableau de signes d'un trinôme et je résous une inéquation du second degré.
\item[$\square$] Je résous un système somme-produit via $X^2-SX+P=0$.
\item[$\square$] Je vérifie qu'un réel $\alpha$ est zéro d'un polynôme en calculant $P(\alpha)$.
\item[$\square$] Je factorise un polynôme de degré 3 par division euclidienne ou coefficients indéterminés.
\item[$\square$] Je dresse le tableau de signes d'un polynôme de degré 3 et je résous l'inéquation associée.
\item[$\square$] Pour $\sqrt{ax+b}=cx+d$, je pose les conditions $ax+b\geq 0$ et $cx+d\geq 0$ avant d'élever au carré.
\item[$\square$] Je vérifie systématiquement les solutions d'une équation irrationnelle dans l'équation initiale.
\item[$\square$] Je rédige proprement : ensemble de définition, conditions, calculs, conclusion avec l'ensemble des solutions.
\end{itemize}
\end{aretenir}

\findecours

\end{document}
